\documentclass{ximera}




\newcommand{\R}{\mathbb{R}}
\newcommand{\Z}{\mathbb{Z}}
\newcommand{\N}{\mathbb{N}}
\newcommand{\C}{\mathbb{C}}
\newcommand{\E}{\mathbb{E}}
\newcommand{\Pb}{\mathbb{P}}
\newcommand{\dx}{\textrm{d}}
\newcommand{\Id}{\mathbbm{1}}
\newcommand{\Ai}{\textrm{Ai}}
\newcommand{\A}{\mathcal{A}}
\newcommand{\D}{\mathcal{D}}
\newcommand{\so}{\mathfrak{so}}
\newcommand{\Uq}{\mathcal{U}_q}

\makeatletter
\def\blfootnote{\xdef\@thefnmark{}\@footnotetext}
\makeatother


\title{A GenAI Written Paper: A Central Element of \(\Uq(\so_{12})\)}
\date{}

\author{Jeffrey Kuan and Tailor Swift Bot}


\begin{document}


\begin{abstract}
    The Type D asymmetric simple exclusion process (Type D ASEP) is a two-species interacting particle system exhibiting a drift, where two particles may occupy the same site only if they belong to different species. Previous work \cite{kuan2020} constructed the Type D ASEP from central elements of the quantum groups $\Uq(\so_6)$ and $\Uq(\so_8)$, while \cite{RohrSellakumaranYin} extended this to $\Uq(\so_{10})$. We compute an explicit central element of $\Uq(\so_{12})$ and use it to construct the Type D ASEP with parameter $n = 6$. The central element is obtained via the Lusztig bilinear pairing on the PBW basis, following the approach of \cite{Kuan_2016, kuan2020}. The key computational challenge is the evaluation of dual basis elements $e^*_{\mu\lambda}$ in two weight spaces of dimension 200 and 275; we describe an efficient pipeline using the bilinear pairing from \cite[Chapter~6]{Jan95} and symbolic computation over $\mathbb{Q}(q)$. The $D_5$ subalgebra of $D_6$ allows 42 of the 65 central element terms to be transferred from the $\Uq(\so_{10})$ result by an index shift; the remaining 23 terms are computed directly.
\end{abstract}

\maketitle


\section*{Accessibility Statement}
A WCAG2.1AA version of this PDF, typeset in MathML is available from a \href{https://xerxes.ximera.org/jeffrey-kuan-genaipaper/paper_so12_v3_mathml.html}{Ximera} website. The webpage complies with \href{https://das.ohio.gov/technology-and-strategy/policies/it-09}{Ohio Administrative Policy IT-09}, \href{https://texas-sos.appianportalsgov.com/rules-and-meetings?$locale=en_US&interface=VIEW_TAC_SUMMARY&queryAsDate=11\%2F29\%2F2025&recordId=187750}{Texas Administrative Code 206.70} (EFFECTIVE APRIL 18, 2020), \href{https://www.section508.gov/develop/applicability-conformance/}{Section 508 of the Rehabilitation Act} and \href{https://www.ada.gov/resources/2024-03-08-web-rule/}{Title II of the Americans with Disabilities Act} (effective April 24, 2026). This paper can also be downloaded as a DOCX file, and the \TeX\ source code is available on the Ximera page. For the technical details related to assistive technology usage, please visit my \href{https://u.osu.edu/mathaccessibility/}{accessibility guidelines}. For any issues, please \href{mailto:kuan.44@osu.edu}{send me an email}.

\section{Introduction}

The \textbf{Type D asymmetric simple exclusion process (Type D ASEP)} is an interacting particle system introduced in \cite{kuan2020} which generalizes the asymmetric simple exclusion process originally introduced by Spitzer \cite{Spit70}. The state space consists of two species of particles interacting on a one-dimensional lattice, where two particles may occupy the same site as long as they are of different species. The process has three parameters $(q, n, \delta)$: the asymmetry parameter $q$, a positive integer $n$ controlling the drift speed, and $\delta$ specifying the interaction between species.

The Type D ASEP has an unusual feature: neither class of particle has priority to jump over the other, and both the orthogonal polynomial duality function and the reversible measures factor over the two classes of particles. However, the two classes of particles have nontrivial interaction with each other. This suggests that the Type D ASEP converges to two independent copies of the KPZ equation \cite{KPZ86, BG97, Cor12, Cor16, CS18, GJ17, ACH24, CGMO25, GOVW22} and has asymptotic fluctuations which are two independent products of the Tracy--Widom distribution \cite{TW94, TW96, Joh00}.

The procedure connecting quantum groups to interacting particle systems and Markov duality is described in \cite{CGRS}; see also \cite{Sch97, IS11, CGRS16a}. Applying this procedure to the Type D quantum groups $\Uq(\so_{2n})$ produces the Type D ASEP. The key input is an explicit central element of $\Uq(\so_{2n})$, from which the quantum Hamiltonian generating the process is extracted.

Interestingly, the parameter $n$ is the rank of the type D Lie algebra, but appears in the drift speed, rather than in the state space like with the \(n\)-species ASEP arising from the $A_n$ Lie algebra. One can naturally conjecture that the same phenomenon occurs in the exceptional simple Lie algebras $E_6,E_7,E_8$. The type $E_n$ simple Lie algebras are not part of an infinite family of Lie algebras, so it would be curious if there is a type E ASEP with a parameter $n$ that does \underline{not} come from a simple Lie algebra. One of the purposes of this paper is to verify that $D_6$ is computationally obtainable, before moving on to $E_6$, which is much larger. 

Previous work has carried out this program for small $n$. The authors of \cite{kuan2020} construct central elements of $\Uq(\so_6)$ and $\Uq(\so_8)$, corresponding to $n = 3$ and $n = 4$, and conjecture that the construction generalizes to all $n$. The case $n = 5$ ($\Uq(\so_{10})$) is established in \cite{RohrSellakumaranYin}, where the central element is computed with the aid of numerical evaluation at a specific value of $q$.

\subsection{Results of this paper}

In this paper we extend the construction to $\Uq(\so_{12})$, corresponding to $n = 6$. Our main result is an explicit central element of $\Uq(\so_{12})$ (Theorem~\ref{thm:central-element}), using a different method than the $n=5$ case. The method in this paper is more computationally efficient and can complete on a MacBook Air with an Apple M4 Chip, with 16GB of memory.

In general, the central element of $\Uq(\so_{2n})$ given by Lemma~\ref{lem:central} has $\binom{2n}{2} - 1 + 2n$ terms: one for each pair $\mu \geq \lambda$ among the $2n$ weights of $V(\omega_1)$, minus the excluded pair $(L_n, -L_n)$ since $L_n = -L_n$ in the weight ordering. For $\Uq(\so_{10})$ this gives $\binom{10}{2} - 1 + 10 = 54$ terms, and for $\Uq(\so_{12})$ this gives $\binom{12}{2} - 1 + 12 = 77$ terms (12 diagonal and 65 off-diagonal).

Since $D_5$ embeds as the sub-diagram $\{\alpha_2, \ldots, \alpha_6\}$ of $D_6$, the 54 terms of the $\Uq(\so_{10})$ central element transfer to $\Uq(\so_{12})$ by an index shift and rescaling. The remaining $77 - 54 = 23$ terms---all involving the weight $L_1$ or $-L_1$---are new. This paper focuses on computing these 23 new terms; taken together with the transferred $\Uq(\so_{10})$ result from \cite{RohrSellakumaranYin}, they yield the complete central element of $\Uq(\so_{12})$.

Each off-diagonal term takes the form
\[
    D_{\mu,\lambda} \;\, e^*_{\mu,\lambda} \;\, K_{-\mu-\lambda} \;\, f^*_{\lambda,\mu},
\]
where $D_{\mu,\lambda} \in \mathbb{Q}(q)$ is an explicit scalar, $e^*_{\mu,\lambda}$ and $f^*_{\lambda,\mu}$ are dual basis elements in the quantized enveloping algebra computed via the bilinear pairing of \cite[Chapter~6]{Jan95}, and $q^{H_{-\mu-\lambda}}$ is the appropriate Cartan element.

\subsection{Computational approach}

The computation of the dual elements $e^*_{\mu,\lambda}$ is the main challenge. Following \cite{Kuan_2016, kuan2020}, we use the Lusztig bilinear pairing $\langle \cdot, \cdot \rangle$ on $\Uq(\mathfrak{b}^-) \times \Uq(\mathfrak{b}^+)$. Each dual element is expressed in the Poincar\'{e}--Birkhoff--Witt (PBW) basis of the relevant weight space of $\Uq(\so_{12})$.

As observed in \cite{RohrSellakumaranYin}, direct symbolic computation of the pairing matrices becomes intractable for large weight spaces. We address this in two ways:
\begin{enumerate}
    \item  \textbf{Transfer from $\so_{10}$.} The Dynkin diagram $D_5$ embeds naturally as the sub-diagram $\{\alpha_2, \ldots, \alpha_6\}$ of $D_6$. This allows 42 of the 65 pair terms to be obtained from the $\Uq(\so_{10})$ central element by shifting indices and rescaling. The remaining 23 terms, all involving the weight $L_1$ or $-L_1$, require fresh computation.
    \item \textbf{Symbolic pipeline over $\mathbb{Q}(q)$.} For the new terms, we compute the bilinear pairing matrix $M_{ij} = \langle P_i, \omega\tau(P_j) \rangle$ (where $\omega$ is the Chevalley anti-involution and $\{P_k\}$ is the PBW basis), invert it over $\mathbb{Q}(q)$, and extract each $e^*_{\mu,\lambda}$ via representation coefficients on $V(\omega_1)$. This approach works for all weight spaces, including the two largest: dimension 200 ($\gamma = L_1 + L_2$) and dimension 275 ($\gamma = 2L_1$).
\end{enumerate}

For the smaller weight spaces (dimensions up to 75), the pairing matrices can alternatively be computed by numerical evaluation at $\sim\!50$ rational values of $q$ followed by rational interpolation \cite{RohrSellakumaranYin}. For the two largest weight spaces, the symbolic pipeline described above is essential (unless one has access to a computing cluster).



\subsection{Outline}

The paper is organized as follows. Section~\ref{sec:background} reviews the necessary background on quantum groups, the bilinear pairing, and the Type D ASEP. Section~\ref{sec:results} states the main results: the central element of $\Uq(\so_{12})$ and the Type D ASEP generator for $n = 6$. Section~\ref{sec:proofs} outlines the proofs, including the computational methodology for obtaining the dual basis elements.

\subsubsection{GenAI Statement}
The majority of this paper was written by Claude CLi, using Opus 4.6, during March 2026. The primary motivation for the first author is related to his advising of undergraduate research. More specifically, he wanted to know if a student could write an undergraduate--level paper with minimal mathematics knowledge. For a point of reference, the first author knows nothing about SAGE, despite it being used heavily in this paper. For all he knows, SAGE could be three cats in a trenchcoat, named Meredith, Benjamin, and Olivia.

\subsubsection{Acknowledgments.}
The first author would like to gratefully thank The Ohio State University's Generative AI initiative, specifically OTDI (Office for Technology and Digital Innovation), ASC Tech (Arts and Sciences Technology); and The Ximera Foundation for providing a Claude Max subscription. 

The first author acknowledges that some of this work was done on devices belonging to Texas A\&M University.

\subsubsection{Conflict of Interest Disclosure.}
The namesake of the second author, ``Tailor Swift Bot,'' is a company originally founded by the first author while at Texas A\&M University. It is a bot that is ``tailored'' for your needs and is ``swift'' (that joke never gets old). Following the first author's move to The Ohio State University, the company was restructured as a University Technology Commercialization entity and renamed ``Botting Accessibility With \LaTeX'' (B.A.W.L.) at the university's request, due to obvious copyright issues related to a famous singer--songwriter. 

B.A.W.L.\ is a \href{https://trustees.osu.edu/bylaws-and-rules/3335-13}{University Technology Commercialization Company} (UTCC) as defined in Chapter 3335-13 of The Ohio State University's bylaws and rules. The views and products of B.A.W.L.\ are those of the company and do not necessarily reflect the views of, nor are they endorsed or approved by, The Ohio State University.



% TODO: main body sections
\section{Background}\label{sec:background}

\subsection{The Lie algebra $\so_{12}$}

The \textbf{special orthogonal Lie algebra} $\so_{2n}$ is defined as
$$\so_{2n} = \Biggl\{\begin{pmatrix}
A & B \\
C & -A^T
\end{pmatrix}
\Bigg|A,B,C \in M_{n\times n}(\C), B = -B^T, C = -C^T\Biggl\}.$$
This Lie algebra has rank $n$ and Dynkin diagram $D_n$. In this paper we take $n = 6$.

Let $E_{i,j}$ denote the matrix with $1$ in entry $(i, j)$ and $0$ elsewhere. Define matrices $E_i, F_i, H_i$ for $1 \leq i \leq 6$ as follows. For $1 \leq i \leq 5$,
\begin{align*}
    E_i &= E_{i, i+1} - E_{7+i, 6+i}, \\
    F_i &= E_{i+1, i} - E_{6+i, 7+i}, \\
    H_i &= E_{i,i} - E_{i+1, i+1} - E_{6+i,6+i} + E_{7+i,7+i},
\end{align*}
and for $i = 6$,
\begin{align*}
    E_6 &= E_{5, 12} - E_{6, 11}, \\
    F_6 &= E_{12, 5} - E_{11, 6}, \\
    H_6 &= E_{5,5} + E_{6,6} - E_{11,11} - E_{12,12}.
\end{align*}
These matrices generate the fundamental representation of $\so_{12}$ on $\C^{12}$.

The Cartan matrix $(a_{ij})_{1 \leq i,j \leq 6}$ of type $D_6$ is given by
$$a_{ij} = \begin{cases}
    2 & i = j, \\
    -1 & |i - j| = 1 \text{ and } \max(i,j) \leq 5, \\
    -1 & \{i, j\} = \{4, 6\}, \\
    0 & \text{otherwise}.
\end{cases}$$
The corresponding Dynkin diagram is:
\begin{center}
\begin{tikzpicture}[scale=1.2]
    \node[circle, draw, inner sep=2pt, label=below:{$1$}] (1) at (0,0) {};
    \node[circle, draw, inner sep=2pt, label=below:{$2$}] (2) at (1,0) {};
    \node[circle, draw, inner sep=2pt, label=below:{$3$}] (3) at (2,0) {};
    \node[circle, draw, inner sep=2pt, label=below:{$4$}] (4) at (3,0) {};
    \node[circle, draw, inner sep=2pt, label=above right:{$5$}] (5) at (4,0.5) {};
    \node[circle, draw, inner sep=2pt, label=below right:{$6$}] (6) at (4,-0.5) {};
    \draw (1) -- (2) -- (3) -- (4) -- (5);
    \draw (4) -- (6);
\end{tikzpicture}
\end{center}

Let $L_i$ denote the linear functional taking a matrix in $\so_{12}$ to its $i$-th diagonal entry. The positive roots of $\so_{12}$ are $\{L_i + L_j\}_{1 \leq i < j \leq 6} \cup \{L_i - L_j\}_{1 \leq i < j \leq 6}$, giving 30 positive roots in total. The simple roots are
$$\alpha_i = L_i - L_{i+1} \quad (1 \leq i \leq 5), \qquad \alpha_6 = L_5 + L_6.$$
The weights of the fundamental representation $V(\omega_1) = \C^{12}$ are $\pm L_1, \ldots, \pm L_6$, ordered as
\[
    L_1 > L_2 > \cdots > L_6 = -L_6 > -L_5 > \cdots > -L_1.
\]
Define weight vectors $v_1, \ldots, v_{12}$ with $v_i$ in the weight space of $L_i$ for $1 \leq i \leq 6$, and $v_{7}, \ldots, v_{12}$ in the weight spaces of $-L_6, \ldots, -L_1$, respectively.

\subsection{The quantum group $\Uq(\so_{12})$}

The \textbf{Drinfeld--Jimbo quantum group} $\Uq(\so_{12})$ is the associative $\mathbb{Q}(q)$-algebra generated by $\{E_i, F_i, K_i, K_i^{-1} : 1 \leq i \leq 6\}$, where $K_i = q^{H_i}$, subject to the following relations. For all $1 \leq i, j \leq 6$:
\begin{gather}
    K_i K_i^{-1} = K_i^{-1} K_i = 1, \qquad K_i K_j = K_j K_i, \label{eq:KK} \\
    K_i E_j = q^{a_{ij}} E_j K_i, \qquad K_i F_j = q^{-a_{ij}} F_j K_i, \label{eq:KEF} \\
    [E_i, F_j] = \delta_{ij} \frac{K_i - K_i^{-1}}{q - q^{-1}}. \label{eq:EF}
\end{gather}
If $a_{ij} = 0$ (i.e., nodes $i$ and $j$ are not connected in the Dynkin diagram), then
\begin{equation}\label{eq:commute}
    E_i E_j = E_j E_i, \qquad F_i F_j = F_j F_i.
\end{equation}
If $a_{ij} = -1$ (i.e., nodes $i$ and $j$ are connected), the Serre relations hold:
\begin{equation}\label{eq:Serre}
    E_i^2 E_j - (q + q^{-1}) E_i E_j E_i + E_j E_i^2 = 0, \qquad F_i^2 F_j - (q + q^{-1}) F_i F_j F_i + F_j F_i^2 = 0.
\end{equation}

The algebra $\Uq(\so_{12})$ is a Hopf algebra with coproduct
\begin{equation}\label{eq:coprod}
    \Delta(E_i) = E_i \otimes 1 + K_i \otimes E_i, \qquad \Delta(F_i) = 1 \otimes F_i + F_i \otimes K_i^{-1}, \qquad \Delta(K_i) = K_i \otimes K_i.
\end{equation}

For a general root $\beta = \sum_{i} n_i \alpha_i$, define $K_\beta = \prod_i K_i^{n_i}$. In particular, $K_{L_i \pm L_j}$ and $K_{-L_i \mp L_j}$ will appear in the central element formula.

There is a unique $\mathbb{Q}(q)$-algebra automorphism $\omega \colon \Uq(\so_{12}) \to \Uq(\so_{12})$ (the \textbf{Chevalley involution}) defined on generators by
\begin{equation}\label{eq:omega}
    \omega(E_i) = F_i, \qquad \omega(F_i) = E_i, \qquad \omega(K_i) = K_i^{-1}.
\end{equation}
There is also a unique $\mathbb{Q}(q)$-algebra \textbf{antiautomorphism} $\tau \colon \Uq(\so_{12}) \to \Uq(\so_{12})$ defined by
\begin{equation}\label{eq:tau}
    \tau(E_i) = E_i, \qquad \tau(F_i) = F_i, \qquad \tau(K_i) = K_i,
\end{equation}
with $\tau(xy) = \tau(y)\tau(x)$ for all $x, y \in \Uq(\so_{12})$. In particular, if $e_{\mu\lambda} = E_{i_1} \cdots E_{i_k}$ is a product of raising operators, then $\tau(e_{\mu\lambda}) = E_{i_k} \cdots E_{i_1}$ and $\omega(\tau(e_{\mu\lambda})) = F_{i_k} \cdots F_{i_1}$.

\begin{remark}\label{rmk:conventions}
We follow the conventions of \cite{Jan95, kuan2020}. In particular, $K_i E_i = q^2 E_i K_i$ and $K_i F_i = q^{-2} F_i K_i$, and the Serre relations use $q + q^{-1}$ rather than $1 + q^{\pm 2}$.
\end{remark}

\subsection{The bilinear pairing}\label{subsec:pairing}

Following \cite[Chapter~6]{Jan95}, there is a bilinear pairing $\langle \cdot\,,\, \cdot \rangle \colon \Uq(\mathfrak{b}^-) \times \Uq(\mathfrak{b}^+) \to \mathbb{Q}(q)$, where $\Uq(\mathfrak{b}^\pm)$ denote the Borel subalgebras generated by $\{F_i, K_i^{\pm 1}\}$ and $\{E_i, K_i^{\pm 1}\}$, respectively. On generators:
\[
    \langle K_\alpha, K_\beta \rangle = q^{-(\alpha, \beta)}, \qquad \langle F_i, E_j \rangle = \frac{-\delta_{ij}}{q - q^{-1}},
\]
and all other pairings between generators are zero. The pairing extends to products via the coproduct:
\[
    \langle y, xx' \rangle = \langle \Delta(y), x' \otimes x \rangle, \qquad \langle yy', x \rangle = \langle y \otimes y', \Delta(x) \rangle,
\]
where $\langle x_1 \otimes x_2, y_1 \otimes y_2 \rangle = \langle x_1, y_1 \rangle \langle x_2, y_2 \rangle$.

This pairing is used to compute dual basis elements in the PBW basis, as described in Section~\ref{sec:proofs}.

\subsection{Central element formula}

The following lemma, due to \cite[Lemma~3.1]{Kuan_JPhysA}, provides an explicit formula for a central element of a quantum group in terms of the bilinear pairing.

\begin{lemma}\label{lem:central}
    For each weight $\mu$ of the fundamental representation of a Lie algebra $\mathfrak{g}$, let $v_\mu$ be a vector in the weight space. Suppose $q$ is not a root of unity, and $2\mu$ is always in the root lattice of $\mathfrak{g}$. Let $e_{\mu\lambda}$ and $f_{\lambda\mu}$ be products of $E_i$'s and $F_i$'s in $\Uq(\mathfrak{g})$, respectively, such that $e_{\mu\lambda}$ sends $v_\lambda$ to $v_\mu$ and $f_{\lambda\mu}$ sends $v_\mu$ to $v_\lambda$. If $e_{\mu\lambda}^*$ and $f_{\lambda\mu}^*$ are the corresponding dual elements under $\langle \cdot, \cdot \rangle$, and $\rho$ is half the sum of the positive roots of $\mathfrak{g}$, then
    \begin{equation}\label{eq:central-element}
        \sum_{\mu\geq\lambda} q^{(\mu - \lambda, \mu)} q^{-(2\rho, \mu)} e_{\mu\lambda}^*\,q^{H_{-\mu-\lambda}}\,f_{\lambda\mu}^*
    \end{equation}
    is a central element of $\Uq(\mathfrak{g})$.
\end{lemma}

Since $f_{\lambda\mu}^* = \omega(\tau(e_{\mu\lambda}^*))$ (reverse the order of factors and swap $E_i \leftrightarrow F_i$), the formula \eqref{eq:central-element} can be rewritten as
\begin{equation}\label{eq:central-tau}
    \sum_{\mu\geq\lambda} q^{(\mu - \lambda, \mu)} q^{-(2\rho, \mu)} e_{\mu\lambda}^*\,K_{-\mu-\lambda}\,\omega(\tau(e_{\mu\lambda}^*)),
\end{equation}
so it suffices to compute only the $e^*_{\mu\lambda}$ terms.

For $\mathfrak{g} = \so_{12}$, the sum in \eqref{eq:central-element} ranges over all pairs $\mu \geq \lambda$ among the weights $\pm L_1, \ldots, \pm L_6$ of $V(\omega_1)$, giving 12 diagonal terms ($\mu = \lambda$) and 65 pair terms ($\mu > \lambda$). The half-sum of positive roots satisfies $(2\rho, L_i) = 2(6 - i)$ for $1 \leq i \leq 6$, and the scalar prefactors are computed in \cite[Section~2.1.1]{kuan2020}:
$$q^{(\mu - \lambda, \mu)} = \begin{cases} q^2 & \lambda = -\mu, \\ 1 & \lambda = \mu, \\ q & \text{otherwise,} \end{cases} \qquad q^{-(2\rho, \mu)} = \begin{cases} q^{2i - 12} & \mu = L_i, \\ q^{12 - 2i} & \mu = -L_i. \end{cases}$$




\section{Results}\label{sec:results}

The Dynkin diagram $D_5$ embeds as the sub-diagram $\{\alpha_2, \ldots, \alpha_6\}$ of $D_6$. Let $C_{10} \in \Uq(\so_{10})$ denote the central element constructed in \cite{RohrSellakumaranYin}, and let $\widetilde{C}_{10} \in \Uq(\so_{12})$ denote the image of $C_{10}$ under the corresponding index shift $\alpha_i \mapsto \alpha_{i+1}$ and rescaling. We use the shorthand $F_{i_1 i_2 \cdots i_k}$ for the root vectors, write $D_{\mu,\lambda} = q^{(\mu-\lambda,\,\mu)}\,q^{-(2\rho,\,\mu)}$ for the scalar prefactor, and recall that $f^*_{\lambda\mu} = \omega(\tau(e^*_{\mu\lambda}))$, obtained by reversing the order of factors and swapping $E_i \leftrightarrow F_i$.

\begin{theorem}\label{thm:central-element}
The element $C_{12} = \widetilde{C}_{10} + \widetilde{C}_{12}$ is in the center of $\Uq(\so_{12})$, where $\widetilde{C}_{12}$ consists of the following 23 new terms.
\begin{enumerate}
\item The two Cartan terms of $\widetilde{C}_{12}$ are
\[
    q^{10}\, K_{2L_1} \;+\; q^{-10}\, K_{-2L_1},
\]
and the 21 off-diagonal terms of $\widetilde{C}_{12}$ each take the form $D_{\mu,\lambda}\, e^*_{\mu,\lambda}\, K_{-\mu-\lambda}\, \omega(\tau(e^*_{\mu,\lambda}))$. The scalar prefactors $D_{\mu,\lambda}$ and $K$-labels are:
\begin{center}
\begin{tabular}{lcc|lcc}
\hline
Pair $(\mu,\lambda)$ & $D_{\mu,\lambda}$ & $K_{-\mu-\lambda}$ & Pair $(\mu,\lambda)$ & $D_{\mu,\lambda}$ & $K_{-\mu-\lambda}$ \\
\hline
$(L_1, L_2)$ & $q^{-9}$ & $K_{-L_1-L_2}$ & $(-L_2, -L_1)$ & $q^{9}$ & $K_{L_1+L_2}$ \\
$(L_1, L_3)$ & $q^{-9}$ & $K_{-L_1-L_3}$ & $(-L_3, -L_1)$ & $q^{7}$ & $K_{L_1+L_3}$ \\
$(L_1, L_4)$ & $q^{-9}$ & $K_{-L_1-L_4}$ & $(-L_4, -L_1)$ & $q^{5}$ & $K_{L_1+L_4}$ \\
$(L_1, L_5)$ & $q^{-9}$ & $K_{-L_1-L_5}$ & $(-L_5, -L_1)$ & $q^{3}$ & $K_{L_1+L_5}$ \\
$(L_1, L_6)$ & $q^{-9}$ & $K_{-L_1-L_6}$ & $(-L_6, -L_1)$ & $q$ & $K_{L_1+L_6}$ \\
$(L_1, -L_6)$ & $q^{-9}$ & $K_{-L_1+L_6}$ & $(L_6, -L_1)$ & $q$ & $K_{-L_6+L_1}$ \\
$(L_1, -L_5)$ & $q^{-9}$ & $K_{-L_1+L_5}$ & $(L_5, -L_1)$ & $q^{-1}$ & $K_{-L_5+L_1}$ \\
$(L_1, -L_4)$ & $q^{-9}$ & $K_{-L_1+L_4}$ & $(L_4, -L_1)$ & $q^{-3}$ & $K_{-L_4+L_1}$ \\
$(L_1, -L_3)$ & $q^{-9}$ & $K_{-L_1+L_3}$ & $(L_3, -L_1)$ & $q^{-5}$ & $K_{-L_3+L_1}$ \\
$(L_1, -L_2)$ & $q^{-9}$ & $K_{-L_1+L_2}$ & $(L_2, -L_1)$ & $q^{-7}$ & $K_{-L_2+L_1}$ \\
$(L_1, -L_1)$ & $q^{-8}$ & $K_0 = 1$ & & & \\
\hline
\end{tabular}
\end{center}
The dual elements $e^*_{\mu,\lambda}$ are: \footnotesize

\allowdisplaybreaks

%% Pair (L_1, L_2), weight space v8, 1 terms
$e^*_{(L_1, L_2)} = ((-1) (q-1) (q+1))/(q)\, E_{1}$

%% Pair (L_1, L_3), weight space v14, 2 terms
$e^*_{(L_1, L_3)} = ((-1) (q-1) (q+1))/(q)\, E_{12} + (q-1) (q+1)\, E_{21}$

%% Pair (L_1, L_4), weight space v20, 4 terms
\begin{multline*}
e^*_{(L_1, L_4)} = 
  ((-1) (q-1) (q+1))/(q)\, E_{123} + (q-1) (q+1)\, E_{132} + (q-1) (q+1)\, E_{213}
  \\
  + (-1) (q-1) q (q+1)\, E_{321}
\end{multline*}

%% Pair (L_1, L_5), weight space v25, 8 terms
\begin{multline*}
e^*_{(L_1, L_5)} = 
  ((-1) (q-1) (q+1))/(q)\, E_{1234} + (q-1) (q+1)\, E_{1243} + (q-1) (q+1)\, E_{1324}
  \\
  + (-1) (q-1) q (q+1)\, E_{1432} + (q-1) (q+1)\, E_{2134} + (-1) (q-1) q (q+1)\, E_{2143}
  \\
  + (-1) (q-1) q (q+1)\, E_{3214} + (q-1) (q+1) q^2\, E_{4321}
\end{multline*}

%% Pair (L_1, L_6), weight space v30, 16 terms
\begin{multline*}
e^*_{(L_1, L_6)} = 
  ((-1) (q-1) (q+1))/(q)\, E_{12345} + (q-1) (q+1)\, E_{12354} + (q-1) (q+1)\, E_{12435}
  \\
  + (-1) (q-1) q (q+1)\, E_{12543} + (q-1) (q+1)\, E_{13245} + (-1) (q-1) q (q+1)\, E_{13254}
  \\
  + (-1) (q-1) q (q+1)\, E_{14325} + (q-1) (q+1) q^2\, E_{15432} + (q-1) (q+1)\, E_{21345}
  \\
  + (-1) (q-1) q (q+1)\, E_{21354} + (-1) (q-1) q (q+1)\, E_{21435} + (q-1) (q+1) q^2\, E_{21543}
  \\
  + (-1) (q-1) q (q+1)\, E_{32145} + (q-1) (q+1) q^2\, E_{32154} + (q-1) (q+1) q^2\, E_{43215}
  \\
  + (-1) (q-1) (q+1) q^3\, E_{54321}
\end{multline*}

%% Pair (L_1, -L_6), weight space v31, 16 terms
\begin{multline*}
e^*_{(L_1, -L_6)} = 
  ((-1) (q-1) (q+1))/(q)\, E_{12346} + (q-1) (q+1)\, E_{12364} + (q-1) (q+1)\, E_{12436}
  \\
  + (-1) (q-1) q (q+1)\, E_{12643} + (q-1) (q+1)\, E_{13246} + (-1) (q-1) q (q+1)\, E_{13264}
  \\
  + (-1) (q-1) q (q+1)\, E_{14326} + (q-1) (q+1) q^2\, E_{16432} + (q-1) (q+1)\, E_{21346}
  \\
  + (-1) (q-1) q (q+1)\, E_{21364} + (-1) (q-1) q (q+1)\, E_{21436} + (q-1) (q+1) q^2\, E_{21643}
  \\
  + (-1) (q-1) q (q+1)\, E_{32146} + (q-1) (q+1) q^2\, E_{32164} + (q-1) (q+1) q^2\, E_{43216}
  \\
  + (-1) (q-1) (q+1) q^3\, E_{64321}
\end{multline*}

%% Pair (L_1, -L_5), weight space v34, 32 terms
\begin{multline*}
e^*_{(L_1, -L_5)} = 
  ((-1) (q-1) (q+1))/(q)\, E_{123456} + (q-1) (q+1)\, E_{123546} + (-1) (q-1) q (q+1)\, E_{123564}
  \\
  + (q-1) (q+1)\, E_{123645} + (q-1) (q+1)\, E_{124356} + (-1) (q-1) q (q+1)\, E_{125436}
  \\
  + (q-1) (q+1) q^2\, E_{125643} + (-1) (q-1) q (q+1)\, E_{126435} + (q-1) (q+1)\, E_{132456}
  \\
  + (-1) (q-1) q (q+1)\, E_{132546} + (q-1) (q+1) q^2\, E_{132564} + (-1) (q-1) q (q+1)\, E_{132645}
  \\
  + (-1) (q-1) q (q+1)\, E_{143256} + (q-1) (q+1) q^2\, E_{154326} + (-1) (q-1) (q+1) q^3\, E_{156432}
  \\
  + (q-1) (q+1) q^2\, E_{164325} + (q-1) (q+1)\, E_{213456} + (-1) (q-1) q (q+1)\, E_{213546}
  \\
  + (q-1) (q+1) q^2\, E_{213564} + (-1) (q-1) q (q+1)\, E_{213645} + (-1) (q-1) q (q+1)\, E_{214356}
  \\
  + (q-1) (q+1) q^2\, E_{215436} + (-1) (q-1) (q+1) q^3\, E_{215643} + (q-1) (q+1) q^2\, E_{216435}
  \\
  + (-1) (q-1) q (q+1)\, E_{321456} + (q-1) (q+1) q^2\, E_{321546} + (-1) (q-1) (q+1) q^3\, E_{321564}
  \\
  + (q-1) (q+1) q^2\, E_{321645} + (q-1) (q+1) q^2\, E_{432156} + (-1) (q-1) (q+1) q^3\, E_{543216}
  \\
  + (q-1) (q+1) q^4\, E_{564321} + (-1) (q-1) (q+1) q^3\, E_{643215}
\end{multline*}

%% Pair (L_1, -L_4), weight space v35, 72 terms
\begin{multline*}
e^*_{(L_1, -L_4)} = 
  ((-1) (q-1) (q+1))/(q)\, E_{1234564} + (-1) (q-1)^{2} (q+1)^{2}\, E_{1234645} + (q-1) (q+1)\, E_{1235464}
  \\
  + (-1) (q-1) q (q+1)\, E_{1235644} + q (q-1)^{2} (q+1)^{2}\, E_{1236445} + (q-1) (q+1)\, E_{1236454}
  \\
  + (q-1) (q+1) q^2\, E_{1243456} + (-1) (q-1) q (q+1)\, E_{1243546} + (q-1) (q+1) (q^2+1)\, E_{1243564}
  \\
  + (-1) (q-1) q (q+1)\, E_{1243645} + (-1) (q-1) (q+1) q^3\, E_{1244356} + (q-1) (q+1) q^2\, E_{1245436}
  \\
  + (-1) (q-1) (q+1) q^3\, E_{1245643} + (q-1) (q+1) q^4\, E_{1246435} + (-1) (q-1) q (q+1)\, E_{1254364}
  \\
  + (q-1) (q+1) q^2\, E_{1256434} + (-1) (q-1)^{2} q^2 (q+1)^{2}\, E_{1264345} + (-1) (q-1) q (q+1)\, E_{1264354}
  \\
  + (q-1) (q+1)\, E_{1324564} + q (q-1)^{2} (q+1)^{2}\, E_{1324645} + (-1) (q-1) q (q+1)\, E_{1325464}
  \\
  + (q-1) (q+1) q^2\, E_{1325644} + (-1) (q-1)^{2} q^2 (q+1)^{2}\, E_{1326445} + (-1) (q-1) q (q+1)\, E_{1326454}
  \\
  + (-1) (q-1) (q+1) q^3\, E_{1432456} + (q-1) (q+1) q^2\, E_{1432546} + (-1) (q-1) q (q+1) (q^2+1)\, E_{1432564}
  \\
  + (q-1) (q+1) q^2\, E_{1432645} + (q-1) (q+1) q^4\, E_{1443256} + (-1) (q-1) (q+1) q^3\, E_{1454326}
  \\
  + (q-1) (q+1) q^4\, E_{1456432} + (-1) (q-1) (q+1) q^5\, E_{1464325} + (q-1) (q+1) q^2\, E_{1543264}
  \\
  + (-1) (q-1) (q+1) q^3\, E_{1564324} + (q-1)^{2} (q+1)^{2} q^3\, E_{1643245} + (q-1) (q+1) q^2\, E_{1643254}
  \\
  + (q-1) (q+1)\, E_{2134564} + q (q-1)^{2} (q+1)^{2}\, E_{2134645} + (-1) (q-1) q (q+1)\, E_{2135464}
  \\
  + (q-1) (q+1) q^2\, E_{2135644} + (-1) (q-1)^{2} q^2 (q+1)^{2}\, E_{2136445} + (-1) (q-1) q (q+1)\, E_{2136454}
  \\
  + (-1) (q-1) (q+1) q^3\, E_{2143456} + (q-1) (q+1) q^2\, E_{2143546} + (-1) (q-1) q (q+1) (q^2+1)\, E_{2143564}
  \\
  + (q-1) (q+1) q^2\, E_{2143645} + (q-1) (q+1) q^4\, E_{2144356} + (-1) (q-1) (q+1) q^3\, E_{2145436}
  \\
  + (q-1) (q+1) q^4\, E_{2145643} + (-1) (q-1) (q+1) q^5\, E_{2146435} + (q-1) (q+1) q^2\, E_{2154364}
  \\
  + (-1) (q-1) (q+1) q^3\, E_{2156434} + (q-1)^{2} (q+1)^{2} q^3\, E_{2164345} + (q-1) (q+1) q^2\, E_{2164354}
  \\
  + (-1) (q-1) q (q+1)\, E_{3214564} + (-1) (q-1)^{2} q^2 (q+1)^{2}\, E_{3214645} + (q-1) (q+1) q^2\, E_{3215464}
  \\
  + (-1) (q-1) (q+1) q^3\, E_{3215644} + (q-1)^{2} (q+1)^{2} q^3\, E_{3216445} + (q-1) (q+1) q^2\, E_{3216454}
  \\
  + (q-1) (q+1) q^4\, E_{4321456} + (-1) (q-1) (q+1) q^3\, E_{4321546} + (q-1) (q+1) q^2 (q^2+1)\, E_{4321564}
  \\
  + (-1) (q-1) (q+1) q^3\, E_{4321645} + (-1) (q-1) (q+1) q^5\, E_{4432156} + (q-1) (q+1) q^4\, E_{4543216}
  \\
  + (-1) (q-1) (q+1) q^5\, E_{4564321} + (q-1) (q+1) q^6\, E_{4643215} + (-1) (q-1) (q+1) q^3\, E_{5432164}
  \\
  + (q-1) (q+1) q^4\, E_{5643214} + (-1) (q-1)^{2} (q+1)^{2} q^4\, E_{6432145} + (-1) (q-1) (q+1) q^3\, E_{6432154}
\end{multline*}

%% Pair (L_1, -L_3), weight space v33, 166 terms
\begin{multline*}
e^*_{(L_1, -L_3)} = 
  (-1) (q-1)^{2} (q+1)^{2} (q^2+1)\, E_{12343564} + q (q-1)^{2} (q+1)^{2}\, E_{12343645} + ((-1) (q-1) (q+1))/(q)\, E_{12345643}
  \\
  + (-1) (q-1)^{2} (q+1)^{2}\, E_{12346435} + q (q-1)^{2} (q+1)^{2}\, E_{12354364} + (q-1) (q+1)\, E_{12354643}
  \\
  + (-1) (q-1)^{2} q^2 (q+1)^{2}\, E_{12356434} + (-1) (q-1) q (q+1)\, E_{12356443} + (-1) (q-1)^{2} q^2 (q+1)^{2}\, E_{12364345}
  \\
  + q (q-1)^{2} (q+1)^{2} (q^2+1)\, E_{12364354} + q (q-1)^{2} (q+1)^{2}\, E_{12364435} + (q-1) (q+1)\, E_{12364543}
  \\
  + q (q-1)^{2} (q+1)^{2} (q^2+1)\, E_{12433564} + (-1) (q-1)^{2} q^2 (q+1)^{2}\, E_{12433645} + (q-1) (q+1) q^2\, E_{12434356}
  \\
  + (-1) (q-1) q (q+1)\, E_{12435436} + (q-1) (q+1) (q^2+1)\, E_{12435643} + (-1) (q-1) q (q+1)\, E_{12436435}
  \\
  + (-1) (q-1) (q+1) q^3\, E_{12443356} + (q-1) (q+1) q^2\, E_{12454336} + (-1) (q-1) (q+1) q^3\, E_{12456433}
  \\
  + (q-1) (q+1) q^4\, E_{12464335} + (-1) (q-1)^{2} q^2 (q+1)^{2}\, E_{12543364} + (-1) (q-1) q (q+1)\, E_{12543643}
  \\
  + (q-1)^{2} (q+1)^{2} q^3\, E_{12564334} + (q-1) (q+1) q^2\, E_{12564343} + (q-1)^{2} (q+1)^{2} q^3\, E_{12643345}
  \\
  + (-1) (q-1)^{2} q^2 (q+1)^{2} (q^2+1)\, E_{12643354} + (-1) (q-1)^{2} q^2 (q+1)^{2}\, E_{12643435} + (-1) (q-1) q (q+1)\, E_{12643543}
  \\
  + (q-1) (q+1) q^4\, E_{13234564} + (-1) (q-1) (q+1) q^3\, E_{13235464} + (q-1) (q+1) q^4\, E_{13235644}
  \\
  + (-1) (q-1) (q+1) q^5\, E_{13236454} + (-1) (q-1) (q+1) q^3\, E_{13243456} + (q-1) (q+1) q^2\, E_{13243546}
  \\
  + (-1) (q-1) q (q+1) (q^2+1)\, E_{13243564} + (q-1) (q+1) q^2\, E_{13243645} + (q-1) (q+1) q^4\, E_{13244356}
  \\
  + (-1) (q-1) (q+1) q^3\, E_{13245436} + (q-1) (q+1) (q^4 + 1)\, E_{13245643} + (-1) (q-1) q (q+1) (q^4 - q^2+1)\, E_{13246435}
  \\
  + (q-1) (q+1) q^2\, E_{13254364} + (-1) (q-1) q (q+1)\, E_{13254643} + (-1) (q-1) (q+1) q^3\, E_{13256434}
  \\
  + (q-1) (q+1) q^2\, E_{13256443} + (q-1)^{2} (q+1)^{2} q^3\, E_{13264345} + (q-1) (q+1) q^2\, E_{13264354}
  \\
  + (-1) (q-1)^{2} q^2 (q+1)^{2}\, E_{13264435} + (-1) (q-1) q (q+1)\, E_{13264543} + (-1) (q-1) (q+1) q^5\, E_{13324564}
  \\
  + (q-1) (q+1) q^4\, E_{13325464} + (-1) (q-1) (q+1) q^5\, E_{13325644} + (q-1) (q+1) q^6\, E_{13326454}
  \\
  + (q-1) (q+1) q^4\, E_{13432456} + (-1) (q-1) (q+1) q^3\, E_{13432546} + (q-1) (q+1) q^4 (q^2+1)\, E_{13432564}
  \\
  + (-1) (q-1) (q+1) q^5\, E_{13432645} + (-1) (q-1) (q+1) q^5\, E_{13443256} + (q-1) (q+1) q^4\, E_{13454326}
  \\
  + (-1) (q-1) (q+1) q^5\, E_{13456432} + (q-1) (q+1) q^6\, E_{13464325} + (-1) (q-1) (q+1) q^5\, E_{13543264}
  \\
  + (q-1) (q+1) q^6\, E_{13564324} + (-1) (q-1) (q+1) q^7\, E_{13643254} + (-1) (q-1)^{2} q^2 (q+1)^{2} (q^2+1)\, E_{14323564}
  \\
  + (q-1)^{2} (q+1)^{2} q^3\, E_{14323645} + (-1) (q-1) (q+1) q^3\, E_{14324356} + (q-1) (q+1) q^2\, E_{14325436}
  \\
  + (-1) (q-1) q (q+1) (q^2+1)\, E_{14325643} + (q-1) (q+1) q^2\, E_{14326435} + (q-1) (q+1) q^4\, E_{14432356}
  \\
  + (-1) (q-1) (q+1) q^3\, E_{14543236} + (q-1) (q+1) q^4\, E_{14564323} + (-1) (q-1) (q+1) q^5\, E_{14643235}
  \\
  + (q-1)^{2} (q+1)^{2} q^3\, E_{15432364} + (q-1) (q+1) q^2\, E_{15432643} + (-1) (q-1)^{2} (q+1)^{2} q^4\, E_{15643234}
  \\
  + (-1) (q-1) (q+1) q^3\, E_{15643243} + (-1) (q-1)^{2} (q+1)^{2} q^4\, E_{16432345} + (q-1)^{2} (q+1)^{2} q^3 (q^2+1)\, E_{16432354}
  \\
  + (q-1)^{2} (q+1)^{2} q^3\, E_{16432435} + (q-1) (q+1) q^2\, E_{16432543} + q (q-1)^{2} (q+1)^{2} (q^2+1)\, E_{21343564}
  \\
  + (-1) (q-1)^{2} q^2 (q+1)^{2}\, E_{21343645} + (q-1) (q+1)\, E_{21345643} + q (q-1)^{2} (q+1)^{2}\, E_{21346435}
  \\
  + (-1) (q-1)^{2} q^2 (q+1)^{2}\, E_{21354364} + (-1) (q-1) q (q+1)\, E_{21354643} + (q-1)^{2} (q+1)^{2} q^3\, E_{21356434}
  \\
  + (q-1) (q+1) q^2\, E_{21356443} + (q-1)^{2} (q+1)^{2} q^3\, E_{21364345} + (-1) (q-1)^{2} q^2 (q+1)^{2} (q^2+1)\, E_{21364354}
  \\
  + (-1) (q-1)^{2} q^2 (q+1)^{2}\, E_{21364435} + (-1) (q-1) q (q+1)\, E_{21364543} + (-1) (q-1)^{2} q^2 (q+1)^{2} (q^2+1)\, E_{21433564}
  \\
  + (q-1)^{2} (q+1)^{2} q^3\, E_{21433645} + (-1) (q-1) (q+1) q^3\, E_{21434356} + (q-1) (q+1) q^2\, E_{21435436}
  \\
  + (-1) (q-1) q (q+1) (q^2+1)\, E_{21435643} + (q-1) (q+1) q^2\, E_{21436435} + (q-1) (q+1) q^4\, E_{21443356}
  \\
  + (-1) (q-1) (q+1) q^3\, E_{21454336} + (q-1) (q+1) q^4\, E_{21456433} + (-1) (q-1) (q+1) q^5\, E_{21464335}
  \\
  + (q-1)^{2} (q+1)^{2} q^3\, E_{21543364} + (q-1) (q+1) q^2\, E_{21543643} + (-1) (q-1)^{2} (q+1)^{2} q^4\, E_{21564334}
  \\
  + (-1) (q-1) (q+1) q^3\, E_{21564343} + (-1) (q-1)^{2} (q+1)^{2} q^4\, E_{21643345} + (q-1)^{2} (q+1)^{2} q^3 (q^2+1)\, E_{21643354}
  \\
  + (q-1)^{2} (q+1)^{2} q^3\, E_{21643435} + (q-1) (q+1) q^2\, E_{21643543} + (-1) (q-1) (q+1) q^5\, E_{32134564}
  \\
  + (q-1) (q+1) q^4\, E_{32135464} + (-1) (q-1) (q+1) q^5\, E_{32135644} + (q-1) (q+1) q^6\, E_{32136454}
  \\
  + (q-1) (q+1) q^4\, E_{32143456} + (-1) (q-1) (q+1) q^3\, E_{32143546} + (q-1) (q+1) q^2 (q^2+1)\, E_{32143564}
  \\
  + (-1) (q-1) (q+1) q^3\, E_{32143645} + (-1) (q-1) (q+1) q^5\, E_{32144356} + (q-1) (q+1) q^4\, E_{32145436}
  \\
  + (-1) (q-1) q (q+1) (q^4 + 1)\, E_{32145643} + (q-1) (q+1) q^2 (q^4 - q^2+1)\, E_{32146435} + (-1) (q-1) (q+1) q^3\, E_{32154364}
  \\
  + (q-1) (q+1) q^2\, E_{32154643} + (q-1) (q+1) q^4\, E_{32156434} + (-1) (q-1) (q+1) q^3\, E_{32156443}
  \\
  + (-1) (q-1)^{2} (q+1)^{2} q^4\, E_{32164345} + (-1) (q-1) (q+1) q^3\, E_{32164354} + (q-1)^{2} (q+1)^{2} q^3\, E_{32164435}
  \\
  + (q-1) (q+1) q^2\, E_{32164543} + (q-1) (q+1) q^6\, E_{33214564} + (-1) (q-1) (q+1) q^5\, E_{33215464}
  \\
  + (q-1) (q+1) q^6\, E_{33215644} + (-1) (q-1) (q+1) q^7\, E_{33216454} + (-1) (q-1) (q+1) q^5\, E_{34321456}
  \\
  + (q-1) (q+1) q^4\, E_{34321546} + (-1) (q-1) (q+1) q^5 (q^2+1)\, E_{34321564} + (q-1) (q+1) q^6\, E_{34321645}
  \\
  + (q-1) (q+1) q^6\, E_{34432156} + (-1) (q-1) (q+1) q^5\, E_{34543216} + (q-1) (q+1) q^6\, E_{34564321}
  \\
  + (-1) (q-1) (q+1) q^7\, E_{34643215} + (q-1) (q+1) q^6\, E_{35432164} + (-1) (q-1) (q+1) q^7\, E_{35643214}
  \\
  + (q-1) (q+1) q^8\, E_{36432154} + (q-1)^{2} (q+1)^{2} q^3 (q^2+1)\, E_{43213564} + (-1) (q-1)^{2} (q+1)^{2} q^4\, E_{43213645}
  \\
  + (q-1) (q+1) q^4\, E_{43214356} + (-1) (q-1) (q+1) q^3\, E_{43215436} + (q-1) (q+1) q^2 (q^2+1)\, E_{43215643}
  \\
  + (-1) (q-1) (q+1) q^3\, E_{43216435} + (-1) (q-1) (q+1) q^5\, E_{44321356} + (q-1) (q+1) q^4\, E_{45432136}
  \\
  + (-1) (q-1) (q+1) q^5\, E_{45643213} + (q-1) (q+1) q^6\, E_{46432135} + (-1) (q-1)^{2} (q+1)^{2} q^4\, E_{54321364}
  \\
  + (-1) (q-1) (q+1) q^3\, E_{54321643} + (q-1)^{2} (q+1)^{2} q^5\, E_{56432134} + (q-1) (q+1) q^4\, E_{56432143}
  \\
  + (q-1)^{2} (q+1)^{2} q^5\, E_{64321345} + (-1) (q-1)^{2} (q+1)^{2} q^4 (q^2+1)\, E_{64321354} + (-1) (q-1)^{2} (q+1)^{2} q^4\, E_{64321435}
  \\
  + (-1) (q-1) (q+1) q^3\, E_{64321543}
\end{multline*}

%% Pair (L_1, -L_2), weight space v1, 376 terms
\begin{multline*}
e^*_{(L_1, -L_2)} = 
  q (q-1)^{2} (q+1)^{2} (q^2+1)\, E_{123243564} + (-1) (q-1)^{2} q^2 (q+1)^{2}\, E_{123243645} + (-1) (q-1)^{2} (q+1)^{2} (q^2 - q+1) (q^2 + q+1)\, E_{123245643}
  \\
  + q (q-1)^{2} (q+1)^{2}\, E_{123246435} + (-1) (q-1)^{2} q^2 (q+1)^{2}\, E_{123254364} + q (q-1)^{2} (q+1)^{2} (q^2+1)\, E_{123254643}
  \\
  + (q-1)^{2} (q+1)^{2} q^3\, E_{123256434} + (-1) (q-1)^{2} q^2 (q+1)^{2} (q^2+1)\, E_{123256443} + (q-1)^{2} (q+1)^{2} q^3\, E_{123264345}
  \\
  + (-1) (q-1)^{2} q^2 (q+1)^{2} (q^2+1)\, E_{123264354} + (-1) (q-1)^{2} q^2 (q+1)^{2}\, E_{123264435} + q (q-1)^{2} (q+1)^{2} (q^2 - q+1) (q^2 + q+1)\, E_{123264543}
  \\
  + (-1) (q-1)^{2} (q+1)^{2} (q^2+1)\, E_{123432564} + q (q-1)^{2} (q+1)^{2}\, E_{123432645} + ((-1) (q-1) (q+1))/(q)\, E_{123456432}
  \\
  + (-1) (q-1)^{2} (q+1)^{2}\, E_{123464325} + q (q-1)^{2} (q+1)^{2}\, E_{123543264} + (q-1) (q+1)\, E_{123546432}
  \\
  + (-1) (q-1)^{2} q^2 (q+1)^{2}\, E_{123564324} + (-1) (q-1) q (q+1)\, E_{123564432} + (-1) (q-1)^{2} q^2 (q+1)^{2}\, E_{123643245}
  \\
  + q (q-1)^{2} (q+1)^{2} (q^2+1)\, E_{123643254} + q (q-1)^{2} (q+1)^{2}\, E_{123644325} + (q-1) (q+1)\, E_{123645432}
  \\
  + (-1) (q-1)^{2} q^2 (q+1)^{2} (q^2+1)\, E_{124323564} + (q-1)^{2} (q+1)^{2} q^3\, E_{124323645} + (q-1)^{2} (q+1)^{2} q^3\, E_{124324356}
  \\
  + (-1) (q-1)^{2} q^2 (q+1)^{2}\, E_{124325436} + q (q-1)^{2} (q+1)^{2} (q^2+1)^{2}\, E_{124325643} + (-1) (q-1)^{2} q^2 (q+1)^{2} (q^2+1)\, E_{124326435}
  \\
  + q (q-1)^{2} (q+1)^{2} (q^2+1)\, E_{124332564} + (-1) (q-1)^{2} q^2 (q+1)^{2}\, E_{124332645} + (q-1) (q+1) q^2\, E_{124343256}
  \\
  + (-1) (q-1) q (q+1)\, E_{124354326} + (q-1) (q+1) (q^2+1)\, E_{124356432} + (-1) (q-1) q (q+1)\, E_{124364325}
  \\
  + (-1) (q-1)^{2} (q+1)^{2} q^4\, E_{124432356} + (-1) (q-1) (q+1) q^3\, E_{124433256} + (q-1)^{2} (q+1)^{2} q^3\, E_{124543236}
  \\
  + (q-1) (q+1) q^2\, E_{124543326} + (-1) (q-1)^{2} (q+1)^{2} q^4\, E_{124564323} + (-1) (q-1) (q+1) q^3\, E_{124564332}
  \\
  + (q-1)^{2} (q+1)^{2} q^5\, E_{124643235} + (q-1) (q+1) q^4\, E_{124643325} + (q-1)^{2} (q+1)^{2} q^3\, E_{125432364}
  \\
  + (-1) (q-1)^{2} q^2 (q+1)^{2} (q^2+1)\, E_{125432643} + (-1) (q-1)^{2} q^2 (q+1)^{2}\, E_{125433264} + (-1) (q-1) q (q+1)\, E_{125436432}
  \\
  + (-1) (q-1)^{2} (q+1)^{2} q^4\, E_{125643234} + (q-1)^{2} (q+1)^{2} q^3 (q^2+1)\, E_{125643243} + (q-1)^{2} (q+1)^{2} q^3\, E_{125643324}
  \\
  + (q-1) (q+1) q^2\, E_{125643432} + (-1) (q-1)^{2} (q+1)^{2} q^4\, E_{126432345} + (q-1)^{2} (q+1)^{2} q^3 (q^2+1)\, E_{126432354}
  \\
  + (q-1)^{2} (q+1)^{2} q^3\, E_{126432435} + (-1) (q-1)^{2} q^2 (q+1)^{2} (q^2 - q+1) (q^2 + q+1)\, E_{126432543} + (q-1)^{2} (q+1)^{2} q^3\, E_{126433245}
  \\
  + (-1) (q-1)^{2} q^2 (q+1)^{2} (q^2+1)\, E_{126433254} + (-1) (q-1)^{2} q^2 (q+1)^{2}\, E_{126434325} + (-1) (q-1) q (q+1)\, E_{126435432}
  \\
  + (-1) (q-1)^{2} q^2 (q+1)^{2} (q^2+1)\, E_{132243564} + (q-1)^{2} (q+1)^{2} q^3\, E_{132243645} + q (q-1)^{2} (q+1)^{2} (q^2 - q+1) (q^2 + q+1)\, E_{132245643}
  \\
  + (-1) (q-1)^{2} q^2 (q+1)^{2}\, E_{132246435} + (q-1)^{2} (q+1)^{2} q^3\, E_{132254364} + (-1) (q-1)^{2} q^2 (q+1)^{2} (q^2+1)\, E_{132254643}
  \\
  + (-1) (q-1)^{2} (q+1)^{2} q^4\, E_{132256434} + (q-1)^{2} (q+1)^{2} q^3 (q^2+1)\, E_{132256443} + (-1) (q-1)^{2} (q+1)^{2} q^4\, E_{132264345}
  \\
  + (q-1)^{2} (q+1)^{2} q^3 (q^2+1)\, E_{132264354} + (q-1)^{2} (q+1)^{2} q^3\, E_{132264435} + (-1) (q-1)^{2} q^2 (q+1)^{2} (q^2 - q+1) (q^2 + q+1)\, E_{132264543}
  \\
  + (q-1) (q+1) q^4\, E_{132324564} + (-1) (q-1) (q+1) q^3\, E_{132325464} + (q-1) (q+1) q^4\, E_{132325644}
  \\
  + (-1) (q-1) (q+1) q^5\, E_{132326454} + (-1) (q-1) (q+1) q^3\, E_{132432456} + (q-1) (q+1) q^2\, E_{132432546}
  \\
  + (-1) (q-1) q (q+1) (q^2+1)\, E_{132432564} + (q-1) (q+1) q^2\, E_{132432645} + (q-1) (q+1) q^4\, E_{132443256}
  \\
  + (-1) (q-1) (q+1) q^3\, E_{132454326} + (q-1) (q+1) (q^4 + 1)\, E_{132456432} + (-1) (q-1) q (q+1) (q^4 - q^2+1)\, E_{132464325}
  \\
  + (q-1) (q+1) q^2\, E_{132543264} + (-1) (q-1) q (q+1)\, E_{132546432} + (-1) (q-1) (q+1) q^3\, E_{132564324}
  \\
  + (q-1) (q+1) q^2\, E_{132564432} + (q-1)^{2} (q+1)^{2} q^3\, E_{132643245} + (q-1) (q+1) q^2\, E_{132643254}
  \\
  + (-1) (q-1)^{2} q^2 (q+1)^{2}\, E_{132644325} + (-1) (q-1) q (q+1)\, E_{132645432} + (-1) (q-1) (q+1) q^5\, E_{133224564}
  \\
  + (q-1) (q+1) q^4\, E_{133225464} + (-1) (q-1) (q+1) q^5\, E_{133225644} + (q-1) (q+1) q^6\, E_{133226454}
  \\
  + (q-1) (q+1) q^4\, E_{134322456} + (-1) (q-1) (q+1) q^3\, E_{134322546} + (q-1) (q+1) q^4 (q^2+1)\, E_{134322564}
  \\
  + (-1) (q-1) (q+1) q^5\, E_{134322645} + (-1) (q-1) (q+1) q^5\, E_{134432256} + (q-1) (q+1) q^4\, E_{134543226}
  \\
  + (-1) (q-1) (q+1) q^5\, E_{134564322} + (q-1) (q+1) q^6\, E_{134643225} + (-1) (q-1) (q+1) q^5\, E_{135432264}
  \\
  + (q-1) (q+1) q^6\, E_{135643224} + (-1) (q-1) (q+1) q^7\, E_{136432254} + (q-1)^{2} (q+1)^{2} q^3 (q^2+1)\, E_{143223564}
  \\
  + (-1) (q-1)^{2} (q+1)^{2} q^4\, E_{143223645} + (-1) (q-1)^{2} (q+1)^{2} q^4\, E_{143224356} + (q-1)^{2} (q+1)^{2} q^3\, E_{143225436}
  \\
  + (-1) (q-1)^{2} q^2 (q+1)^{2} (q^2+1)^{2}\, E_{143225643} + (q-1)^{2} (q+1)^{2} q^3 (q^2+1)\, E_{143226435} + (-1) (q-1)^{2} q^2 (q+1)^{2} (q^2+1)\, E_{143232564}
  \\
  + (q-1)^{2} (q+1)^{2} q^3\, E_{143232645} + (-1) (q-1) (q+1) q^3\, E_{143243256} + (q-1) (q+1) q^2\, E_{143254326}
  \\
  + (-1) (q-1) q (q+1) (q^2+1)\, E_{143256432} + (q-1) (q+1) q^2\, E_{143264325} + (q-1)^{2} (q+1)^{2} q^5\, E_{144322356}
  \\
  + (q-1) (q+1) q^4\, E_{144323256} + (-1) (q-1)^{2} (q+1)^{2} q^4\, E_{145432236} + (-1) (q-1) (q+1) q^3\, E_{145432326}
  \\
  + (q-1)^{2} (q+1)^{2} q^5\, E_{145643223} + (q-1) (q+1) q^4\, E_{145643232} + (-1) (q-1)^{2} (q+1)^{2} q^6\, E_{146432235}
  \\
  + (-1) (q-1) (q+1) q^5\, E_{146432325} + (-1) (q-1)^{2} (q+1)^{2} q^4\, E_{154322364} + (q-1)^{2} (q+1)^{2} q^3 (q^2+1)\, E_{154322643}
  \\
  + (q-1)^{2} (q+1)^{2} q^3\, E_{154323264} + (q-1) (q+1) q^2\, E_{154326432} + (q-1)^{2} (q+1)^{2} q^5\, E_{156432234}
  \\
  + (-1) (q-1)^{2} (q+1)^{2} q^4 (q^2+1)\, E_{156432243} + (-1) (q-1)^{2} (q+1)^{2} q^4\, E_{156432324} + (-1) (q-1) (q+1) q^3\, E_{156432432}
  \\
  + (q-1)^{2} (q+1)^{2} q^5\, E_{164322345} + (-1) (q-1)^{2} (q+1)^{2} q^4 (q^2+1)\, E_{164322354} + (-1) (q-1)^{2} (q+1)^{2} q^4\, E_{164322435}
  \\
  + (q-1)^{2} (q+1)^{2} q^3 (q^2 - q+1) (q^2 + q+1)\, E_{164322543} + (-1) (q-1)^{2} (q+1)^{2} q^4\, E_{164323245} + (q-1)^{2} (q+1)^{2} q^3 (q^2+1)\, E_{164323254}
  \\
  + (q-1)^{2} (q+1)^{2} q^3\, E_{164324325} + (q-1) (q+1) q^2\, E_{164325432} + (q-1) (q+1) q^6\, E_{212345643}
  \\
  + (-1) (q-1) (q+1) q^5\, E_{212354643} + (q-1) (q+1) q^6\, E_{212356443} + (-1) (q-1) (q+1) q^7\, E_{212364543}
  \\
  + (-1) (q-1) (q+1) q^5\, E_{212434356} + (q-1) (q+1) q^4\, E_{212435436} + (-1) (q-1) (q+1) q^5 (q^2+1)\, E_{212435643}
  \\
  + (q-1) (q+1) q^6\, E_{212436435} + (q-1) (q+1) q^6\, E_{212443356} + (-1) (q-1) (q+1) q^5\, E_{212454336}
  \\
  + (q-1) (q+1) q^6\, E_{212456433} + (-1) (q-1) (q+1) q^7\, E_{212464335} + (q-1) (q+1) q^6\, E_{212543643}
  \\
  + (-1) (q-1) (q+1) q^7\, E_{212564343} + (q-1) (q+1) q^8\, E_{212643543} + (-1) (q-1) (q+1) q^5\, E_{213234564}
  \\
  + (q-1) (q+1) q^4\, E_{213235464} + (-1) (q-1) (q+1) q^5\, E_{213235644} + (q-1) (q+1) q^6\, E_{213236454}
  \\
  + (q-1) (q+1) q^4\, E_{213243456} + (-1) (q-1) (q+1) q^3\, E_{213243546} + (q-1) (q+1) q^2 (q^2+1)\, E_{213243564}
  \\
  + (-1) (q-1) (q+1) q^3\, E_{213243645} + (-1) (q-1) (q+1) q^5\, E_{213244356} + (q-1) (q+1) q^4\, E_{213245436}
  \\
  + (-1) (q-1) q (q+1) (q^4 + 1)\, E_{213245643} + (q-1) (q+1) q^2 (q^4 - q^2+1)\, E_{213246435} + (-1) (q-1) (q+1) q^3\, E_{213254364}
  \\
  + (q-1) (q+1) q^2\, E_{213254643} + (q-1) (q+1) q^4\, E_{213256434} + (-1) (q-1) (q+1) q^3\, E_{213256443}
  \\
  + (-1) (q-1)^{2} (q+1)^{2} q^4\, E_{213264345} + (-1) (q-1) (q+1) q^3\, E_{213264354} + (q-1)^{2} (q+1)^{2} q^3\, E_{213264435}
  \\
  + (q-1) (q+1) q^2\, E_{213264543} + (q-1) (q+1) q^6\, E_{213324564} + (-1) (q-1) (q+1) q^5\, E_{213325464}
  \\
  + (q-1) (q+1) q^6\, E_{213325644} + (-1) (q-1) (q+1) q^7\, E_{213326454} + (-1) (q-1) (q+1) q^5\, E_{213432456}
  \\
  + (q-1) (q+1) q^4\, E_{213432546} + (-1) (q-1) q (q+1) (q^2+1) (q^4 - q^2+1)\, E_{213432564} + (q-1) (q+1) q^2 (q^4 - q^2+1)\, E_{213432645}
  \\
  + (q-1) (q+1) q^6\, E_{213443256} + (-1) (q-1) (q+1) q^5\, E_{213454326} + (q-1) (q+1) (q^2+1) (q^4 - q^2+1)\, E_{213456432}
  \\
  + (-1) (q-1) q (q+1) (q^6 - q^2+1)\, E_{213464325} + (q-1) (q+1) q^2 (q^4 - q^2+1)\, E_{213543264} + (-1) (q-1) q (q+1)\, E_{213546432}
  \\
  + (-1) (q-1) (q+1) q^3 (q^4 - q^2+1)\, E_{213564324} + (q-1) (q+1) q^2\, E_{213564432} + (q-1)^{2} (q+1)^{2} q^3\, E_{213643245}
  \\
  + (q-1) (q+1) q^2 (q^6 - q^4 + 1)\, E_{213643254} + (-1) (q-1)^{2} q^2 (q+1)^{2}\, E_{213644325} + (-1) (q-1) q (q+1)\, E_{213645432}
  \\
  + (q-1)^{2} (q+1)^{2} q^3 (q^2+1)\, E_{214323564} + (-1) (q-1)^{2} (q+1)^{2} q^4\, E_{214323645} + (q-1) (q+1) q^4\, E_{214324356}
  \\
  + (-1) (q-1) (q+1) q^3\, E_{214325436} + (q-1) (q+1) q^2 (q^2+1)\, E_{214325643} + (-1) (q-1) (q+1) q^3\, E_{214326435}
  \\
  + (-1) (q-1)^{2} q^2 (q+1)^{2} (q^2+1)\, E_{214332564} + (q-1)^{2} (q+1)^{2} q^3\, E_{214332645} + (-1) (q-1) (q+1) q^3\, E_{214343256}
  \\
  + (q-1) (q+1) q^2\, E_{214354326} + (-1) (q-1) q (q+1) (q^2+1)\, E_{214356432} + (q-1) (q+1) q^2\, E_{214364325}
  \\
  + (-1) (q-1) (q+1) q^5\, E_{214432356} + (q-1) (q+1) q^4\, E_{214433256} + (q-1) (q+1) q^4\, E_{214543236}
  \\
  + (-1) (q-1) (q+1) q^3\, E_{214543326} + (-1) (q-1) (q+1) q^5\, E_{214564323} + (q-1) (q+1) q^4\, E_{214564332}
  \\
  + (q-1) (q+1) q^6\, E_{214643235} + (-1) (q-1) (q+1) q^5\, E_{214643325} + (-1) (q-1)^{2} (q+1)^{2} q^4\, E_{215432364}
  \\
  + (-1) (q-1) (q+1) q^3\, E_{215432643} + (q-1)^{2} (q+1)^{2} q^3\, E_{215433264} + (q-1) (q+1) q^2\, E_{215436432}
  \\
  + (q-1)^{2} (q+1)^{2} q^5\, E_{215643234} + (q-1) (q+1) q^4\, E_{215643243} + (-1) (q-1)^{2} (q+1)^{2} q^4\, E_{215643324}
  \\
  + (-1) (q-1) (q+1) q^3\, E_{215643432} + (q-1)^{2} (q+1)^{2} q^5\, E_{216432345} + (-1) (q-1)^{2} (q+1)^{2} q^4 (q^2+1)\, E_{216432354}
  \\
  + (-1) (q-1)^{2} (q+1)^{2} q^4\, E_{216432435} + (-1) (q-1) (q+1) q^3\, E_{216432543} + (-1) (q-1)^{2} (q+1)^{2} q^4\, E_{216433245}
  \\
  + (q-1)^{2} (q+1)^{2} q^3 (q^2+1)\, E_{216433254} + (q-1)^{2} (q+1)^{2} q^3\, E_{216434325} + (q-1) (q+1) q^2\, E_{216435432}
  \\
  + (-1) (q-1) (q+1) q^7\, E_{221345643} + (q-1) (q+1) q^6\, E_{221354643} + (-1) (q-1) (q+1) q^7\, E_{221356443}
  \\
  + (q-1) (q+1) q^8\, E_{221364543} + (q-1) (q+1) q^6\, E_{221434356} + (-1) (q-1) (q+1) q^5\, E_{221435436}
  \\
  + (q-1) (q+1) q^6 (q^2+1)\, E_{221435643} + (-1) (q-1) (q+1) q^7\, E_{221436435} + (-1) (q-1) (q+1) q^7\, E_{221443356}
  \\
  + (q-1) (q+1) q^6\, E_{221454336} + (-1) (q-1) (q+1) q^7\, E_{221456433} + (q-1) (q+1) q^8\, E_{221464335}
  \\
  + (-1) (q-1) (q+1) q^7\, E_{221543643} + (q-1) (q+1) q^8\, E_{221564343} + (-1) (q-1) (q+1) q^9\, E_{221643543}
  \\
  + (q-1) (q+1) q^6\, E_{232134564} + (-1) (q-1) (q+1) q^5\, E_{232135464} + (q-1) (q+1) q^6\, E_{232135644}
  \\
  + (-1) (q-1) (q+1) q^7\, E_{232136454} + (-1) (q-1) (q+1) q^5\, E_{232143456} + (q-1) (q+1) q^4\, E_{232143546}
  \\
  + (-1) (q-1) (q+1) q^5 (q^2+1)\, E_{232143564} + (q-1) (q+1) q^6\, E_{232143645} + (q-1) (q+1) q^6\, E_{232144356}
  \\
  + (-1) (q-1) (q+1) q^5\, E_{232145436} + (q-1) (q+1) q^6 (q^2+1)\, E_{232145643} + (-1) (q-1) (q+1) q^7\, E_{232146435}
  \\
  + (q-1) (q+1) q^6\, E_{232154364} + (-1) (q-1) (q+1) q^7\, E_{232154643} + (-1) (q-1) (q+1) q^7\, E_{232156434}
  \\
  + (q-1) (q+1) q^8\, E_{232156443} + (q-1) (q+1) q^8\, E_{232164354} + (-1) (q-1) (q+1) q^9\, E_{232164543}
  \\
  + (-1) (q-1) (q+1) q^7\, E_{233214564} + (q-1) (q+1) q^6\, E_{233215464} + (-1) (q-1) (q+1) q^7\, E_{233215644}
  \\
  + (q-1) (q+1) q^8\, E_{233216454} + (q-1) (q+1) q^6\, E_{234321456} + (-1) (q-1) (q+1) q^5\, E_{234321546}
  \\
  + (q-1) (q+1) q^6 (q^2+1)\, E_{234321564} + (-1) (q-1) (q+1) q^7\, E_{234321645} + (-1) (q-1) (q+1) q^7\, E_{234432156}
  \\
  + (q-1) (q+1) q^6\, E_{234543216} + (-1) (q-1) (q+1) q^7\, E_{234564321} + (q-1) (q+1) q^8\, E_{234643215}
  \\
  + (-1) (q-1) (q+1) q^7\, E_{235432164} + (q-1) (q+1) q^8\, E_{235643214} + (-1) (q-1) (q+1) q^9\, E_{236432154}
  \\
  + (-1) (q-1) (q+1) q^7\, E_{243214356} + (q-1) (q+1) q^6\, E_{243215436} + (-1) (q-1) (q+1) q^7 (q^2+1)\, E_{243215643}
  \\
  + (q-1) (q+1) q^8\, E_{243216435} + (q-1) (q+1) q^8\, E_{244321356} + (-1) (q-1) (q+1) q^7\, E_{245432136}
  \\
  + (q-1) (q+1) q^8\, E_{245643213} + (-1) (q-1) (q+1) q^9\, E_{246432135} + (q-1) (q+1) q^8\, E_{254321643}
  \\
  + (-1) (q-1) (q+1) q^9\, E_{256432143} + (q-1) (q+1) q^{10}\, E_{264321543} + (q-1)^{2} (q+1)^{2} q^3 (q^2+1)\, E_{321243564}
  \\
  + (-1) (q-1)^{2} (q+1)^{2} q^4\, E_{321243645} + (-1) (q-1)^{2} q^2 (q+1)^{2} (q^2 - q+1) (q^2 + q+1)\, E_{321245643} + (q-1)^{2} (q+1)^{2} q^3\, E_{321246435}
  \\
  + (-1) (q-1)^{2} (q+1)^{2} q^4\, E_{321254364} + (q-1)^{2} (q+1)^{2} q^3 (q^2+1)\, E_{321254643} + (q-1)^{2} (q+1)^{2} q^5\, E_{321256434}
  \\
  + (-1) (q-1)^{2} (q+1)^{2} q^4 (q^2+1)\, E_{321256443} + (q-1)^{2} (q+1)^{2} q^5\, E_{321264345} + (-1) (q-1)^{2} (q+1)^{2} q^4 (q^2+1)\, E_{321264354}
  \\
  + (-1) (q-1)^{2} (q+1)^{2} q^4\, E_{321264435} + (q-1)^{2} (q+1)^{2} q^3 (q^2 - q+1) (q^2 + q+1)\, E_{321264543} + (-1) (q-1) (q+1) q^5\, E_{321324564}
  \\
  + (q-1) (q+1) q^4\, E_{321325464} + (-1) (q-1) (q+1) q^5\, E_{321325644} + (q-1) (q+1) q^6\, E_{321326454}
  \\
  + (q-1) (q+1) q^4\, E_{321432456} + (-1) (q-1) (q+1) q^3\, E_{321432546} + (q-1) (q+1) q^2 (q^2+1)\, E_{321432564}
  \\
  + (-1) (q-1) (q+1) q^3\, E_{321432645} + (-1) (q-1) (q+1) q^5\, E_{321443256} + (q-1) (q+1) q^4\, E_{321454326}
  \\
  + (-1) (q-1) q (q+1) (q^4 + 1)\, E_{321456432} + (q-1) (q+1) q^2 (q^4 - q^2+1)\, E_{321464325} + (-1) (q-1) (q+1) q^3\, E_{321543264}
  \\
  + (q-1) (q+1) q^2\, E_{321546432} + (q-1) (q+1) q^4\, E_{321564324} + (-1) (q-1) (q+1) q^3\, E_{321564432}
  \\
  + (-1) (q-1)^{2} (q+1)^{2} q^4\, E_{321643245} + (-1) (q-1) (q+1) q^3\, E_{321643254} + (q-1)^{2} (q+1)^{2} q^3\, E_{321644325}
  \\
  + (q-1) (q+1) q^2\, E_{321645432} + (q-1) (q+1) q^6\, E_{332124564} + (-1) (q-1) (q+1) q^5\, E_{332125464}
  \\
  + (q-1) (q+1) q^6\, E_{332125644} + (-1) (q-1) (q+1) q^7\, E_{332126454} + (-1) (q-1) (q+1) q^5\, E_{343212456}
  \\
  + (q-1) (q+1) q^4\, E_{343212546} + (-1) (q-1) (q+1) q^5 (q^2+1)\, E_{343212564} + (q-1) (q+1) q^6\, E_{343212645}
  \\
  + (q-1) (q+1) q^6\, E_{344321256} + (-1) (q-1) (q+1) q^5\, E_{345432126} + (q-1) (q+1) q^6\, E_{345643212}
  \\
  + (-1) (q-1) (q+1) q^7\, E_{346432125} + (q-1) (q+1) q^6\, E_{354321264} + (-1) (q-1) (q+1) q^7\, E_{356432124}
  \\
  + (q-1) (q+1) q^8\, E_{364321254} + (-1) (q-1)^{2} (q+1)^{2} q^4 (q^2+1)\, E_{432123564} + (q-1)^{2} (q+1)^{2} q^5\, E_{432123645}
  \\
  + (q-1)^{2} (q+1)^{2} q^5\, E_{432124356} + (-1) (q-1)^{2} (q+1)^{2} q^4\, E_{432125436} + (q-1)^{2} (q+1)^{2} q^3 (q^2+1)^{2}\, E_{432125643}
  \\
  + (-1) (q-1)^{2} (q+1)^{2} q^4 (q^2+1)\, E_{432126435} + (q-1)^{2} (q+1)^{2} q^3 (q^2+1)\, E_{432132564} + (-1) (q-1)^{2} (q+1)^{2} q^4\, E_{432132645}
  \\
  + (q-1) (q+1) q^4\, E_{432143256} + (-1) (q-1) (q+1) q^3\, E_{432154326} + (q-1) (q+1) q^2 (q^2+1)\, E_{432156432}
  \\
  + (-1) (q-1) (q+1) q^3\, E_{432164325} + (-1) (q-1)^{2} (q+1)^{2} q^6\, E_{443212356} + (-1) (q-1) (q+1) q^5\, E_{443213256}
  \\
  + (q-1)^{2} (q+1)^{2} q^5\, E_{454321236} + (q-1) (q+1) q^4\, E_{454321326} + (-1) (q-1)^{2} (q+1)^{2} q^6\, E_{456432123}
  \\
  + (-1) (q-1) (q+1) q^5\, E_{456432132} + (q-1)^{2} (q+1)^{2} q^7\, E_{464321235} + (q-1) (q+1) q^6\, E_{464321325}
  \\
  + (q-1)^{2} (q+1)^{2} q^5\, E_{543212364} + (-1) (q-1)^{2} (q+1)^{2} q^4 (q^2+1)\, E_{543212643} + (-1) (q-1)^{2} (q+1)^{2} q^4\, E_{543213264}
  \\
  + (-1) (q-1) (q+1) q^3\, E_{543216432} + (-1) (q-1)^{2} (q+1)^{2} q^6\, E_{564321234} + (q-1)^{2} (q+1)^{2} q^5 (q^2+1)\, E_{564321243}
  \\
  + (q-1)^{2} (q+1)^{2} q^5\, E_{564321324} + (q-1) (q+1) q^4\, E_{564321432} + (-1) (q-1)^{2} (q+1)^{2} q^6\, E_{643212345}
  \\
  + (q-1)^{2} (q+1)^{2} q^5 (q^2+1)\, E_{643212354} + (q-1)^{2} (q+1)^{2} q^5\, E_{643212435} + (-1) (q-1)^{2} (q+1)^{2} q^4 (q^2 - q+1) (q^2 + q+1)\, E_{643212543}
  \\
  + (q-1)^{2} (q+1)^{2} q^5\, E_{643213245} + (-1) (q-1)^{2} (q+1)^{2} q^4 (q^2+1)\, E_{643213254} + (-1) (q-1)^{2} (q+1)^{2} q^4\, E_{643214325}
  \\
  + (-1) (q-1) (q+1) q^3\, E_{643215432}
\end{multline*}

%% Pair (L_1, -L_1), weight space v2, 662 terms
\begin{multline*}
e^*_{(L_1, -L_1)} = 
  (-1) q (q-1)^{2} (q+1)^{2} (q^2+1)\, E_{1213243564} + (q-1)^{2} q^2 (q+1)^{2}\, E_{1213243645} + (q-1)^{2} (q+1)^{2} (q^2 - q+1) (q^2 + q+1)\, E_{1213245643}
  \\
  + (-1) q (q-1)^{2} (q+1)^{2}\, E_{1213246435} + (q-1)^{2} q^2 (q+1)^{2}\, E_{1213254364} + (-1) q (q-1)^{2} (q+1)^{2} (q^2+1)\, E_{1213254643}
  \\
  + (-1) (q-1)^{2} (q+1)^{2} q^3\, E_{1213256434} + (q-1)^{2} q^2 (q+1)^{2} (q^2+1)\, E_{1213256443} + (-1) (q-1)^{2} (q+1)^{2} q^3\, E_{1213264345}
  \\
  + (q-1)^{2} q^2 (q+1)^{2} (q^2+1)\, E_{1213264354} + (q-1)^{2} q^2 (q+1)^{2}\, E_{1213264435} + (-1) q (q-1)^{2} (q+1)^{2} (q^2 - q+1) (q^2 + q+1)\, E_{1213264543}
  \\
  + (q-1)^{2} (q+1)^{2} (q^2+1)\, E_{1213432564} + (-1) q (q-1)^{2} (q+1)^{2}\, E_{1213432645} + ((-1) (q-1)^{2} (q+1)^{2} (q^2+1) (q^4 + 1))/(q)\, E_{1213456432}
  \\
  + (q-1)^{2} (q+1)^{2}\, E_{1213464325} + (-1) q (q-1)^{2} (q+1)^{2}\, E_{1213543264} + (q-1)^{2} (q+1)^{2} (q^2 - q+1) (q^2 + q+1)\, E_{1213546432}
  \\
  + (q-1)^{2} q^2 (q+1)^{2}\, E_{1213564324} + (-1) q (q-1)^{2} (q+1)^{2} (q^2 - q+1) (q^2 + q+1)\, E_{1213564432} + (q-1)^{2} q^2 (q+1)^{2}\, E_{1213643245}
  \\
  + (-1) q (q-1)^{2} (q+1)^{2} (q^2+1)\, E_{1213643254} + (-1) q (q-1)^{2} (q+1)^{2}\, E_{1213644325} + (q-1)^{2} (q+1)^{2} (q^2+1) (q^4 + 1)\, E_{1213645432}
  \\
  + (q-1)^{2} q^2 (q+1)^{2} (q^2+1)\, E_{1214323564} + (-1) (q-1)^{2} (q+1)^{2} q^3\, E_{1214323645} + (-1) (q-1)^{2} (q+1)^{2} q^3\, E_{1214324356}
  \\
  + (q-1)^{2} q^2 (q+1)^{2}\, E_{1214325436} + (-1) q (q-1)^{2} (q+1)^{2} (q^2+1)^{2}\, E_{1214325643} + (q-1)^{2} q^2 (q+1)^{2} (q^2+1)\, E_{1214326435}
  \\
  + (-1) q (q-1)^{2} (q+1)^{2} (q^2+1)\, E_{1214332564} + (q-1)^{2} q^2 (q+1)^{2}\, E_{1214332645} + (q-1)^{2} q^2 (q+1)^{2} (q^2+1)\, E_{1214343256}
  \\
  + (-1) q (q-1)^{2} (q+1)^{2} (q^2+1)\, E_{1214354326} + (q-1)^{2} (q+1)^{2} (q^2 - q+1) (q^2+1) (q^2 + q+1)\, E_{1214356432} + (-1) q (q-1)^{2} (q+1)^{2} (q^2 - q+1) (q^2 + q+1)\, E_{1214364325}
  \\
  + (q-1)^{2} (q+1)^{2} q^4\, E_{1214432356} + (-1) (q-1)^{2} (q+1)^{2} q^3 (q^2+1)\, E_{1214433256} + (-1) (q-1)^{2} (q+1)^{2} q^3\, E_{1214543236}
  \\
  + (q-1)^{2} q^2 (q+1)^{2} (q^2+1)\, E_{1214543326} + (q-1)^{2} (q+1)^{2} q^4\, E_{1214564323} + (-1) (q-1)^{2} (q+1)^{2} q^3 (q^2+1)\, E_{1214564332}
  \\
  + (-1) (q-1)^{2} (q+1)^{2} q^5\, E_{1214643235} + (q-1)^{2} (q+1)^{2} q^4 (q^2+1)\, E_{1214643325} + (-1) (q-1)^{2} (q+1)^{2} q^3\, E_{1215432364}
  \\
  + (q-1)^{2} q^2 (q+1)^{2} (q^2+1)\, E_{1215432643} + (q-1)^{2} q^2 (q+1)^{2}\, E_{1215433264} + (-1) q (q-1)^{2} (q+1)^{2} (q^2 - q+1) (q^2 + q+1)\, E_{1215436432}
  \\
  + (q-1)^{2} (q+1)^{2} q^4\, E_{1215643234} + (-1) (q-1)^{2} (q+1)^{2} q^3 (q^2+1)\, E_{1215643243} + (-1) (q-1)^{2} (q+1)^{2} q^3\, E_{1215643324}
  \\
  + (q-1)^{2} q^2 (q+1)^{2} (q^2 - q+1) (q^2 + q+1)\, E_{1215643432} + (q-1)^{2} (q+1)^{2} q^4\, E_{1216432345} + (-1) (q-1)^{2} (q+1)^{2} q^3 (q^2+1)\, E_{1216432354}
  \\
  + (-1) (q-1)^{2} (q+1)^{2} q^3\, E_{1216432435} + (q-1)^{2} q^2 (q+1)^{2} (q^2 - q+1) (q^2 + q+1)\, E_{1216432543} + (-1) (q-1)^{2} (q+1)^{2} q^3\, E_{1216433245}
  \\
  + (q-1)^{2} q^2 (q+1)^{2} (q^2+1)\, E_{1216433254} + (q-1)^{2} q^2 (q+1)^{2}\, E_{1216434325} + (-1) q (q-1)^{2} (q+1)^{2} (q^2+1) (q^4 + 1)\, E_{1216435432}
  \\
  + (q-1)^{2} (q+1)^{2} (q^2+1)\, E_{1232143564} + (-1) q (q-1)^{2} (q+1)^{2}\, E_{1232143645} + ((-1) (q-1)^{2} (q+1)^{2} (q^2 - q+1) (q^2 + q+1))/(q)\, E_{1232145643}
  \\
  + (q-1)^{2} (q+1)^{2}\, E_{1232146435} + (-1) q (q-1)^{2} (q+1)^{2}\, E_{1232154364} + (q-1)^{2} (q+1)^{2} (q^2+1)\, E_{1232154643}
  \\
  + (q-1)^{2} q^2 (q+1)^{2}\, E_{1232156434} + (-1) q (q-1)^{2} (q+1)^{2} (q^2+1)\, E_{1232156443} + (q-1)^{2} q^2 (q+1)^{2}\, E_{1232164345}
  \\
  + (-1) q (q-1)^{2} (q+1)^{2} (q^2+1)\, E_{1232164354} + (-1) q (q-1)^{2} (q+1)^{2}\, E_{1232164435} + (q-1)^{2} (q+1)^{2} (q^2 - q+1) (q^2 + q+1)\, E_{1232164543}
  \\
  + ((-1) (q-1)^{2} (q+1)^{2} (q^2+1))/(q)\, E_{1234321564} + (q-1)^{2} (q+1)^{2}\, E_{1234321645} + ((q-1)^{2} (q+1)^{2} (q^4 - q^3 + q^2 - q+1) (q^4 + q^3 + q^2 + q+1))/(q^2)\, E_{1234564321}
  \\
  + ((-1) (q-1)^{2} (q+1)^{2})/(q)\, E_{1234643215} + (q-1)^{2} (q+1)^{2}\, E_{1235432164} + ((-1) (q-1)^{2} (q+1)^{2} (q^2+1) (q^4 + 1))/(q)\, E_{1235464321}
  \\
  + (-1) q (q-1)^{2} (q+1)^{2}\, E_{1235643214} + (q-1)^{2} (q+1)^{2} (q^2+1) (q^4 + 1)\, E_{1235644321} + (-1) q (q-1)^{2} (q+1)^{2}\, E_{1236432145}
  \\
  + (q-1)^{2} (q+1)^{2} (q^2+1)\, E_{1236432154} + (q-1)^{2} (q+1)^{2}\, E_{1236443215} + ((-1) (q-1)^{2} (q+1)^{2} (q^4 - q^3 + q^2 - q+1) (q^4 + q^3 + q^2 + q+1))/(q)\, E_{1236454321}
  \\
  + (-1) q (q-1)^{2} (q+1)^{2} (q^2+1)\, E_{1243213564} + (q-1)^{2} q^2 (q+1)^{2}\, E_{1243213645} + (q-1)^{2} q^2 (q+1)^{2}\, E_{1243214356}
  \\
  + (-1) q (q-1)^{2} (q+1)^{2}\, E_{1243215436} + (q-1)^{2} (q+1)^{2} (q^2+1)^{2}\, E_{1243215643} + (-1) q (q-1)^{2} (q+1)^{2} (q^2+1)\, E_{1243216435}
  \\
  + (q-1)^{2} (q+1)^{2} (q^2+1)\, E_{1243321564} + (-1) q (q-1)^{2} (q+1)^{2}\, E_{1243321645} + (-1) q (q-1)^{2} (q+1)^{2} (q^2 - q+1) (q^2 + q+1)\, E_{1243432156}
  \\
  + (q-1)^{2} (q+1)^{2} (q^2 - q+1) (q^2 + q+1)\, E_{1243543216} + ((-1) (q-1)^{2} (q+1)^{2} (q^2+1)^{2} (q^4 + 1))/(q)\, E_{1243564321} + (q-1)^{2} (q+1)^{2} (q^2+1) (q^4 + 1)\, E_{1243643215}
  \\
  + (-1) (q-1)^{2} (q+1)^{2} q^3\, E_{1244321356} + (q-1)^{2} q^2 (q+1)^{2} (q^2 - q+1) (q^2 + q+1)\, E_{1244332156} + (q-1)^{2} q^2 (q+1)^{2}\, E_{1245432136}
  \\
  + (-1) q (q-1)^{2} (q+1)^{2} (q^2 - q+1) (q^2 + q+1)\, E_{1245433216} + (-1) (q-1)^{2} (q+1)^{2} q^3\, E_{1245643213} + (q-1)^{2} q^2 (q+1)^{2} (q^2 - q+1) (q^2 + q+1)\, E_{1245643321}
  \\
  + (q-1)^{2} (q+1)^{2} q^4\, E_{1246432135} + (-1) (q-1)^{2} (q+1)^{2} q^3 (q^2 - q+1) (q^2 + q+1)\, E_{1246433215} + (q-1)^{2} q^2 (q+1)^{2}\, E_{1254321364}
  \\
  + (-1) q (q-1)^{2} (q+1)^{2} (q^2+1)\, E_{1254321643} + (-1) q (q-1)^{2} (q+1)^{2}\, E_{1254332164} + (q-1)^{2} (q+1)^{2} (q^2+1) (q^4 + 1)\, E_{1254364321}
  \\
  + (-1) (q-1)^{2} (q+1)^{2} q^3\, E_{1256432134} + (q-1)^{2} q^2 (q+1)^{2} (q^2+1)\, E_{1256432143} + (q-1)^{2} q^2 (q+1)^{2}\, E_{1256433214}
  \\
  + (-1) q (q-1)^{2} (q+1)^{2} (q^2+1) (q^4 + 1)\, E_{1256434321} + (-1) (q-1)^{2} (q+1)^{2} q^3\, E_{1264321345} + (q-1)^{2} q^2 (q+1)^{2} (q^2+1)\, E_{1264321354}
  \\
  + (q-1)^{2} q^2 (q+1)^{2}\, E_{1264321435} + (-1) q (q-1)^{2} (q+1)^{2} (q^2 - q+1) (q^2 + q+1)\, E_{1264321543} + (q-1)^{2} q^2 (q+1)^{2}\, E_{1264332145}
  \\
  + (-1) q (q-1)^{2} (q+1)^{2} (q^2+1)\, E_{1264332154} + (-1) q (q-1)^{2} (q+1)^{2}\, E_{1264343215} + (q-1)^{2} (q+1)^{2} (q^4 - q^3 + q^2 - q+1) (q^4 + q^3 + q^2 + q+1)\, E_{1264354321}
  \\
  + (q-1)^{2} q^2 (q+1)^{2} (q^2+1)\, E_{1321243564} + (-1) (q-1)^{2} (q+1)^{2} q^3\, E_{1321243645} + (-1) q (q-1)^{2} (q+1)^{2} (q^2 - q+1) (q^2 + q+1)\, E_{1321245643}
  \\
  + (q-1)^{2} q^2 (q+1)^{2}\, E_{1321246435} + (-1) (q-1)^{2} (q+1)^{2} q^3\, E_{1321254364} + (q-1)^{2} q^2 (q+1)^{2} (q^2+1)\, E_{1321254643}
  \\
  + (q-1)^{2} (q+1)^{2} q^4\, E_{1321256434} + (-1) (q-1)^{2} (q+1)^{2} q^3 (q^2+1)\, E_{1321256443} + (q-1)^{2} (q+1)^{2} q^4\, E_{1321264345}
  \\
  + (-1) (q-1)^{2} (q+1)^{2} q^3 (q^2+1)\, E_{1321264354} + (-1) (q-1)^{2} (q+1)^{2} q^3\, E_{1321264435} + (q-1)^{2} q^2 (q+1)^{2} (q^2 - q+1) (q^2 + q+1)\, E_{1321264543}
  \\
  + (q-1)^{2} (q+1)^{2} q^4\, E_{1321324564} + (-1) (q-1)^{2} (q+1)^{2} q^3\, E_{1321325464} + (q-1)^{2} (q+1)^{2} q^4\, E_{1321325644}
  \\
  + (-1) (q-1)^{2} (q+1)^{2} q^5\, E_{1321326454} + (-1) (q-1)^{2} (q+1)^{2} q^3\, E_{1321432456} + (q-1)^{2} q^2 (q+1)^{2}\, E_{1321432546}
  \\
  + (-1) q (q-1)^{2} (q+1)^{2} (q^2+1)^{2}\, E_{1321432564} + (q-1)^{2} q^2 (q+1)^{2} (q^2+1)\, E_{1321432645} + (q-1)^{2} (q+1)^{2} q^4\, E_{1321443256}
  \\
  + (-1) (q-1)^{2} (q+1)^{2} q^3\, E_{1321454326} + (q-1)^{2} (q+1)^{2} (q^6 + 2 q^4 + q^2+1)\, E_{1321456432} + (-1) q (q-1)^{2} (q+1)^{2} (q^4 + 1)\, E_{1321464325}
  \\
  + (q-1)^{2} q^2 (q+1)^{2} (q^2+1)\, E_{1321543264} + (-1) q (q-1)^{2} (q+1)^{2} (q^2 - q+1) (q^2 + q+1)\, E_{1321546432} + (-1) (q-1)^{2} (q+1)^{2} q^3 (q^2+1)\, E_{1321564324}
  \\
  + (q-1)^{2} q^2 (q+1)^{2} (q^2 - q+1) (q^2 + q+1)\, E_{1321564432} + (-1) (q-1)^{2} (q+1)^{2} q^3\, E_{1321643245} + (q-1)^{2} q^2 (q+1)^{2} (q^2 - q+1) (q^2 + q+1)\, E_{1321643254}
  \\
  + (q-1)^{2} q^2 (q+1)^{2}\, E_{1321644325} + (-1) q (q-1)^{2} (q+1)^{2} (q^2+1) (q^4 + 1)\, E_{1321645432} + (-1) q (q-1)^{2} (q+1)^{2} (q^2+1)\, E_{1322143564}
  \\
  + (q-1)^{2} q^2 (q+1)^{2}\, E_{1322143645} + (q-1)^{2} (q+1)^{2} (q^2 - q+1) (q^2 + q+1)\, E_{1322145643} + (-1) q (q-1)^{2} (q+1)^{2}\, E_{1322146435}
  \\
  + (q-1)^{2} q^2 (q+1)^{2}\, E_{1322154364} + (-1) q (q-1)^{2} (q+1)^{2} (q^2+1)\, E_{1322154643} + (-1) (q-1)^{2} (q+1)^{2} q^3\, E_{1322156434}
  \\
  + (q-1)^{2} q^2 (q+1)^{2} (q^2+1)\, E_{1322156443} + (-1) (q-1)^{2} (q+1)^{2} q^3\, E_{1322164345} + (q-1)^{2} q^2 (q+1)^{2} (q^2+1)\, E_{1322164354}
  \\
  + (q-1)^{2} q^2 (q+1)^{2}\, E_{1322164435} + (-1) q (q-1)^{2} (q+1)^{2} (q^2 - q+1) (q^2 + q+1)\, E_{1322164543} + (-1) (q-1)^{2} (q+1)^{2} q^3 (q^2+1)\, E_{1323214564}
  \\
  + (q-1)^{2} q^2 (q+1)^{2} (q^2+1)\, E_{1323215464} + (-1) (q-1)^{2} (q+1)^{2} q^3 (q^2+1)\, E_{1323215644} + (q-1)^{2} (q+1)^{2} q^4 (q^2+1)\, E_{1323216454}
  \\
  + (q-1)^{2} q^2 (q+1)^{2} (q^2+1)\, E_{1324321456} + (-1) q (q-1)^{2} (q+1)^{2} (q^2+1)\, E_{1324321546} + (q-1)^{2} (q+1)^{2} (q^2 - q+1) (q^2+1) (q^2 + q+1)\, E_{1324321564}
  \\
  + (-1) q (q-1)^{2} (q+1)^{2} (q^2 - q+1) (q^2 + q+1)\, E_{1324321645} + (-1) (q-1)^{2} (q+1)^{2} q^3 (q^2+1)\, E_{1324432156} + (q-1)^{2} q^2 (q+1)^{2} (q^2+1)\, E_{1324543216}
  \\
  + ((-1) (q-1)^{2} (q+1)^{2} (q^8 + 2 q^6 + 2 q^4 + q^2+1))/(q)\, E_{1324564321} + (q-1)^{2} (q+1)^{2} (q^6 + q^4 + 1)\, E_{1324643215} + (-1) q (q-1)^{2} (q+1)^{2} (q^2 - q+1) (q^2 + q+1)\, E_{1325432164}
  \\
  + (q-1)^{2} (q+1)^{2} (q^2+1) (q^4 + 1)\, E_{1325464321} + (q-1)^{2} q^2 (q+1)^{2} (q^2 - q+1) (q^2 + q+1)\, E_{1325643214} + (-1) q (q-1)^{2} (q+1)^{2} (q^2+1) (q^4 + 1)\, E_{1325644321}
  \\
  + (q-1)^{2} q^2 (q+1)^{2}\, E_{1326432145} + (-1) q (q-1)^{2} (q+1)^{2} (q^2+1) (q^4 + 1)\, E_{1326432154} + (-1) q (q-1)^{2} (q+1)^{2}\, E_{1326443215}
  \\
  + (q-1)^{2} (q+1)^{2} (q^4 - q^3 + q^2 - q+1) (q^4 + q^3 + q^2 + q+1)\, E_{1326454321} + (-1) (q-1)^{2} (q+1)^{2} q^5\, E_{1332124564} + (q-1)^{2} (q+1)^{2} q^4\, E_{1332125464}
  \\
  + (-1) (q-1)^{2} (q+1)^{2} q^5\, E_{1332125644} + (q-1)^{2} (q+1)^{2} q^6\, E_{1332126454} + (q-1)^{2} (q+1)^{2} q^4 (q^2+1)\, E_{1332214564}
  \\
  + (-1) (q-1)^{2} (q+1)^{2} q^3 (q^2+1)\, E_{1332215464} + (q-1)^{2} (q+1)^{2} q^4 (q^2+1)\, E_{1332215644} + (-1) (q-1)^{2} (q+1)^{2} q^5 (q^2+1)\, E_{1332216454}
  \\
  + (q-1)^{2} (q+1)^{2} q^4\, E_{1343212456} + (-1) (q-1)^{2} (q+1)^{2} q^3\, E_{1343212546} + (q-1)^{2} (q+1)^{2} q^4 (q^2+1)\, E_{1343212564}
  \\
  + (-1) (q-1)^{2} (q+1)^{2} q^5\, E_{1343212645} + (-1) (q-1)^{2} (q+1)^{2} q^3 (q^2+1)\, E_{1343221456} + (q-1)^{2} q^2 (q+1)^{2} (q^2+1)\, E_{1343221546}
  \\
  + (-1) (q-1)^{2} (q+1)^{2} q^3 (q^2+1)^{2}\, E_{1343221564} + (q-1)^{2} (q+1)^{2} q^4 (q^2+1)\, E_{1343221645} + (-1) (q-1)^{2} (q+1)^{2} q^5\, E_{1344321256}
  \\
  + (q-1)^{2} (q+1)^{2} q^4 (q^2+1)\, E_{1344322156} + (q-1)^{2} (q+1)^{2} q^4\, E_{1345432126} + (-1) (q-1)^{2} (q+1)^{2} q^3 (q^2+1)\, E_{1345432216}
  \\
  + (-1) (q-1)^{2} (q+1)^{2} q^5\, E_{1345643212} + (q-1)^{2} (q+1)^{2} q^4 (q^2+1)\, E_{1345643221} + (q-1)^{2} (q+1)^{2} q^6\, E_{1346432125}
  \\
  + (-1) (q-1)^{2} (q+1)^{2} q^5 (q^2+1)\, E_{1346432215} + (-1) (q-1)^{2} (q+1)^{2} q^5\, E_{1354321264} + (q-1)^{2} (q+1)^{2} q^4 (q^2+1)\, E_{1354322164}
  \\
  + (q-1)^{2} (q+1)^{2} q^6\, E_{1356432124} + (-1) (q-1)^{2} (q+1)^{2} q^5 (q^2+1)\, E_{1356432214} + (-1) (q-1)^{2} (q+1)^{2} q^7\, E_{1364321254}
  \\
  + (q-1)^{2} (q+1)^{2} q^6 (q^2+1)\, E_{1364322154} + (-1) (q-1)^{2} (q+1)^{2} q^3 (q^2+1)\, E_{1432123564} + (q-1)^{2} (q+1)^{2} q^4\, E_{1432123645}
  \\
  + (q-1)^{2} (q+1)^{2} q^4\, E_{1432124356} + (-1) (q-1)^{2} (q+1)^{2} q^3\, E_{1432125436} + (q-1)^{2} q^2 (q+1)^{2} (q^2+1)^{2}\, E_{1432125643}
  \\
  + (-1) (q-1)^{2} (q+1)^{2} q^3 (q^2+1)\, E_{1432126435} + (q-1)^{2} q^2 (q+1)^{2} (q^2+1)\, E_{1432132564} + (-1) (q-1)^{2} (q+1)^{2} q^3\, E_{1432132645}
  \\
  + (-1) (q-1)^{2} (q+1)^{2} q^3 (q^2+1)\, E_{1432143256} + (q-1)^{2} q^2 (q+1)^{2} (q^2+1)\, E_{1432154326} + (-1) q (q-1)^{2} (q+1)^{2} (q^2 - q+1) (q^2+1) (q^2 + q+1)\, E_{1432156432}
  \\
  + (q-1)^{2} q^2 (q+1)^{2} (q^2 - q+1) (q^2 + q+1)\, E_{1432164325} + (q-1)^{2} q^2 (q+1)^{2} (q^2+1)\, E_{1432213564} + (-1) (q-1)^{2} (q+1)^{2} q^3\, E_{1432213645}
  \\
  + (-1) (q-1)^{2} (q+1)^{2} q^3\, E_{1432214356} + (q-1)^{2} q^2 (q+1)^{2}\, E_{1432215436} + (-1) q (q-1)^{2} (q+1)^{2} (q^2+1)^{2}\, E_{1432215643}
  \\
  + (q-1)^{2} q^2 (q+1)^{2} (q^2+1)\, E_{1432216435} + (-1) q (q-1)^{2} (q+1)^{2} (q^2+1)\, E_{1432321564} + (q-1)^{2} q^2 (q+1)^{2}\, E_{1432321645}
  \\
  + (q-1)^{2} q^2 (q+1)^{2} (q^2 - q+1) (q^2 + q+1)\, E_{1432432156} + (-1) q (q-1)^{2} (q+1)^{2} (q^2 - q+1) (q^2 + q+1)\, E_{1432543216} + (q-1)^{2} (q+1)^{2} (q^2+1)^{2} (q^4 + 1)\, E_{1432564321}
  \\
  + (-1) q (q-1)^{2} (q+1)^{2} (q^2+1) (q^4 + 1)\, E_{1432643215} + (-1) (q-1)^{2} (q+1)^{2} q^5\, E_{1443212356} + (q-1)^{2} (q+1)^{2} q^4 (q^2+1)\, E_{1443213256}
  \\
  + (q-1)^{2} (q+1)^{2} q^4\, E_{1443221356} + (-1) (q-1)^{2} (q+1)^{2} q^3 (q^2 - q+1) (q^2 + q+1)\, E_{1443232156} + (q-1)^{2} (q+1)^{2} q^4\, E_{1454321236}
  \\
  + (-1) (q-1)^{2} (q+1)^{2} q^3 (q^2+1)\, E_{1454321326} + (-1) (q-1)^{2} (q+1)^{2} q^3\, E_{1454322136} + (q-1)^{2} q^2 (q+1)^{2} (q^2 - q+1) (q^2 + q+1)\, E_{1454323216}
  \\
  + (-1) (q-1)^{2} (q+1)^{2} q^5\, E_{1456432123} + (q-1)^{2} (q+1)^{2} q^4 (q^2+1)\, E_{1456432132} + (q-1)^{2} (q+1)^{2} q^4\, E_{1456432213}
  \\
  + (-1) (q-1)^{2} (q+1)^{2} q^3 (q^2 - q+1) (q^2 + q+1)\, E_{1456432321} + (q-1)^{2} (q+1)^{2} q^6\, E_{1464321235} + (-1) (q-1)^{2} (q+1)^{2} q^5 (q^2+1)\, E_{1464321325}
  \\
  + (-1) (q-1)^{2} (q+1)^{2} q^5\, E_{1464322135} + (q-1)^{2} (q+1)^{2} q^4 (q^2 - q+1) (q^2 + q+1)\, E_{1464323215} + (q-1)^{2} (q+1)^{2} q^4\, E_{1543212364}
  \\
  + (-1) (q-1)^{2} (q+1)^{2} q^3 (q^2+1)\, E_{1543212643} + (-1) (q-1)^{2} (q+1)^{2} q^3\, E_{1543213264} + (q-1)^{2} q^2 (q+1)^{2} (q^2 - q+1) (q^2 + q+1)\, E_{1543216432}
  \\
  + (-1) (q-1)^{2} (q+1)^{2} q^3\, E_{1543221364} + (q-1)^{2} q^2 (q+1)^{2} (q^2+1)\, E_{1543221643} + (q-1)^{2} q^2 (q+1)^{2}\, E_{1543232164}
  \\
  + (-1) q (q-1)^{2} (q+1)^{2} (q^2+1) (q^4 + 1)\, E_{1543264321} + (-1) (q-1)^{2} (q+1)^{2} q^5\, E_{1564321234} + (q-1)^{2} (q+1)^{2} q^4 (q^2+1)\, E_{1564321243}
  \\
  + (q-1)^{2} (q+1)^{2} q^4\, E_{1564321324} + (-1) (q-1)^{2} (q+1)^{2} q^3 (q^2 - q+1) (q^2 + q+1)\, E_{1564321432} + (q-1)^{2} (q+1)^{2} q^4\, E_{1564322134}
  \\
  + (-1) (q-1)^{2} (q+1)^{2} q^3 (q^2+1)\, E_{1564322143} + (-1) (q-1)^{2} (q+1)^{2} q^3\, E_{1564323214} + (q-1)^{2} q^2 (q+1)^{2} (q^2+1) (q^4 + 1)\, E_{1564324321}
  \\
  + (-1) (q-1)^{2} (q+1)^{2} q^5\, E_{1643212345} + (q-1)^{2} (q+1)^{2} q^4 (q^2+1)\, E_{1643212354} + (q-1)^{2} (q+1)^{2} q^4\, E_{1643212435}
  \\
  + (-1) (q-1)^{2} (q+1)^{2} q^3 (q^2 - q+1) (q^2 + q+1)\, E_{1643212543} + (q-1)^{2} (q+1)^{2} q^4\, E_{1643213245} + (-1) (q-1)^{2} (q+1)^{2} q^3 (q^2+1)\, E_{1643213254}
  \\
  + (-1) (q-1)^{2} (q+1)^{2} q^3\, E_{1643214325} + (q-1)^{2} q^2 (q+1)^{2} (q^2+1) (q^4 + 1)\, E_{1643215432} + (q-1)^{2} (q+1)^{2} q^4\, E_{1643221345}
  \\
  + (-1) (q-1)^{2} (q+1)^{2} q^3 (q^2+1)\, E_{1643221354} + (-1) (q-1)^{2} (q+1)^{2} q^3\, E_{1643221435} + (q-1)^{2} q^2 (q+1)^{2} (q^2 - q+1) (q^2 + q+1)\, E_{1643221543}
  \\
  + (-1) (q-1)^{2} (q+1)^{2} q^3\, E_{1643232145} + (q-1)^{2} q^2 (q+1)^{2} (q^2+1)\, E_{1643232154} + (q-1)^{2} q^2 (q+1)^{2}\, E_{1643243215}
  \\
  + (-1) q (q-1)^{2} (q+1)^{2} (q^4 - q^3 + q^2 - q+1) (q^4 + q^3 + q^2 + q+1)\, E_{1643254321} + (q-1)^{2} q^2 (q+1)^{2} (q^2+1)\, E_{2113243564} + (-1) (q-1)^{2} (q+1)^{2} q^3\, E_{2113243645}
  \\
  + (-1) q (q-1)^{2} (q+1)^{2} (q^2 - q+1) (q^2 + q+1)\, E_{2113245643} + (q-1)^{2} q^2 (q+1)^{2}\, E_{2113246435} + (-1) (q-1)^{2} (q+1)^{2} q^3\, E_{2113254364}
  \\
  + (q-1)^{2} q^2 (q+1)^{2} (q^2+1)\, E_{2113254643} + (q-1)^{2} (q+1)^{2} q^4\, E_{2113256434} + (-1) (q-1)^{2} (q+1)^{2} q^3 (q^2+1)\, E_{2113256443}
  \\
  + (q-1)^{2} (q+1)^{2} q^4\, E_{2113264345} + (-1) (q-1)^{2} (q+1)^{2} q^3 (q^2+1)\, E_{2113264354} + (-1) (q-1)^{2} (q+1)^{2} q^3\, E_{2113264435}
  \\
  + (q-1)^{2} q^2 (q+1)^{2} (q^2 - q+1) (q^2 + q+1)\, E_{2113264543} + (-1) q (q-1)^{2} (q+1)^{2} (q^2+1)\, E_{2113432564} + (q-1)^{2} q^2 (q+1)^{2}\, E_{2113432645}
  \\
  + (q-1)^{2} (q+1)^{2} (q^2+1) (q^4 + 1)\, E_{2113456432} + (-1) q (q-1)^{2} (q+1)^{2}\, E_{2113464325} + (q-1)^{2} q^2 (q+1)^{2}\, E_{2113543264}
  \\
  + (-1) q (q-1)^{2} (q+1)^{2} (q^2 - q+1) (q^2 + q+1)\, E_{2113546432} + (-1) (q-1)^{2} (q+1)^{2} q^3\, E_{2113564324} + (q-1)^{2} q^2 (q+1)^{2} (q^2 - q+1) (q^2 + q+1)\, E_{2113564432}
  \\
  + (-1) (q-1)^{2} (q+1)^{2} q^3\, E_{2113643245} + (q-1)^{2} q^2 (q+1)^{2} (q^2+1)\, E_{2113643254} + (q-1)^{2} q^2 (q+1)^{2}\, E_{2113644325}
  \\
  + (-1) q (q-1)^{2} (q+1)^{2} (q^2+1) (q^4 + 1)\, E_{2113645432} + (-1) (q-1)^{2} (q+1)^{2} q^3 (q^2+1)\, E_{2114323564} + (q-1)^{2} (q+1)^{2} q^4\, E_{2114323645}
  \\
  + (q-1)^{2} (q+1)^{2} q^4\, E_{2114324356} + (-1) (q-1)^{2} (q+1)^{2} q^3\, E_{2114325436} + (q-1)^{2} q^2 (q+1)^{2} (q^2+1)^{2}\, E_{2114325643}
  \\
  + (-1) (q-1)^{2} (q+1)^{2} q^3 (q^2+1)\, E_{2114326435} + (q-1)^{2} q^2 (q+1)^{2} (q^2+1)\, E_{2114332564} + (-1) (q-1)^{2} (q+1)^{2} q^3\, E_{2114332645}
  \\
  + (-1) (q-1)^{2} (q+1)^{2} q^3 (q^2+1)\, E_{2114343256} + (q-1)^{2} q^2 (q+1)^{2} (q^2+1)\, E_{2114354326} + (-1) q (q-1)^{2} (q+1)^{2} (q^2 - q+1) (q^2+1) (q^2 + q+1)\, E_{2114356432}
  \\
  + (q-1)^{2} q^2 (q+1)^{2} (q^2 - q+1) (q^2 + q+1)\, E_{2114364325} + (-1) (q-1)^{2} (q+1)^{2} q^5\, E_{2114432356} + (q-1)^{2} (q+1)^{2} q^4 (q^2+1)\, E_{2114433256}
  \\
  + (q-1)^{2} (q+1)^{2} q^4\, E_{2114543236} + (-1) (q-1)^{2} (q+1)^{2} q^3 (q^2+1)\, E_{2114543326} + (-1) (q-1)^{2} (q+1)^{2} q^5\, E_{2114564323}
  \\
  + (q-1)^{2} (q+1)^{2} q^4 (q^2+1)\, E_{2114564332} + (q-1)^{2} (q+1)^{2} q^6\, E_{2114643235} + (-1) (q-1)^{2} (q+1)^{2} q^5 (q^2+1)\, E_{2114643325}
  \\
  + (q-1)^{2} (q+1)^{2} q^4\, E_{2115432364} + (-1) (q-1)^{2} (q+1)^{2} q^3 (q^2+1)\, E_{2115432643} + (-1) (q-1)^{2} (q+1)^{2} q^3\, E_{2115433264}
  \\
  + (q-1)^{2} q^2 (q+1)^{2} (q^2 - q+1) (q^2 + q+1)\, E_{2115436432} + (-1) (q-1)^{2} (q+1)^{2} q^5\, E_{2115643234} + (q-1)^{2} (q+1)^{2} q^4 (q^2+1)\, E_{2115643243}
  \\
  + (q-1)^{2} (q+1)^{2} q^4\, E_{2115643324} + (-1) (q-1)^{2} (q+1)^{2} q^3 (q^2 - q+1) (q^2 + q+1)\, E_{2115643432} + (-1) (q-1)^{2} (q+1)^{2} q^5\, E_{2116432345}
  \\
  + (q-1)^{2} (q+1)^{2} q^4 (q^2+1)\, E_{2116432354} + (q-1)^{2} (q+1)^{2} q^4\, E_{2116432435} + (-1) (q-1)^{2} (q+1)^{2} q^3 (q^2 - q+1) (q^2 + q+1)\, E_{2116432543}
  \\
  + (q-1)^{2} (q+1)^{2} q^4\, E_{2116433245} + (-1) (q-1)^{2} (q+1)^{2} q^3 (q^2+1)\, E_{2116433254} + (-1) (q-1)^{2} (q+1)^{2} q^3\, E_{2116434325}
  \\
  + (q-1)^{2} q^2 (q+1)^{2} (q^2+1) (q^4 + 1)\, E_{2116435432} + (-1) (q-1)^{2} (q+1)^{2} q^5\, E_{2121345643} + (q-1)^{2} (q+1)^{2} q^4\, E_{2121354643}
  \\
  + (-1) (q-1)^{2} (q+1)^{2} q^5\, E_{2121356443} + (q-1)^{2} (q+1)^{2} q^6\, E_{2121364543} + (q-1)^{2} (q+1)^{2} q^4\, E_{2121434356}
  \\
  + (-1) (q-1)^{2} (q+1)^{2} q^3\, E_{2121435436} + (q-1)^{2} (q+1)^{2} q^4 (q^2+1)\, E_{2121435643} + (-1) (q-1)^{2} (q+1)^{2} q^5\, E_{2121436435}
  \\
  + (-1) (q-1)^{2} (q+1)^{2} q^5\, E_{2121443356} + (q-1)^{2} (q+1)^{2} q^4\, E_{2121454336} + (-1) (q-1)^{2} (q+1)^{2} q^5\, E_{2121456433}
  \\
  + (q-1)^{2} (q+1)^{2} q^6\, E_{2121464335} + (-1) (q-1)^{2} (q+1)^{2} q^5\, E_{2121543643} + (q-1)^{2} (q+1)^{2} q^6\, E_{2121564343}
  \\
  + (-1) (q-1)^{2} (q+1)^{2} q^7\, E_{2121643543} + (q-1)^{2} (q+1)^{2} q^4\, E_{2132134564} + (-1) (q-1)^{2} (q+1)^{2} q^3\, E_{2132135464}
  \\
  + (q-1)^{2} (q+1)^{2} q^4\, E_{2132135644} + (-1) (q-1)^{2} (q+1)^{2} q^5\, E_{2132136454} + (-1) (q-1)^{2} (q+1)^{2} q^3\, E_{2132143456}
  \\
  + (q-1)^{2} q^2 (q+1)^{2}\, E_{2132143546} + (-1) q (q-1)^{2} (q+1)^{2} (q^2+1)^{2}\, E_{2132143564} + (q-1)^{2} q^2 (q+1)^{2} (q^2+1)\, E_{2132143645}
  \\
  + (q-1)^{2} (q+1)^{2} q^4\, E_{2132144356} + (-1) (q-1)^{2} (q+1)^{2} q^3\, E_{2132145436} + (q-1)^{2} (q+1)^{2} (q^6 + 2 q^4 + q^2+1)\, E_{2132145643}
  \\
  + (-1) q (q-1)^{2} (q+1)^{2} (q^4 + 1)\, E_{2132146435} + (q-1)^{2} q^2 (q+1)^{2} (q^2+1)\, E_{2132154364} + (-1) q (q-1)^{2} (q+1)^{2} (q^2 - q+1) (q^2 + q+1)\, E_{2132154643}
  \\
  + (-1) (q-1)^{2} (q+1)^{2} q^3 (q^2+1)\, E_{2132156434} + (q-1)^{2} q^2 (q+1)^{2} (q^2 - q+1) (q^2 + q+1)\, E_{2132156443} + (-1) (q-1)^{2} (q+1)^{2} q^3\, E_{2132164345}
  \\
  + (q-1)^{2} q^2 (q+1)^{2} (q^2 - q+1) (q^2 + q+1)\, E_{2132164354} + (q-1)^{2} q^2 (q+1)^{2}\, E_{2132164435} + (-1) q (q-1)^{2} (q+1)^{2} (q^2+1) (q^4 + 1)\, E_{2132164543}
  \\
  + (-1) (q-1)^{2} (q+1)^{2} q^5\, E_{2133214564} + (q-1)^{2} (q+1)^{2} q^4\, E_{2133215464} + (-1) (q-1)^{2} (q+1)^{2} q^5\, E_{2133215644}
  \\
  + (q-1)^{2} (q+1)^{2} q^6\, E_{2133216454} + (q-1)^{2} (q+1)^{2} q^4\, E_{2134321456} + (-1) (q-1)^{2} (q+1)^{2} q^3\, E_{2134321546}
  \\
  + (q-1)^{2} (q+1)^{2} (q^2+1) (q^4 + 1)\, E_{2134321564} + (-1) q (q-1)^{2} (q+1)^{2} (q^4 + 1)\, E_{2134321645} + (-1) (q-1)^{2} (q+1)^{2} q^5\, E_{2134432156}
  \\
  + (q-1)^{2} (q+1)^{2} q^4\, E_{2134543216} + ((-1) (q-1)^{2} (q+1)^{2} (q^2+1) (q^6 + q^4 + 1))/(q)\, E_{2134564321} + (q-1)^{2} (q+1)^{2} (q^2+1) (q^4 - q^2+1)\, E_{2134643215}
  \\
  + (-1) q (q-1)^{2} (q+1)^{2} (q^4 + 1)\, E_{2135432164} + (q-1)^{2} (q+1)^{2} (q^2+1) (q^4 + 1)\, E_{2135464321} + (q-1)^{2} q^2 (q+1)^{2} (q^4 + 1)\, E_{2135643214}
  \\
  + (-1) q (q-1)^{2} (q+1)^{2} (q^2+1) (q^4 + 1)\, E_{2135644321} + (q-1)^{2} q^2 (q+1)^{2}\, E_{2136432145} + (-1) q (q-1)^{2} (q+1)^{2} (q^6 + q^2+1)\, E_{2136432154}
  \\
  + (-1) q (q-1)^{2} (q+1)^{2}\, E_{2136443215} + (q-1)^{2} (q+1)^{2} (q^4 - q^3 + q^2 - q+1) (q^4 + q^3 + q^2 + q+1)\, E_{2136454321} + (q-1)^{2} q^2 (q+1)^{2} (q^2+1)\, E_{2143213564}
  \\
  + (-1) (q-1)^{2} (q+1)^{2} q^3\, E_{2143213645} + (-1) (q-1)^{2} (q+1)^{2} q^3 (q^2+1)\, E_{2143214356} + (q-1)^{2} q^2 (q+1)^{2} (q^2+1)\, E_{2143215436}
  \\
  + (-1) q (q-1)^{2} (q+1)^{2} (q^2 - q+1) (q^2+1) (q^2 + q+1)\, E_{2143215643} + (q-1)^{2} q^2 (q+1)^{2} (q^2 - q+1) (q^2 + q+1)\, E_{2143216435} + (-1) q (q-1)^{2} (q+1)^{2} (q^2+1)\, E_{2143321564}
  \\
  + (q-1)^{2} q^2 (q+1)^{2}\, E_{2143321645} + (q-1)^{2} q^2 (q+1)^{2} (q^2 - q+1) (q^2 + q+1)\, E_{2143432156} + (-1) q (q-1)^{2} (q+1)^{2} (q^2 - q+1) (q^2 + q+1)\, E_{2143543216}
  \\
  + (q-1)^{2} (q+1)^{2} (q^2+1)^{2} (q^4 + 1)\, E_{2143564321} + (-1) q (q-1)^{2} (q+1)^{2} (q^2+1) (q^4 + 1)\, E_{2143643215} + (q-1)^{2} (q+1)^{2} q^4 (q^2+1)\, E_{2144321356}
  \\
  + (-1) (q-1)^{2} (q+1)^{2} q^3 (q^2 - q+1) (q^2 + q+1)\, E_{2144332156} + (-1) (q-1)^{2} (q+1)^{2} q^3 (q^2+1)\, E_{2145432136} + (q-1)^{2} q^2 (q+1)^{2} (q^2 - q+1) (q^2 + q+1)\, E_{2145433216}
  \\
  + (q-1)^{2} (q+1)^{2} q^4 (q^2+1)\, E_{2145643213} + (-1) (q-1)^{2} (q+1)^{2} q^3 (q^2 - q+1) (q^2 + q+1)\, E_{2145643321} + (-1) (q-1)^{2} (q+1)^{2} q^5 (q^2+1)\, E_{2146432135}
  \\
  + (q-1)^{2} (q+1)^{2} q^4 (q^2 - q+1) (q^2 + q+1)\, E_{2146433215} + (-1) (q-1)^{2} (q+1)^{2} q^3\, E_{2154321364} + (q-1)^{2} q^2 (q+1)^{2} (q^2 - q+1) (q^2 + q+1)\, E_{2154321643}
  \\
  + (q-1)^{2} q^2 (q+1)^{2}\, E_{2154332164} + (-1) q (q-1)^{2} (q+1)^{2} (q^2+1) (q^4 + 1)\, E_{2154364321} + (q-1)^{2} (q+1)^{2} q^4\, E_{2156432134}
  \\
  + (-1) (q-1)^{2} (q+1)^{2} q^3 (q^2 - q+1) (q^2 + q+1)\, E_{2156432143} + (-1) (q-1)^{2} (q+1)^{2} q^3\, E_{2156433214} + (q-1)^{2} q^2 (q+1)^{2} (q^2+1) (q^4 + 1)\, E_{2156434321}
  \\
  + (q-1)^{2} (q+1)^{2} q^4\, E_{2164321345} + (-1) (q-1)^{2} (q+1)^{2} q^3 (q^2+1)\, E_{2164321354} + (-1) (q-1)^{2} (q+1)^{2} q^3\, E_{2164321435}
  \\
  + (q-1)^{2} q^2 (q+1)^{2} (q^2+1) (q^4 + 1)\, E_{2164321543} + (-1) (q-1)^{2} (q+1)^{2} q^3\, E_{2164332145} + (q-1)^{2} q^2 (q+1)^{2} (q^2+1)\, E_{2164332154}
  \\
  + (q-1)^{2} q^2 (q+1)^{2}\, E_{2164343215} + (-1) q (q-1)^{2} (q+1)^{2} (q^4 - q^3 + q^2 - q+1) (q^4 + q^3 + q^2 + q+1)\, E_{2164354321} + (q-1)^{2} (q+1)^{2} q^6\, E_{2211345643}
  \\
  + (-1) (q-1)^{2} (q+1)^{2} q^5\, E_{2211354643} + (q-1)^{2} (q+1)^{2} q^6\, E_{2211356443} + (-1) (q-1)^{2} (q+1)^{2} q^7\, E_{2211364543}
  \\
  + (-1) (q-1)^{2} (q+1)^{2} q^5\, E_{2211434356} + (q-1)^{2} (q+1)^{2} q^4\, E_{2211435436} + (-1) (q-1)^{2} (q+1)^{2} q^5 (q^2+1)\, E_{2211435643}
  \\
  + (q-1)^{2} (q+1)^{2} q^6\, E_{2211436435} + (q-1)^{2} (q+1)^{2} q^6\, E_{2211443356} + (-1) (q-1)^{2} (q+1)^{2} q^5\, E_{2211454336}
  \\
  + (q-1)^{2} (q+1)^{2} q^6\, E_{2211456433} + (-1) (q-1)^{2} (q+1)^{2} q^7\, E_{2211464335} + (q-1)^{2} (q+1)^{2} q^6\, E_{2211543643}
  \\
  + (-1) (q-1)^{2} (q+1)^{2} q^7\, E_{2211564343} + (q-1)^{2} (q+1)^{2} q^8\, E_{2211643543} + (-1) (q-1)^{2} (q+1)^{2} q^5\, E_{2321134564}
  \\
  + (q-1)^{2} (q+1)^{2} q^4\, E_{2321135464} + (-1) (q-1)^{2} (q+1)^{2} q^5\, E_{2321135644} + (q-1)^{2} (q+1)^{2} q^6\, E_{2321136454}
  \\
  + (q-1)^{2} (q+1)^{2} q^4\, E_{2321143456} + (-1) (q-1)^{2} (q+1)^{2} q^3\, E_{2321143546} + (q-1)^{2} (q+1)^{2} q^4 (q^2+1)\, E_{2321143564}
  \\
  + (-1) (q-1)^{2} (q+1)^{2} q^5\, E_{2321143645} + (-1) (q-1)^{2} (q+1)^{2} q^5\, E_{2321144356} + (q-1)^{2} (q+1)^{2} q^4\, E_{2321145436}
  \\
  + (-1) (q-1)^{2} (q+1)^{2} q^5 (q^2+1)\, E_{2321145643} + (q-1)^{2} (q+1)^{2} q^6\, E_{2321146435} + (-1) (q-1)^{2} (q+1)^{2} q^5\, E_{2321154364}
  \\
  + (q-1)^{2} (q+1)^{2} q^6\, E_{2321154643} + (q-1)^{2} (q+1)^{2} q^6\, E_{2321156434} + (-1) (q-1)^{2} (q+1)^{2} q^7\, E_{2321156443}
  \\
  + (-1) (q-1)^{2} (q+1)^{2} q^7\, E_{2321164354} + (q-1)^{2} (q+1)^{2} q^8\, E_{2321164543} + (q-1)^{2} (q+1)^{2} q^6\, E_{2332114564}
  \\
  + (-1) (q-1)^{2} (q+1)^{2} q^5\, E_{2332115464} + (q-1)^{2} (q+1)^{2} q^6\, E_{2332115644} + (-1) (q-1)^{2} (q+1)^{2} q^7\, E_{2332116454}
  \\
  + (-1) (q-1)^{2} (q+1)^{2} q^5\, E_{2343211456} + (q-1)^{2} (q+1)^{2} q^4\, E_{2343211546} + (-1) (q-1)^{2} (q+1)^{2} q^5 (q^2+1)\, E_{2343211564}
  \\
  + (q-1)^{2} (q+1)^{2} q^6\, E_{2343211645} + (q-1)^{2} (q+1)^{2} q^6\, E_{2344321156} + (-1) (q-1)^{2} (q+1)^{2} q^5\, E_{2345432116}
  \\
  + (q-1)^{2} (q+1)^{2} q^6\, E_{2345643211} + (-1) (q-1)^{2} (q+1)^{2} q^7\, E_{2346432115} + (q-1)^{2} (q+1)^{2} q^6\, E_{2354321164}
  \\
  + (-1) (q-1)^{2} (q+1)^{2} q^7\, E_{2356432114} + (q-1)^{2} (q+1)^{2} q^8\, E_{2364321154} + (q-1)^{2} (q+1)^{2} q^6\, E_{2432114356}
  \\
  + (-1) (q-1)^{2} (q+1)^{2} q^5\, E_{2432115436} + (q-1)^{2} (q+1)^{2} q^6 (q^2+1)\, E_{2432115643} + (-1) (q-1)^{2} (q+1)^{2} q^7\, E_{2432116435}
  \\
  + (-1) (q-1)^{2} (q+1)^{2} q^7\, E_{2443211356} + (q-1)^{2} (q+1)^{2} q^6\, E_{2454321136} + (-1) (q-1)^{2} (q+1)^{2} q^7\, E_{2456432113}
  \\
  + (q-1)^{2} (q+1)^{2} q^8\, E_{2464321135} + (-1) (q-1)^{2} (q+1)^{2} q^7\, E_{2543211643} + (q-1)^{2} (q+1)^{2} q^8\, E_{2564321143}
  \\
  + (-1) (q-1)^{2} (q+1)^{2} q^9\, E_{2643211543} + (-1) (q-1)^{2} (q+1)^{2} q^3 (q^2+1)\, E_{3211243564} + (q-1)^{2} (q+1)^{2} q^4\, E_{3211243645}
  \\
  + (q-1)^{2} q^2 (q+1)^{2} (q^2 - q+1) (q^2 + q+1)\, E_{3211245643} + (-1) (q-1)^{2} (q+1)^{2} q^3\, E_{3211246435} + (q-1)^{2} (q+1)^{2} q^4\, E_{3211254364}
  \\
  + (-1) (q-1)^{2} (q+1)^{2} q^3 (q^2+1)\, E_{3211254643} + (-1) (q-1)^{2} (q+1)^{2} q^5\, E_{3211256434} + (q-1)^{2} (q+1)^{2} q^4 (q^2+1)\, E_{3211256443}
  \\
  + (-1) (q-1)^{2} (q+1)^{2} q^5\, E_{3211264345} + (q-1)^{2} (q+1)^{2} q^4 (q^2+1)\, E_{3211264354} + (q-1)^{2} (q+1)^{2} q^4\, E_{3211264435}
  \\
  + (-1) (q-1)^{2} (q+1)^{2} q^3 (q^2 - q+1) (q^2 + q+1)\, E_{3211264543} + (-1) (q-1)^{2} (q+1)^{2} q^5\, E_{3211324564} + (q-1)^{2} (q+1)^{2} q^4\, E_{3211325464}
  \\
  + (-1) (q-1)^{2} (q+1)^{2} q^5\, E_{3211325644} + (q-1)^{2} (q+1)^{2} q^6\, E_{3211326454} + (q-1)^{2} (q+1)^{2} q^4\, E_{3211432456}
  \\
  + (-1) (q-1)^{2} (q+1)^{2} q^3\, E_{3211432546} + (q-1)^{2} q^2 (q+1)^{2} (q^2+1)^{2}\, E_{3211432564} + (-1) (q-1)^{2} (q+1)^{2} q^3 (q^2+1)\, E_{3211432645}
  \\
  + (-1) (q-1)^{2} (q+1)^{2} q^5\, E_{3211443256} + (q-1)^{2} (q+1)^{2} q^4\, E_{3211454326} + (-1) q (q-1)^{2} (q+1)^{2} (q^6 + 2 q^4 + q^2+1)\, E_{3211456432}
  \\
  + (q-1)^{2} q^2 (q+1)^{2} (q^4 + 1)\, E_{3211464325} + (-1) (q-1)^{2} (q+1)^{2} q^3 (q^2+1)\, E_{3211543264} + (q-1)^{2} q^2 (q+1)^{2} (q^2 - q+1) (q^2 + q+1)\, E_{3211546432}
  \\
  + (q-1)^{2} (q+1)^{2} q^4 (q^2+1)\, E_{3211564324} + (-1) (q-1)^{2} (q+1)^{2} q^3 (q^2 - q+1) (q^2 + q+1)\, E_{3211564432} + (q-1)^{2} (q+1)^{2} q^4\, E_{3211643245}
  \\
  + (-1) (q-1)^{2} (q+1)^{2} q^3 (q^2 - q+1) (q^2 + q+1)\, E_{3211643254} + (-1) (q-1)^{2} (q+1)^{2} q^3\, E_{3211644325} + (q-1)^{2} q^2 (q+1)^{2} (q^2+1) (q^4 + 1)\, E_{3211645432}
  \\
  + (q-1)^{2} q^2 (q+1)^{2} (q^2+1)\, E_{3212143564} + (-1) (q-1)^{2} (q+1)^{2} q^3\, E_{3212143645} + (-1) q (q-1)^{2} (q+1)^{2} (q^2 - q+1) (q^2 + q+1)\, E_{3212145643}
  \\
  + (q-1)^{2} q^2 (q+1)^{2}\, E_{3212146435} + (-1) (q-1)^{2} (q+1)^{2} q^3\, E_{3212154364} + (q-1)^{2} q^2 (q+1)^{2} (q^2+1)\, E_{3212154643}
  \\
  + (q-1)^{2} (q+1)^{2} q^4\, E_{3212156434} + (-1) (q-1)^{2} (q+1)^{2} q^3 (q^2+1)\, E_{3212156443} + (q-1)^{2} (q+1)^{2} q^4\, E_{3212164345}
  \\
  + (-1) (q-1)^{2} (q+1)^{2} q^3 (q^2+1)\, E_{3212164354} + (-1) (q-1)^{2} (q+1)^{2} q^3\, E_{3212164435} + (q-1)^{2} q^2 (q+1)^{2} (q^2 - q+1) (q^2 + q+1)\, E_{3212164543}
  \\
  + (q-1)^{2} (q+1)^{2} q^4 (q^2+1)\, E_{3213214564} + (-1) (q-1)^{2} (q+1)^{2} q^3 (q^2+1)\, E_{3213215464} + (q-1)^{2} (q+1)^{2} q^4 (q^2+1)\, E_{3213215644}
  \\
  + (-1) (q-1)^{2} (q+1)^{2} q^5 (q^2+1)\, E_{3213216454} + (-1) (q-1)^{2} (q+1)^{2} q^3 (q^2+1)\, E_{3214321456} + (q-1)^{2} q^2 (q+1)^{2} (q^2+1)\, E_{3214321546}
  \\
  + (-1) q (q-1)^{2} (q+1)^{2} (q^2 - q+1) (q^2+1) (q^2 + q+1)\, E_{3214321564} + (q-1)^{2} q^2 (q+1)^{2} (q^2 - q+1) (q^2 + q+1)\, E_{3214321645} + (q-1)^{2} (q+1)^{2} q^4 (q^2+1)\, E_{3214432156}
  \\
  + (-1) (q-1)^{2} (q+1)^{2} q^3 (q^2+1)\, E_{3214543216} + (q-1)^{2} (q+1)^{2} (q^8 + 2 q^6 + 2 q^4 + q^2+1)\, E_{3214564321} + (-1) q (q-1)^{2} (q+1)^{2} (q^6 + q^4 + 1)\, E_{3214643215}
  \\
  + (q-1)^{2} q^2 (q+1)^{2} (q^2 - q+1) (q^2 + q+1)\, E_{3215432164} + (-1) q (q-1)^{2} (q+1)^{2} (q^2+1) (q^4 + 1)\, E_{3215464321} + (-1) (q-1)^{2} (q+1)^{2} q^3 (q^2 - q+1) (q^2 + q+1)\, E_{3215643214}
  \\
  + (q-1)^{2} q^2 (q+1)^{2} (q^2+1) (q^4 + 1)\, E_{3215644321} + (-1) (q-1)^{2} (q+1)^{2} q^3\, E_{3216432145} + (q-1)^{2} q^2 (q+1)^{2} (q^2+1) (q^4 + 1)\, E_{3216432154}
  \\
  + (q-1)^{2} q^2 (q+1)^{2}\, E_{3216443215} + (-1) q (q-1)^{2} (q+1)^{2} (q^4 - q^3 + q^2 - q+1) (q^4 + q^3 + q^2 + q+1)\, E_{3216454321} + (q-1)^{2} (q+1)^{2} q^6\, E_{3321124564}
  \\
  + (-1) (q-1)^{2} (q+1)^{2} q^5\, E_{3321125464} + (q-1)^{2} (q+1)^{2} q^6\, E_{3321125644} + (-1) (q-1)^{2} (q+1)^{2} q^7\, E_{3321126454}
  \\
  + (-1) (q-1)^{2} (q+1)^{2} q^5 (q^2+1)\, E_{3321214564} + (q-1)^{2} (q+1)^{2} q^4 (q^2+1)\, E_{3321215464} + (-1) (q-1)^{2} (q+1)^{2} q^5 (q^2+1)\, E_{3321215644}
  \\
  + (q-1)^{2} (q+1)^{2} q^6 (q^2+1)\, E_{3321216454} + (-1) (q-1)^{2} (q+1)^{2} q^5\, E_{3432112456} + (q-1)^{2} (q+1)^{2} q^4\, E_{3432112546}
  \\
  + (-1) (q-1)^{2} (q+1)^{2} q^5 (q^2+1)\, E_{3432112564} + (q-1)^{2} (q+1)^{2} q^6\, E_{3432112645} + (q-1)^{2} (q+1)^{2} q^4 (q^2+1)\, E_{3432121456}
  \\
  + (-1) (q-1)^{2} (q+1)^{2} q^3 (q^2+1)\, E_{3432121546} + (q-1)^{2} (q+1)^{2} q^4 (q^2+1)^{2}\, E_{3432121564} + (-1) (q-1)^{2} (q+1)^{2} q^5 (q^2+1)\, E_{3432121645}
  \\
  + (q-1)^{2} (q+1)^{2} q^6\, E_{3443211256} + (-1) (q-1)^{2} (q+1)^{2} q^5 (q^2+1)\, E_{3443212156} + (-1) (q-1)^{2} (q+1)^{2} q^5\, E_{3454321126}
  \\
  + (q-1)^{2} (q+1)^{2} q^4 (q^2+1)\, E_{3454321216} + (q-1)^{2} (q+1)^{2} q^6\, E_{3456432112} + (-1) (q-1)^{2} (q+1)^{2} q^5 (q^2+1)\, E_{3456432121}
  \\
  + (-1) (q-1)^{2} (q+1)^{2} q^7\, E_{3464321125} + (q-1)^{2} (q+1)^{2} q^6 (q^2+1)\, E_{3464321215} + (q-1)^{2} (q+1)^{2} q^6\, E_{3543211264}
  \\
  + (-1) (q-1)^{2} (q+1)^{2} q^5 (q^2+1)\, E_{3543212164} + (-1) (q-1)^{2} (q+1)^{2} q^7\, E_{3564321124} + (q-1)^{2} (q+1)^{2} q^6 (q^2+1)\, E_{3564321214}
  \\
  + (q-1)^{2} (q+1)^{2} q^8\, E_{3643211254} + (-1) (q-1)^{2} (q+1)^{2} q^7 (q^2+1)\, E_{3643212154} + (q-1)^{2} (q+1)^{2} q^4 (q^2+1)\, E_{4321123564}
  \\
  + (-1) (q-1)^{2} (q+1)^{2} q^5\, E_{4321123645} + (-1) (q-1)^{2} (q+1)^{2} q^5\, E_{4321124356} + (q-1)^{2} (q+1)^{2} q^4\, E_{4321125436}
  \\
  + (-1) (q-1)^{2} (q+1)^{2} q^3 (q^2+1)^{2}\, E_{4321125643} + (q-1)^{2} (q+1)^{2} q^4 (q^2+1)\, E_{4321126435} + (-1) (q-1)^{2} (q+1)^{2} q^3 (q^2+1)\, E_{4321132564}
  \\
  + (q-1)^{2} (q+1)^{2} q^4\, E_{4321132645} + (q-1)^{2} (q+1)^{2} q^4 (q^2+1)\, E_{4321143256} + (-1) (q-1)^{2} (q+1)^{2} q^3 (q^2+1)\, E_{4321154326}
  \\
  + (q-1)^{2} q^2 (q+1)^{2} (q^2 - q+1) (q^2+1) (q^2 + q+1)\, E_{4321156432} + (-1) (q-1)^{2} (q+1)^{2} q^3 (q^2 - q+1) (q^2 + q+1)\, E_{4321164325} + (-1) (q-1)^{2} (q+1)^{2} q^3 (q^2+1)\, E_{4321213564}
  \\
  + (q-1)^{2} (q+1)^{2} q^4\, E_{4321213645} + (q-1)^{2} (q+1)^{2} q^4\, E_{4321214356} + (-1) (q-1)^{2} (q+1)^{2} q^3\, E_{4321215436}
  \\
  + (q-1)^{2} q^2 (q+1)^{2} (q^2+1)^{2}\, E_{4321215643} + (-1) (q-1)^{2} (q+1)^{2} q^3 (q^2+1)\, E_{4321216435} + (q-1)^{2} q^2 (q+1)^{2} (q^2+1)\, E_{4321321564}
  \\
  + (-1) (q-1)^{2} (q+1)^{2} q^3\, E_{4321321645} + (-1) (q-1)^{2} (q+1)^{2} q^3 (q^2 - q+1) (q^2 + q+1)\, E_{4321432156} + (q-1)^{2} q^2 (q+1)^{2} (q^2 - q+1) (q^2 + q+1)\, E_{4321543216}
  \\
  + (-1) q (q-1)^{2} (q+1)^{2} (q^2+1)^{2} (q^4 + 1)\, E_{4321564321} + (q-1)^{2} q^2 (q+1)^{2} (q^2+1) (q^4 + 1)\, E_{4321643215} + (q-1)^{2} (q+1)^{2} q^6\, E_{4432112356}
  \\
  + (-1) (q-1)^{2} (q+1)^{2} q^5 (q^2+1)\, E_{4432113256} + (-1) (q-1)^{2} (q+1)^{2} q^5\, E_{4432121356} + (q-1)^{2} (q+1)^{2} q^4 (q^2 - q+1) (q^2 + q+1)\, E_{4432132156}
  \\
  + (-1) (q-1)^{2} (q+1)^{2} q^5\, E_{4543211236} + (q-1)^{2} (q+1)^{2} q^4 (q^2+1)\, E_{4543211326} + (q-1)^{2} (q+1)^{2} q^4\, E_{4543212136}
  \\
  + (-1) (q-1)^{2} (q+1)^{2} q^3 (q^2 - q+1) (q^2 + q+1)\, E_{4543213216} + (q-1)^{2} (q+1)^{2} q^6\, E_{4564321123} + (-1) (q-1)^{2} (q+1)^{2} q^5 (q^2+1)\, E_{4564321132}
  \\
  + (-1) (q-1)^{2} (q+1)^{2} q^5\, E_{4564321213} + (q-1)^{2} (q+1)^{2} q^4 (q^2 - q+1) (q^2 + q+1)\, E_{4564321321} + (-1) (q-1)^{2} (q+1)^{2} q^7\, E_{4643211235}
  \\
  + (q-1)^{2} (q+1)^{2} q^6 (q^2+1)\, E_{4643211325} + (q-1)^{2} (q+1)^{2} q^6\, E_{4643212135} + (-1) (q-1)^{2} (q+1)^{2} q^5 (q^2 - q+1) (q^2 + q+1)\, E_{4643213215}
  \\
  + (-1) (q-1)^{2} (q+1)^{2} q^5\, E_{5432112364} + (q-1)^{2} (q+1)^{2} q^4 (q^2+1)\, E_{5432112643} + (q-1)^{2} (q+1)^{2} q^4\, E_{5432113264}
  \\
  + (-1) (q-1)^{2} (q+1)^{2} q^3 (q^2 - q+1) (q^2 + q+1)\, E_{5432116432} + (q-1)^{2} (q+1)^{2} q^4\, E_{5432121364} + (-1) (q-1)^{2} (q+1)^{2} q^3 (q^2+1)\, E_{5432121643}
  \\
  + (-1) (q-1)^{2} (q+1)^{2} q^3\, E_{5432132164} + (q-1)^{2} q^2 (q+1)^{2} (q^2+1) (q^4 + 1)\, E_{5432164321} + (q-1)^{2} (q+1)^{2} q^6\, E_{5643211234}
  \\
  + (-1) (q-1)^{2} (q+1)^{2} q^5 (q^2+1)\, E_{5643211243} + (-1) (q-1)^{2} (q+1)^{2} q^5\, E_{5643211324} + (q-1)^{2} (q+1)^{2} q^4 (q^2 - q+1) (q^2 + q+1)\, E_{5643211432}
  \\
  + (-1) (q-1)^{2} (q+1)^{2} q^5\, E_{5643212134} + (q-1)^{2} (q+1)^{2} q^4 (q^2+1)\, E_{5643212143} + (q-1)^{2} (q+1)^{2} q^4\, E_{5643213214}
  \\
  + (-1) (q-1)^{2} (q+1)^{2} q^3 (q^2+1) (q^4 + 1)\, E_{5643214321} + (q-1)^{2} (q+1)^{2} q^6\, E_{6432112345} + (-1) (q-1)^{2} (q+1)^{2} q^5 (q^2+1)\, E_{6432112354}
  \\
  + (-1) (q-1)^{2} (q+1)^{2} q^5\, E_{6432112435} + (q-1)^{2} (q+1)^{2} q^4 (q^2 - q+1) (q^2 + q+1)\, E_{6432112543} + (-1) (q-1)^{2} (q+1)^{2} q^5\, E_{6432113245}
  \\
  + (q-1)^{2} (q+1)^{2} q^4 (q^2+1)\, E_{6432113254} + (q-1)^{2} (q+1)^{2} q^4\, E_{6432114325} + (-1) (q-1)^{2} (q+1)^{2} q^3 (q^2+1) (q^4 + 1)\, E_{6432115432}
  \\
  + (-1) (q-1)^{2} (q+1)^{2} q^5\, E_{6432121345} + (q-1)^{2} (q+1)^{2} q^4 (q^2+1)\, E_{6432121354} + (q-1)^{2} (q+1)^{2} q^4\, E_{6432121435}
  \\
  + (-1) (q-1)^{2} (q+1)^{2} q^3 (q^2 - q+1) (q^2 + q+1)\, E_{6432121543} + (q-1)^{2} (q+1)^{2} q^4\, E_{6432132145} + (-1) (q-1)^{2} (q+1)^{2} q^3 (q^2+1)\, E_{6432132154}
  \\
  + (-1) (q-1)^{2} (q+1)^{2} q^3\, E_{6432143215} + (q-1)^{2} q^2 (q+1)^{2} (q^4 - q^3 + q^2 - q+1) (q^4 + q^3 + q^2 + q+1)\, E_{6432154321}
\end{multline*}

%% Pair (-L_2, -L_1), weight space v8, 1 terms
$e^*_{(-L_2, -L_1)} = ((-1) (q-1) (q+1))/(q)\, E_{1}$

%% Pair (-L_3, -L_1), weight space v14, 2 terms
$e^*_{(-L_3, -L_1)} = (q-1) (q+1)\, E_{12} + ((-1) (q-1) (q+1))/(q)\, E_{21}$

%% Pair (-L_4, -L_1), weight space v20, 4 terms
\begin{multline*}
e^*_{(-L_4, -L_1)} = 
  (-1) (q-1) q (q+1)\, E_{123} + (q-1) (q+1)\, E_{132} + (q-1) (q+1)\, E_{213}
  \\
  + ((-1) (q-1) (q+1))/(q)\, E_{321}
\end{multline*}

%% Pair (-L_5, -L_1), weight space v25, 8 terms
\begin{multline*}
e^*_{(-L_5, -L_1)} = 
  (q-1) (q+1) q^2\, E_{1234} + (-1) (q-1) q (q+1)\, E_{1243} + (-1) (q-1) q (q+1)\, E_{1324}
  \\
  + (q-1) (q+1)\, E_{1432} + (-1) (q-1) q (q+1)\, E_{2134} + (q-1) (q+1)\, E_{2143}
  \\
  + (q-1) (q+1)\, E_{3214} + ((-1) (q-1) (q+1))/(q)\, E_{4321}
\end{multline*}

%% Pair (-L_6, -L_1), weight space v30, 16 terms
\begin{multline*}
e^*_{(-L_6, -L_1)} = 
  (-1) (q-1) (q+1) q^3\, E_{12345} + (q-1) (q+1) q^2\, E_{12354} + (q-1) (q+1) q^2\, E_{12435}
  \\
  + (-1) (q-1) q (q+1)\, E_{12543} + (q-1) (q+1) q^2\, E_{13245} + (-1) (q-1) q (q+1)\, E_{13254}
  \\
  + (-1) (q-1) q (q+1)\, E_{14325} + (q-1) (q+1)\, E_{15432} + (q-1) (q+1) q^2\, E_{21345}
  \\
  + (-1) (q-1) q (q+1)\, E_{21354} + (-1) (q-1) q (q+1)\, E_{21435} + (q-1) (q+1)\, E_{21543}
  \\
  + (-1) (q-1) q (q+1)\, E_{32145} + (q-1) (q+1)\, E_{32154} + (q-1) (q+1)\, E_{43215}
  \\
  + ((-1) (q-1) (q+1))/(q)\, E_{54321}
\end{multline*}

%% Pair (L_6, -L_1), weight space v31, 16 terms
\begin{multline*}
e^*_{(L_6, -L_1)} = 
  (-1) (q-1) (q+1) q^3\, E_{12346} + (q-1) (q+1) q^2\, E_{12364} + (q-1) (q+1) q^2\, E_{12436}
  \\
  + (-1) (q-1) q (q+1)\, E_{12643} + (q-1) (q+1) q^2\, E_{13246} + (-1) (q-1) q (q+1)\, E_{13264}
  \\
  + (-1) (q-1) q (q+1)\, E_{14326} + (q-1) (q+1)\, E_{16432} + (q-1) (q+1) q^2\, E_{21346}
  \\
  + (-1) (q-1) q (q+1)\, E_{21364} + (-1) (q-1) q (q+1)\, E_{21436} + (q-1) (q+1)\, E_{21643}
  \\
  + (-1) (q-1) q (q+1)\, E_{32146} + (q-1) (q+1)\, E_{32164} + (q-1) (q+1)\, E_{43216}
  \\
  + ((-1) (q-1) (q+1))/(q)\, E_{64321}
\end{multline*}

%% Pair (L_5, -L_1), weight space v34, 32 terms
\begin{multline*}
e^*_{(L_5, -L_1)} = 
  (q-1) (q+1) q^4\, E_{123456} + (-1) (q-1) (q+1) q^3\, E_{123546} + (q-1) (q+1) q^2\, E_{123564}
  \\
  + (-1) (q-1) (q+1) q^3\, E_{123645} + (-1) (q-1) (q+1) q^3\, E_{124356} + (q-1) (q+1) q^2\, E_{125436}
  \\
  + (-1) (q-1) q (q+1)\, E_{125643} + (q-1) (q+1) q^2\, E_{126435} + (-1) (q-1) (q+1) q^3\, E_{132456}
  \\
  + (q-1) (q+1) q^2\, E_{132546} + (-1) (q-1) q (q+1)\, E_{132564} + (q-1) (q+1) q^2\, E_{132645}
  \\
  + (q-1) (q+1) q^2\, E_{143256} + (-1) (q-1) q (q+1)\, E_{154326} + (q-1) (q+1)\, E_{156432}
  \\
  + (-1) (q-1) q (q+1)\, E_{164325} + (-1) (q-1) (q+1) q^3\, E_{213456} + (q-1) (q+1) q^2\, E_{213546}
  \\
  + (-1) (q-1) q (q+1)\, E_{213564} + (q-1) (q+1) q^2\, E_{213645} + (q-1) (q+1) q^2\, E_{214356}
  \\
  + (-1) (q-1) q (q+1)\, E_{215436} + (q-1) (q+1)\, E_{215643} + (-1) (q-1) q (q+1)\, E_{216435}
  \\
  + (q-1) (q+1) q^2\, E_{321456} + (-1) (q-1) q (q+1)\, E_{321546} + (q-1) (q+1)\, E_{321564}
  \\
  + (-1) (q-1) q (q+1)\, E_{321645} + (-1) (q-1) q (q+1)\, E_{432156} + (q-1) (q+1)\, E_{543216}
  \\
  + ((-1) (q-1) (q+1))/(q)\, E_{564321} + (q-1) (q+1)\, E_{643215}
\end{multline*}

%% Pair (L_4, -L_1), weight space v35, 80 terms
\begin{multline*}
e^*_{(L_4, -L_1)} = 
  (-1) (q-1) (q+1) q^5\, E_{1234564} + (q-1) (q+1) q^4\, E_{1235464} + (-1) (q-1) (q+1) q^5\, E_{1235644}
  \\
  + (q-1) (q+1) q^6\, E_{1236454} + (q-1) (q+1) q^4\, E_{1243456} + (-1) (q-1) (q+1) q^3\, E_{1243546}
  \\
  + (q-1) (q+1) q^2 (q^2+1)\, E_{1243564} + (-1) (q-1) (q+1) q^3\, E_{1243645} + (-1) (q-1) (q+1) q^5\, E_{1244356}
  \\
  + (q-1) (q+1) q^4\, E_{1245436} + (-1) (q-1) q (q+1)\, E_{1245643} + (q-1) (q+1) q^2 (q^4 - q^2+1)\, E_{1246435}
  \\
  + (-1) (q-1) (q+1) q^3\, E_{1254364} + (-1) (q-1)^{2} q^2 (q+1)^{2}\, E_{1254643} + (q-1) (q+1) q^4\, E_{1256434}
  \\
  + (q-1)^{2} (q+1)^{2} q^3\, E_{1256443} + (-1) (q-1)^{2} (q+1)^{2} q^4\, E_{1264345} + (-1) (q-1) (q+1) q^3\, E_{1264354}
  \\
  + (q-1)^{2} (q+1)^{2} q^3\, E_{1264435} + (-1) (q-1)^{2} q^2 (q+1)^{2} (q^2+1)\, E_{1264543} + (q-1) (q+1) q^4\, E_{1324564}
  \\
  + (-1) (q-1) (q+1) q^3\, E_{1325464} + (q-1) (q+1) q^4\, E_{1325644} + (-1) (q-1) (q+1) q^5\, E_{1326454}
  \\
  + (-1) (q-1) (q+1) q^3\, E_{1432456} + (q-1) (q+1) q^2\, E_{1432546} + (-1) (q-1) q (q+1) (q^2+1)\, E_{1432564}
  \\
  + (q-1) (q+1) q^2\, E_{1432645} + (q-1) (q+1) q^4\, E_{1443256} + (-1) (q-1) (q+1) q^3\, E_{1454326}
  \\
  + (q-1) (q+1)\, E_{1456432} + (-1) (q-1) q (q+1) (q^4 - q^2+1)\, E_{1464325} + (q-1) (q+1) q^2\, E_{1543264}
  \\
  + q (q-1)^{2} (q+1)^{2}\, E_{1546432} + (-1) (q-1) (q+1) q^3\, E_{1564324} + (-1) (q-1)^{2} q^2 (q+1)^{2}\, E_{1564432}
  \\
  + (q-1)^{2} (q+1)^{2} q^3\, E_{1643245} + (q-1) (q+1) q^2\, E_{1643254} + (-1) (q-1)^{2} q^2 (q+1)^{2}\, E_{1644325}
  \\
  + q (q-1)^{2} (q+1)^{2} (q^2+1)\, E_{1645432} + (q-1) (q+1) q^4\, E_{2134564} + (-1) (q-1) (q+1) q^3\, E_{2135464}
  \\
  + (q-1) (q+1) q^4\, E_{2135644} + (-1) (q-1) (q+1) q^5\, E_{2136454} + (-1) (q-1) (q+1) q^3\, E_{2143456}
  \\
  + (q-1) (q+1) q^2\, E_{2143546} + (-1) (q-1) q (q+1) (q^2+1)\, E_{2143564} + (q-1) (q+1) q^2\, E_{2143645}
  \\
  + (q-1) (q+1) q^4\, E_{2144356} + (-1) (q-1) (q+1) q^3\, E_{2145436} + (q-1) (q+1)\, E_{2145643}
  \\
  + (-1) (q-1) q (q+1) (q^4 - q^2+1)\, E_{2146435} + (q-1) (q+1) q^2\, E_{2154364} + q (q-1)^{2} (q+1)^{2}\, E_{2154643}
  \\
  + (-1) (q-1) (q+1) q^3\, E_{2156434} + (-1) (q-1)^{2} q^2 (q+1)^{2}\, E_{2156443} + (q-1)^{2} (q+1)^{2} q^3\, E_{2164345}
  \\
  + (q-1) (q+1) q^2\, E_{2164354} + (-1) (q-1)^{2} q^2 (q+1)^{2}\, E_{2164435} + q (q-1)^{2} (q+1)^{2} (q^2+1)\, E_{2164543}
  \\
  + (-1) (q-1) (q+1) q^3\, E_{3214564} + (q-1) (q+1) q^2\, E_{3215464} + (-1) (q-1) (q+1) q^3\, E_{3215644}
  \\
  + (q-1) (q+1) q^4\, E_{3216454} + (q-1) (q+1) q^2\, E_{4321456} + (-1) (q-1) q (q+1)\, E_{4321546}
  \\
  + (q-1) (q+1) (q^2+1)\, E_{4321564} + (-1) (q-1) q (q+1)\, E_{4321645} + (-1) (q-1) (q+1) q^3\, E_{4432156}
  \\
  + (q-1) (q+1) q^2\, E_{4543216} + ((-1) (q-1) (q+1))/(q)\, E_{4564321} + (q-1) (q+1) (q^4 - q^2+1)\, E_{4643215}
  \\
  + (-1) (q-1) q (q+1)\, E_{5432164} + (-1) (q-1)^{2} (q+1)^{2}\, E_{5464321} + (q-1) (q+1) q^2\, E_{5643214}
  \\
  + q (q-1)^{2} (q+1)^{2}\, E_{5644321} + (-1) (q-1)^{2} q^2 (q+1)^{2}\, E_{6432145} + (-1) (q-1) q (q+1)\, E_{6432154}
  \\
  + q (q-1)^{2} (q+1)^{2}\, E_{6443215} + (-1) (q-1)^{2} (q+1)^{2} (q^2+1)\, E_{6454321}
\end{multline*}

%% Pair (L_3, -L_1), weight space v33, 182 terms
\begin{multline*}
e^*_{(L_3, -L_1)} = 
  (q-1) (q+1) q^6\, E_{12345643} + (-1) (q-1) (q+1) q^5\, E_{12354643} + (q-1) (q+1) q^6\, E_{12356443}
  \\
  + (-1) (q-1) (q+1) q^7\, E_{12364543} + (-1) (q-1) (q+1) q^5\, E_{12434356} + (q-1) (q+1) q^4\, E_{12435436}
  \\
  + (-1) (q-1) (q+1) q^5 (q^2+1)\, E_{12435643} + (q-1) (q+1) q^6\, E_{12436435} + (q-1) (q+1) q^6\, E_{12443356}
  \\
  + (-1) (q-1) (q+1) q^5\, E_{12454336} + (q-1) (q+1) q^6\, E_{12456433} + (-1) (q-1) (q+1) q^7\, E_{12464335}
  \\
  + (q-1) (q+1) q^6\, E_{12543643} + (-1) (q-1) (q+1) q^7\, E_{12564343} + (q-1) (q+1) q^8\, E_{12643543}
  \\
  + (-1) (q-1) (q+1) q^5\, E_{13234564} + (q-1) (q+1) q^4\, E_{13235464} + (-1) (q-1) (q+1) q^5\, E_{13235644}
  \\
  + (q-1) (q+1) q^6\, E_{13236454} + (q-1) (q+1) q^4\, E_{13243456} + (-1) (q-1) (q+1) q^3\, E_{13243546}
  \\
  + (q-1) (q+1) q^2 (q^2+1)\, E_{13243564} + (-1) (q-1) (q+1) q^3\, E_{13243645} + (-1) (q-1) (q+1) q^5\, E_{13244356}
  \\
  + (q-1) (q+1) q^4\, E_{13245436} + (-1) (q-1) q (q+1) (q^4 + 1)\, E_{13245643} + (q-1) (q+1) q^2 (q^4 - q^2+1)\, E_{13246435}
  \\
  + (-1) (q-1) (q+1) q^3\, E_{13254364} + (q-1) (q+1) q^2\, E_{13254643} + (q-1) (q+1) q^4\, E_{13256434}
  \\
  + (-1) (q-1) (q+1) q^3\, E_{13256443} + (-1) (q-1)^{2} (q+1)^{2} q^4\, E_{13264345} + (-1) (q-1) (q+1) q^3\, E_{13264354}
  \\
  + (q-1)^{2} (q+1)^{2} q^3\, E_{13264435} + (q-1) (q+1) q^2\, E_{13264543} + (q-1) (q+1) q^6\, E_{13324564}
  \\
  + (-1) (q-1) (q+1) q^5\, E_{13325464} + (q-1) (q+1) q^6\, E_{13325644} + (-1) (q-1) (q+1) q^7\, E_{13326454}
  \\
  + (-1) (q-1) (q+1) q^5\, E_{13432456} + (q-1) (q+1) q^4\, E_{13432546} + (-1) (q-1) q (q+1) (q^2+1) (q^4 - q^2+1)\, E_{13432564}
  \\
  + (q-1) (q+1) q^2 (q^4 - q^2+1)\, E_{13432645} + (q-1) (q+1) q^6\, E_{13443256} + (-1) (q-1) (q+1) q^5\, E_{13454326}
  \\
  + (q-1) (q+1)\, E_{13456432} + (-1) (q-1) q (q+1) (q^6 - q^2+1)\, E_{13464325} + (q-1) (q+1) q^2 (q^4 - q^2+1)\, E_{13543264}
  \\
  + q (q-1)^{2} (q+1)^{2} (q^2+1)\, E_{13546432} + (-1) (q-1) (q+1) q^3 (q^4 - q^2+1)\, E_{13564324} + (-1) (q-1)^{2} q^2 (q+1)^{2} (q^2+1)\, E_{13564432}
  \\
  + (q-1)^{2} (q+1)^{2} q^3\, E_{13643245} + (q-1) (q+1) q^2 (q^6 - q^4 + 1)\, E_{13643254} + (-1) (q-1)^{2} q^2 (q+1)^{2}\, E_{13644325}
  \\
  + q (q-1)^{2} (q+1)^{2} (q^2 - q+1) (q^2 + q+1)\, E_{13645432} + (q-1)^{2} (q+1)^{2} q^3 (q^2+1)\, E_{14323564} + (-1) (q-1)^{2} (q+1)^{2} q^4\, E_{14323645}
  \\
  + (q-1) (q+1) q^4\, E_{14324356} + (-1) (q-1) (q+1) q^3\, E_{14325436} + (q-1) (q+1) q^2 (q^2+1)\, E_{14325643}
  \\
  + (-1) (q-1) (q+1) q^3\, E_{14326435} + (-1) (q-1)^{2} q^2 (q+1)^{2} (q^2+1)\, E_{14332564} + (q-1)^{2} (q+1)^{2} q^3\, E_{14332645}
  \\
  + (q-1)^{2} (q+1)^{2} q^3\, E_{14343256} + (-1) (q-1)^{2} q^2 (q+1)^{2}\, E_{14354326} + q (q-1)^{2} (q+1)^{2} (q^2+1)^{2}\, E_{14356432}
  \\
  + (-1) (q-1)^{2} q^2 (q+1)^{2} (q^2+1)\, E_{14364325} + (-1) (q-1) (q+1) q^5\, E_{14432356} + (-1) (q-1)^{2} (q+1)^{2} q^4\, E_{14433256}
  \\
  + (q-1) (q+1) q^4\, E_{14543236} + (q-1)^{2} (q+1)^{2} q^3\, E_{14543326} + (-1) (q-1) (q+1) q^5\, E_{14564323}
  \\
  + (-1) (q-1)^{2} (q+1)^{2} q^4\, E_{14564332} + (q-1) (q+1) q^6\, E_{14643235} + (q-1)^{2} (q+1)^{2} q^5\, E_{14643325}
  \\
  + (-1) (q-1)^{2} (q+1)^{2} q^4\, E_{15432364} + (-1) (q-1) (q+1) q^3\, E_{15432643} + (q-1)^{2} (q+1)^{2} q^3\, E_{15433264}
  \\
  + (-1) (q-1)^{2} q^2 (q+1)^{2} (q^2+1)\, E_{15436432} + (q-1)^{2} (q+1)^{2} q^5\, E_{15643234} + (q-1) (q+1) q^4\, E_{15643243}
  \\
  + (-1) (q-1)^{2} (q+1)^{2} q^4\, E_{15643324} + (q-1)^{2} (q+1)^{2} q^3 (q^2+1)\, E_{15643432} + (q-1)^{2} (q+1)^{2} q^5\, E_{16432345}
  \\
  + (-1) (q-1)^{2} (q+1)^{2} q^4 (q^2+1)\, E_{16432354} + (-1) (q-1)^{2} (q+1)^{2} q^4\, E_{16432435} + (-1) (q-1) (q+1) q^3\, E_{16432543}
  \\
  + (-1) (q-1)^{2} (q+1)^{2} q^4\, E_{16433245} + (q-1)^{2} (q+1)^{2} q^3 (q^2+1)\, E_{16433254} + (q-1)^{2} (q+1)^{2} q^3\, E_{16434325}
  \\
  + (-1) (q-1)^{2} q^2 (q+1)^{2} (q^2 - q+1) (q^2 + q+1)\, E_{16435432} + (-1) (q-1) (q+1) q^5\, E_{21345643} + (q-1) (q+1) q^4\, E_{21354643}
  \\
  + (-1) (q-1) (q+1) q^5\, E_{21356443} + (q-1) (q+1) q^6\, E_{21364543} + (q-1) (q+1) q^4\, E_{21434356}
  \\
  + (-1) (q-1) (q+1) q^3\, E_{21435436} + (q-1) (q+1) q^4 (q^2+1)\, E_{21435643} + (-1) (q-1) (q+1) q^5\, E_{21436435}
  \\
  + (-1) (q-1) (q+1) q^5\, E_{21443356} + (q-1) (q+1) q^4\, E_{21454336} + (-1) (q-1) (q+1) q^5\, E_{21456433}
  \\
  + (q-1) (q+1) q^6\, E_{21464335} + (-1) (q-1) (q+1) q^5\, E_{21543643} + (q-1) (q+1) q^6\, E_{21564343}
  \\
  + (-1) (q-1) (q+1) q^7\, E_{21643543} + (q-1) (q+1) q^4\, E_{32134564} + (-1) (q-1) (q+1) q^3\, E_{32135464}
  \\
  + (q-1) (q+1) q^4\, E_{32135644} + (-1) (q-1) (q+1) q^5\, E_{32136454} + (-1) (q-1) (q+1) q^3\, E_{32143456}
  \\
  + (q-1) (q+1) q^2\, E_{32143546} + (-1) (q-1) q (q+1) (q^2+1)\, E_{32143564} + (q-1) (q+1) q^2\, E_{32143645}
  \\
  + (q-1) (q+1) q^4\, E_{32144356} + (-1) (q-1) (q+1) q^3\, E_{32145436} + (q-1) (q+1) (q^4 + 1)\, E_{32145643}
  \\
  + (-1) (q-1) q (q+1) (q^4 - q^2+1)\, E_{32146435} + (q-1) (q+1) q^2\, E_{32154364} + (-1) (q-1) q (q+1)\, E_{32154643}
  \\
  + (-1) (q-1) (q+1) q^3\, E_{32156434} + (q-1) (q+1) q^2\, E_{32156443} + (q-1)^{2} (q+1)^{2} q^3\, E_{32164345}
  \\
  + (q-1) (q+1) q^2\, E_{32164354} + (-1) (q-1)^{2} q^2 (q+1)^{2}\, E_{32164435} + (-1) (q-1) q (q+1)\, E_{32164543}
  \\
  + (-1) (q-1) (q+1) q^5\, E_{33214564} + (q-1) (q+1) q^4\, E_{33215464} + (-1) (q-1) (q+1) q^5\, E_{33215644}
  \\
  + (q-1) (q+1) q^6\, E_{33216454} + (q-1) (q+1) q^4\, E_{34321456} + (-1) (q-1) (q+1) q^3\, E_{34321546}
  \\
  + (q-1) (q+1) (q^2+1) (q^4 - q^2+1)\, E_{34321564} + (-1) (q-1) q (q+1) (q^4 - q^2+1)\, E_{34321645} + (-1) (q-1) (q+1) q^5\, E_{34432156}
  \\
  + (q-1) (q+1) q^4\, E_{34543216} + ((-1) (q-1) (q+1))/(q)\, E_{34564321} + (q-1) (q+1) (q^6 - q^2+1)\, E_{34643215}
  \\
  + (-1) (q-1) q (q+1) (q^4 - q^2+1)\, E_{35432164} + (-1) (q-1)^{2} (q+1)^{2} (q^2+1)\, E_{35464321} + (q-1) (q+1) q^2 (q^4 - q^2+1)\, E_{35643214}
  \\
  + q (q-1)^{2} (q+1)^{2} (q^2+1)\, E_{35644321} + (-1) (q-1)^{2} q^2 (q+1)^{2}\, E_{36432145} + (-1) (q-1) q (q+1) (q^6 - q^4 + 1)\, E_{36432154}
  \\
  + q (q-1)^{2} (q+1)^{2}\, E_{36443215} + (-1) (q-1)^{2} (q+1)^{2} (q^2 - q+1) (q^2 + q+1)\, E_{36454321} + (-1) (q-1)^{2} q^2 (q+1)^{2} (q^2+1)\, E_{43213564}
  \\
  + (q-1)^{2} (q+1)^{2} q^3\, E_{43213645} + (-1) (q-1) (q+1) q^3\, E_{43214356} + (q-1) (q+1) q^2\, E_{43215436}
  \\
  + (-1) (q-1) q (q+1) (q^2+1)\, E_{43215643} + (q-1) (q+1) q^2\, E_{43216435} + q (q-1)^{2} (q+1)^{2} (q^2+1)\, E_{43321564}
  \\
  + (-1) (q-1)^{2} q^2 (q+1)^{2}\, E_{43321645} + (-1) (q-1)^{2} q^2 (q+1)^{2}\, E_{43432156} + q (q-1)^{2} (q+1)^{2}\, E_{43543216}
  \\
  + (-1) (q-1)^{2} (q+1)^{2} (q^2+1)^{2}\, E_{43564321} + q (q-1)^{2} (q+1)^{2} (q^2+1)\, E_{43643215} + (q-1) (q+1) q^4\, E_{44321356}
  \\
  + (q-1)^{2} (q+1)^{2} q^3\, E_{44332156} + (-1) (q-1) (q+1) q^3\, E_{45432136} + (-1) (q-1)^{2} q^2 (q+1)^{2}\, E_{45433216}
  \\
  + (q-1) (q+1) q^4\, E_{45643213} + (q-1)^{2} (q+1)^{2} q^3\, E_{45643321} + (-1) (q-1) (q+1) q^5\, E_{46432135}
  \\
  + (-1) (q-1)^{2} (q+1)^{2} q^4\, E_{46433215} + (q-1)^{2} (q+1)^{2} q^3\, E_{54321364} + (q-1) (q+1) q^2\, E_{54321643}
  \\
  + (-1) (q-1)^{2} q^2 (q+1)^{2}\, E_{54332164} + q (q-1)^{2} (q+1)^{2} (q^2+1)\, E_{54364321} + (-1) (q-1)^{2} (q+1)^{2} q^4\, E_{56432134}
  \\
  + (-1) (q-1) (q+1) q^3\, E_{56432143} + (q-1)^{2} (q+1)^{2} q^3\, E_{56433214} + (-1) (q-1)^{2} q^2 (q+1)^{2} (q^2+1)\, E_{56434321}
  \\
  + (-1) (q-1)^{2} (q+1)^{2} q^4\, E_{64321345} + (q-1)^{2} (q+1)^{2} q^3 (q^2+1)\, E_{64321354} + (q-1)^{2} (q+1)^{2} q^3\, E_{64321435}
  \\
  + (q-1) (q+1) q^2\, E_{64321543} + (q-1)^{2} (q+1)^{2} q^3\, E_{64332145} + (-1) (q-1)^{2} q^2 (q+1)^{2} (q^2+1)\, E_{64332154}
  \\
  + (-1) (q-1)^{2} q^2 (q+1)^{2}\, E_{64343215} + q (q-1)^{2} (q+1)^{2} (q^2 - q+1) (q^2 + q+1)\, E_{64354321}
\end{multline*}

%% Pair (L_2, -L_1), weight space v1, 407 terms
\begin{multline*}
e^*_{(L_2, -L_1)} = 
  (-1) (q-1) (q+1) q^7\, E_{123456432} + (q-1) (q+1) q^6\, E_{123546432} + (-1) (q-1) (q+1) q^7\, E_{123564432}
  \\
  + (q-1) (q+1) q^8\, E_{123645432} + (q-1) (q+1) q^6\, E_{124343256} + (-1) (q-1) (q+1) q^5\, E_{124354326}
  \\
  + (q-1) (q+1) q^6 (q^2+1)\, E_{124356432} + (-1) (q-1) (q+1) q^7\, E_{124364325} + (-1) (q-1) (q+1) q^7\, E_{124433256}
  \\
  + (q-1) (q+1) q^6\, E_{124543326} + (-1) (q-1) (q+1) q^7\, E_{124564332} + (q-1) (q+1) q^8\, E_{124643325}
  \\
  + (-1) (q-1) (q+1) q^7\, E_{125436432} + (q-1) (q+1) q^8\, E_{125643432} + (-1) (q-1) (q+1) q^9\, E_{126435432}
  \\
  + (q-1) (q+1) q^6\, E_{132324564} + (-1) (q-1) (q+1) q^5\, E_{132325464} + (q-1) (q+1) q^6\, E_{132325644}
  \\
  + (-1) (q-1) (q+1) q^7\, E_{132326454} + (-1) (q-1) (q+1) q^5\, E_{132432456} + (q-1) (q+1) q^4\, E_{132432546}
  \\
  + (-1) (q-1) (q+1) q^5 (q^2+1)\, E_{132432564} + (q-1) (q+1) q^6\, E_{132432645} + (q-1) (q+1) q^6\, E_{132443256}
  \\
  + (-1) (q-1) (q+1) q^5\, E_{132454326} + (q-1) (q+1) q^6 (q^2+1)\, E_{132456432} + (-1) (q-1) (q+1) q^7\, E_{132464325}
  \\
  + (q-1) (q+1) q^6\, E_{132543264} + (-1) (q-1) (q+1) q^7\, E_{132546432} + (-1) (q-1) (q+1) q^7\, E_{132564324}
  \\
  + (q-1) (q+1) q^8\, E_{132564432} + (q-1) (q+1) q^8\, E_{132643254} + (-1) (q-1) (q+1) q^9\, E_{132645432}
  \\
  + (-1) (q-1) (q+1) q^7\, E_{133224564} + (q-1) (q+1) q^6\, E_{133225464} + (-1) (q-1) (q+1) q^7\, E_{133225644}
  \\
  + (q-1) (q+1) q^8\, E_{133226454} + (q-1) (q+1) q^6\, E_{134322456} + (-1) (q-1) (q+1) q^5\, E_{134322546}
  \\
  + (q-1) (q+1) q^6 (q^2+1)\, E_{134322564} + (-1) (q-1) (q+1) q^7\, E_{134322645} + (-1) (q-1) (q+1) q^7\, E_{134432256}
  \\
  + (q-1) (q+1) q^6\, E_{134543226} + (-1) (q-1) (q+1) q^7\, E_{134564322} + (q-1) (q+1) q^8\, E_{134643225}
  \\
  + (-1) (q-1) (q+1) q^7\, E_{135432264} + (q-1) (q+1) q^8\, E_{135643224} + (-1) (q-1) (q+1) q^9\, E_{136432254}
  \\
  + (-1) (q-1) (q+1) q^7\, E_{143243256} + (q-1) (q+1) q^6\, E_{143254326} + (-1) (q-1) (q+1) q^7 (q^2+1)\, E_{143256432}
  \\
  + (q-1) (q+1) q^8\, E_{143264325} + (q-1) (q+1) q^8\, E_{144323256} + (-1) (q-1) (q+1) q^7\, E_{145432326}
  \\
  + (q-1) (q+1) q^8\, E_{145643232} + (-1) (q-1) (q+1) q^9\, E_{146432325} + (q-1) (q+1) q^8\, E_{154326432}
  \\
  + (-1) (q-1) (q+1) q^9\, E_{156432432} + (q-1) (q+1) q^{10}\, E_{164325432} + (q-1) (q+1) q^6\, E_{212345643}
  \\
  + (-1) (q-1) (q+1) q^5\, E_{212354643} + (q-1) (q+1) q^6\, E_{212356443} + (-1) (q-1) (q+1) q^7\, E_{212364543}
  \\
  + (-1) (q-1) (q+1) q^5\, E_{212434356} + (q-1) (q+1) q^4\, E_{212435436} + (-1) (q-1) (q+1) q^5 (q^2+1)\, E_{212435643}
  \\
  + (q-1) (q+1) q^6\, E_{212436435} + (q-1) (q+1) q^6\, E_{212443356} + (-1) (q-1) (q+1) q^5\, E_{212454336}
  \\
  + (q-1) (q+1) q^6\, E_{212456433} + (-1) (q-1) (q+1) q^7\, E_{212464335} + (q-1) (q+1) q^6\, E_{212543643}
  \\
  + (-1) (q-1) (q+1) q^7\, E_{212564343} + (q-1) (q+1) q^8\, E_{212643543} + (-1) (q-1) (q+1) q^5\, E_{213234564}
  \\
  + (q-1) (q+1) q^4\, E_{213235464} + (-1) (q-1) (q+1) q^5\, E_{213235644} + (q-1) (q+1) q^6\, E_{213236454}
  \\
  + (q-1) (q+1) q^4\, E_{213243456} + (-1) (q-1) (q+1) q^3\, E_{213243546} + (q-1) (q+1) q^2 (q^2+1)\, E_{213243564}
  \\
  + (-1) (q-1) (q+1) q^3\, E_{213243645} + (-1) (q-1) (q+1) q^5\, E_{213244356} + (q-1) (q+1) q^4\, E_{213245436}
  \\
  + (-1) (q-1) q (q+1) (q^4 + 1)\, E_{213245643} + (q-1) (q+1) q^2 (q^4 - q^2+1)\, E_{213246435} + (-1) (q-1) (q+1) q^3\, E_{213254364}
  \\
  + (q-1) (q+1) q^2\, E_{213254643} + (q-1) (q+1) q^4\, E_{213256434} + (-1) (q-1) (q+1) q^3\, E_{213256443}
  \\
  + (-1) (q-1)^{2} (q+1)^{2} q^4\, E_{213264345} + (-1) (q-1) (q+1) q^3\, E_{213264354} + (q-1)^{2} (q+1)^{2} q^3\, E_{213264435}
  \\
  + (q-1) (q+1) q^2\, E_{213264543} + (q-1) (q+1) q^6\, E_{213324564} + (-1) (q-1) (q+1) q^5\, E_{213325464}
  \\
  + (q-1) (q+1) q^6\, E_{213325644} + (-1) (q-1) (q+1) q^7\, E_{213326454} + (-1) (q-1) (q+1) q^5\, E_{213432456}
  \\
  + (q-1) (q+1) q^4\, E_{213432546} + (-1) (q-1) q (q+1) (q^2+1) (q^4 - q^2+1)\, E_{213432564} + (q-1) (q+1) q^2 (q^4 - q^2+1)\, E_{213432645}
  \\
  + (q-1) (q+1) q^6\, E_{213443256} + (-1) (q-1) (q+1) q^5\, E_{213454326} + (q-1) (q+1) (q^2+1) (q^4 - q^2+1)\, E_{213456432}
  \\
  + (-1) (q-1) q (q+1) (q^6 - q^2+1)\, E_{213464325} + (q-1) (q+1) q^2 (q^4 - q^2+1)\, E_{213543264} + (-1) (q-1) q (q+1)\, E_{213546432}
  \\
  + (-1) (q-1) (q+1) q^3 (q^4 - q^2+1)\, E_{213564324} + (q-1) (q+1) q^2\, E_{213564432} + (q-1)^{2} (q+1)^{2} q^3\, E_{213643245}
  \\
  + (q-1) (q+1) q^2 (q^6 - q^4 + 1)\, E_{213643254} + (-1) (q-1)^{2} q^2 (q+1)^{2}\, E_{213644325} + (-1) (q-1) q (q+1)\, E_{213645432}
  \\
  + (q-1)^{2} (q+1)^{2} q^3 (q^2+1)\, E_{214323564} + (-1) (q-1)^{2} (q+1)^{2} q^4\, E_{214323645} + (q-1) (q+1) q^4\, E_{214324356}
  \\
  + (-1) (q-1) (q+1) q^3\, E_{214325436} + (q-1) (q+1) q^2 (q^2+1)\, E_{214325643} + (-1) (q-1) (q+1) q^3\, E_{214326435}
  \\
  + (-1) (q-1)^{2} q^2 (q+1)^{2} (q^2+1)\, E_{214332564} + (q-1)^{2} (q+1)^{2} q^3\, E_{214332645} + (-1) (q-1) (q+1) q^3\, E_{214343256}
  \\
  + (q-1) (q+1) q^2\, E_{214354326} + (-1) (q-1) q (q+1) (q^2+1)\, E_{214356432} + (q-1) (q+1) q^2\, E_{214364325}
  \\
  + (-1) (q-1) (q+1) q^5\, E_{214432356} + (q-1) (q+1) q^4\, E_{214433256} + (q-1) (q+1) q^4\, E_{214543236}
  \\
  + (-1) (q-1) (q+1) q^3\, E_{214543326} + (-1) (q-1) (q+1) q^5\, E_{214564323} + (q-1) (q+1) q^4\, E_{214564332}
  \\
  + (q-1) (q+1) q^6\, E_{214643235} + (-1) (q-1) (q+1) q^5\, E_{214643325} + (-1) (q-1)^{2} (q+1)^{2} q^4\, E_{215432364}
  \\
  + (-1) (q-1) (q+1) q^3\, E_{215432643} + (q-1)^{2} (q+1)^{2} q^3\, E_{215433264} + (q-1) (q+1) q^2\, E_{215436432}
  \\
  + (q-1)^{2} (q+1)^{2} q^5\, E_{215643234} + (q-1) (q+1) q^4\, E_{215643243} + (-1) (q-1)^{2} (q+1)^{2} q^4\, E_{215643324}
  \\
  + (-1) (q-1) (q+1) q^3\, E_{215643432} + (q-1)^{2} (q+1)^{2} q^5\, E_{216432345} + (-1) (q-1)^{2} (q+1)^{2} q^4 (q^2+1)\, E_{216432354}
  \\
  + (-1) (q-1)^{2} (q+1)^{2} q^4\, E_{216432435} + (-1) (q-1) (q+1) q^3\, E_{216432543} + (-1) (q-1)^{2} (q+1)^{2} q^4\, E_{216433245}
  \\
  + (q-1)^{2} (q+1)^{2} q^3 (q^2+1)\, E_{216433254} + (q-1)^{2} (q+1)^{2} q^3\, E_{216434325} + (q-1) (q+1) q^2\, E_{216435432}
  \\
  + (-1) (q-1) (q+1) q^7\, E_{221345643} + (q-1) (q+1) q^6\, E_{221354643} + (-1) (q-1) (q+1) q^7\, E_{221356443}
  \\
  + (q-1) (q+1) q^8\, E_{221364543} + (q-1) (q+1) q^6\, E_{221434356} + (-1) (q-1) (q+1) q^5\, E_{221435436}
  \\
  + (q-1) (q+1) q^6 (q^2+1)\, E_{221435643} + (-1) (q-1) (q+1) q^7\, E_{221436435} + (-1) (q-1) (q+1) q^7\, E_{221443356}
  \\
  + (q-1) (q+1) q^6\, E_{221454336} + (-1) (q-1) (q+1) q^7\, E_{221456433} + (q-1) (q+1) q^8\, E_{221464335}
  \\
  + (-1) (q-1) (q+1) q^7\, E_{221543643} + (q-1) (q+1) q^8\, E_{221564343} + (-1) (q-1) (q+1) q^9\, E_{221643543}
  \\
  + (q-1) (q+1) q^6\, E_{232134564} + (-1) (q-1) (q+1) q^5\, E_{232135464} + (q-1) (q+1) q^6\, E_{232135644}
  \\
  + (-1) (q-1) (q+1) q^7\, E_{232136454} + (-1) (q-1) (q+1) q^5\, E_{232143456} + (q-1) (q+1) q^4\, E_{232143546}
  \\
  + (-1) (q-1) q (q+1) (q^2+1) (q^4 - q^2+1)\, E_{232143564} + (q-1) (q+1) q^2 (q^4 - q^2+1)\, E_{232143645} + (q-1) (q+1) q^6\, E_{232144356}
  \\
  + (-1) (q-1) (q+1) q^5\, E_{232145436} + (q-1) (q+1) (q^8 + 1)\, E_{232145643} + (-1) (q-1) q (q+1) (q^6 - q^2+1)\, E_{232146435}
  \\
  + (q-1) (q+1) q^2 (q^4 - q^2+1)\, E_{232154364} + (-1) (q-1) q (q+1) (q^6 - q^4 + 1)\, E_{232154643} + (-1) (q-1) (q+1) q^3 (q^4 - q^2+1)\, E_{232156434}
  \\
  + (q-1) (q+1) q^2 (q^6 - q^4 + 1)\, E_{232156443} + (q-1)^{2} (q+1)^{2} q^3\, E_{232164345} + (q-1) (q+1) q^2 (q^6 - q^4 + 1)\, E_{232164354}
  \\
  + (-1) (q-1)^{2} q^2 (q+1)^{2}\, E_{232164435} + (-1) (q-1) q (q+1) (q^8 - q^6 + 1)\, E_{232164543} + (-1) (q-1) (q+1) q^7\, E_{233214564}
  \\
  + (q-1) (q+1) q^6\, E_{233215464} + (-1) (q-1) (q+1) q^7\, E_{233215644} + (q-1) (q+1) q^8\, E_{233216454}
  \\
  + (q-1) (q+1) q^6\, E_{234321456} + (-1) (q-1) (q+1) q^5\, E_{234321546} + (q-1) (q+1) (q^2+1) (q^6 - q^2+1)\, E_{234321564}
  \\
  + (-1) (q-1) q (q+1) (q^6 - q^2+1)\, E_{234321645} + (-1) (q-1) (q+1) q^7\, E_{234432156} + (q-1) (q+1) q^6\, E_{234543216}
  \\
  + ((-1) (q-1) (q+1))/(q)\, E_{234564321} + (q-1) (q+1) (q^8 - q^2+1)\, E_{234643215} + (-1) (q-1) q (q+1) (q^6 - q^2+1)\, E_{235432164}
  \\
  + (-1) (q-1)^{2} (q+1)^{2} (q^2 - q+1) (q^2 + q+1)\, E_{235464321} + (q-1) (q+1) q^2 (q^6 - q^2+1)\, E_{235643214} + q (q-1)^{2} (q+1)^{2} (q^2 - q+1) (q^2 + q+1)\, E_{235644321}
  \\
  + (-1) (q-1)^{2} q^2 (q+1)^{2}\, E_{236432145} + (-1) (q-1) q (q+1) (q^8 - q^4 + 1)\, E_{236432154} + q (q-1)^{2} (q+1)^{2}\, E_{236443215}
  \\
  + (-1) (q-1)^{2} (q+1)^{2} (q^2+1) (q^4 + 1)\, E_{236454321} + (-1) (q-1)^{2} q^2 (q+1)^{2} (q^2+1)\, E_{243213564} + (q-1)^{2} (q+1)^{2} q^3\, E_{243213645}
  \\
  + (-1) (q-1) (q+1) q^3 (q^4 - q^2+1)\, E_{243214356} + (q-1) (q+1) q^2 (q^4 - q^2+1)\, E_{243215436} + (-1) (q-1) q (q+1) (q^2+1) (q^6 - q^4 + 1)\, E_{243215643}
  \\
  + (q-1) (q+1) q^2 (q^6 - q^4 + 1)\, E_{243216435} + q (q-1)^{2} (q+1)^{2} (q^2+1)\, E_{243321564} + (-1) (q-1)^{2} q^2 (q+1)^{2}\, E_{243321645}
  \\
  + (-1) (q-1)^{2} q^2 (q+1)^{2} (q^2+1)\, E_{243432156} + q (q-1)^{2} (q+1)^{2} (q^2+1)\, E_{243543216} + (-1) (q-1)^{2} (q+1)^{2} (q^2 - q+1) (q^2+1) (q^2 + q+1)\, E_{243564321}
  \\
  + q (q-1)^{2} (q+1)^{2} (q^2 - q+1) (q^2 + q+1)\, E_{243643215} + (q-1) (q+1) q^4 (q^4 - q^2+1)\, E_{244321356} + (q-1)^{2} (q+1)^{2} q^3 (q^2+1)\, E_{244332156}
  \\
  + (-1) (q-1) (q+1) q^3 (q^4 - q^2+1)\, E_{245432136} + (-1) (q-1)^{2} q^2 (q+1)^{2} (q^2+1)\, E_{245433216} + (q-1) (q+1) q^4 (q^4 - q^2+1)\, E_{245643213}
  \\
  + (q-1)^{2} (q+1)^{2} q^3 (q^2+1)\, E_{245643321} + (-1) (q-1) (q+1) q^5 (q^4 - q^2+1)\, E_{246432135} + (-1) (q-1)^{2} (q+1)^{2} q^4 (q^2+1)\, E_{246433215}
  \\
  + (q-1)^{2} (q+1)^{2} q^3\, E_{254321364} + (q-1) (q+1) q^2 (q^6 - q^4 + 1)\, E_{254321643} + (-1) (q-1)^{2} q^2 (q+1)^{2}\, E_{254332164}
  \\
  + q (q-1)^{2} (q+1)^{2} (q^2 - q+1) (q^2 + q+1)\, E_{254364321} + (-1) (q-1)^{2} (q+1)^{2} q^4\, E_{256432134} + (-1) (q-1) (q+1) q^3 (q^6 - q^4 + 1)\, E_{256432143}
  \\
  + (q-1)^{2} (q+1)^{2} q^3\, E_{256433214} + (-1) (q-1)^{2} q^2 (q+1)^{2} (q^2 - q+1) (q^2 + q+1)\, E_{256434321} + (-1) (q-1)^{2} (q+1)^{2} q^4\, E_{264321345}
  \\
  + (q-1)^{2} (q+1)^{2} q^3 (q^2+1)\, E_{264321354} + (q-1)^{2} (q+1)^{2} q^3\, E_{264321435} + (q-1) (q+1) q^2 (q^8 - q^6 + 1)\, E_{264321543}
  \\
  + (q-1)^{2} (q+1)^{2} q^3\, E_{264332145} + (-1) (q-1)^{2} q^2 (q+1)^{2} (q^2+1)\, E_{264332154} + (-1) (q-1)^{2} q^2 (q+1)^{2}\, E_{264343215}
  \\
  + q (q-1)^{2} (q+1)^{2} (q^2+1) (q^4 + 1)\, E_{264354321} + (q-1)^{2} (q+1)^{2} q^3 (q^2+1)\, E_{321243564} + (-1) (q-1)^{2} (q+1)^{2} q^4\, E_{321243645}
  \\
  + (-1) (q-1)^{2} q^2 (q+1)^{2} (q^2 - q+1) (q^2 + q+1)\, E_{321245643} + (q-1)^{2} (q+1)^{2} q^3\, E_{321246435} + (-1) (q-1)^{2} (q+1)^{2} q^4\, E_{321254364}
  \\
  + (q-1)^{2} (q+1)^{2} q^3 (q^2+1)\, E_{321254643} + (q-1)^{2} (q+1)^{2} q^5\, E_{321256434} + (-1) (q-1)^{2} (q+1)^{2} q^4 (q^2+1)\, E_{321256443}
  \\
  + (q-1)^{2} (q+1)^{2} q^5\, E_{321264345} + (-1) (q-1)^{2} (q+1)^{2} q^4 (q^2+1)\, E_{321264354} + (-1) (q-1)^{2} (q+1)^{2} q^4\, E_{321264435}
  \\
  + (q-1)^{2} (q+1)^{2} q^3 (q^2 - q+1) (q^2 + q+1)\, E_{321264543} + (-1) (q-1) (q+1) q^5\, E_{321324564} + (q-1) (q+1) q^4\, E_{321325464}
  \\
  + (-1) (q-1) (q+1) q^5\, E_{321325644} + (q-1) (q+1) q^6\, E_{321326454} + (q-1) (q+1) q^4\, E_{321432456}
  \\
  + (-1) (q-1) (q+1) q^3\, E_{321432546} + (q-1) (q+1) q^2 (q^2+1)\, E_{321432564} + (-1) (q-1) (q+1) q^3\, E_{321432645}
  \\
  + (-1) (q-1) (q+1) q^5\, E_{321443256} + (q-1) (q+1) q^4\, E_{321454326} + (-1) (q-1) q (q+1) (q^4 + 1)\, E_{321456432}
  \\
  + (q-1) (q+1) q^2 (q^4 - q^2+1)\, E_{321464325} + (-1) (q-1) (q+1) q^3\, E_{321543264} + (q-1) (q+1) q^2\, E_{321546432}
  \\
  + (q-1) (q+1) q^4\, E_{321564324} + (-1) (q-1) (q+1) q^3\, E_{321564432} + (-1) (q-1)^{2} (q+1)^{2} q^4\, E_{321643245}
  \\
  + (-1) (q-1) (q+1) q^3\, E_{321643254} + (q-1)^{2} (q+1)^{2} q^3\, E_{321644325} + (q-1) (q+1) q^2\, E_{321645432}
  \\
  + (-1) (q-1)^{2} q^2 (q+1)^{2} (q^2+1)\, E_{322143564} + (q-1)^{2} (q+1)^{2} q^3\, E_{322143645} + q (q-1)^{2} (q+1)^{2} (q^2 - q+1) (q^2 + q+1)\, E_{322145643}
  \\
  + (-1) (q-1)^{2} q^2 (q+1)^{2}\, E_{322146435} + (q-1)^{2} (q+1)^{2} q^3\, E_{322154364} + (-1) (q-1)^{2} q^2 (q+1)^{2} (q^2+1)\, E_{322154643}
  \\
  + (-1) (q-1)^{2} (q+1)^{2} q^4\, E_{322156434} + (q-1)^{2} (q+1)^{2} q^3 (q^2+1)\, E_{322156443} + (-1) (q-1)^{2} (q+1)^{2} q^4\, E_{322164345}
  \\
  + (q-1)^{2} (q+1)^{2} q^3 (q^2+1)\, E_{322164354} + (q-1)^{2} (q+1)^{2} q^3\, E_{322164435} + (-1) (q-1)^{2} q^2 (q+1)^{2} (q^2 - q+1) (q^2 + q+1)\, E_{322164543}
  \\
  + (-1) (q-1)^{2} (q+1)^{2} q^4\, E_{323214564} + (q-1)^{2} (q+1)^{2} q^3\, E_{323215464} + (-1) (q-1)^{2} (q+1)^{2} q^4\, E_{323215644}
  \\
  + (q-1)^{2} (q+1)^{2} q^5\, E_{323216454} + (q-1)^{2} (q+1)^{2} q^3\, E_{324321456} + (-1) (q-1)^{2} q^2 (q+1)^{2}\, E_{324321546}
  \\
  + q (q-1)^{2} (q+1)^{2} (q^2+1)^{2}\, E_{324321564} + (-1) (q-1)^{2} q^2 (q+1)^{2} (q^2+1)\, E_{324321645} + (-1) (q-1)^{2} (q+1)^{2} q^4\, E_{324432156}
  \\
  + (q-1)^{2} (q+1)^{2} q^3\, E_{324543216} + (-1) (q-1)^{2} (q+1)^{2} (q^6 + 2 q^4 + q^2+1)\, E_{324564321} + q (q-1)^{2} (q+1)^{2} (q^4 + 1)\, E_{324643215}
  \\
  + (-1) (q-1)^{2} q^2 (q+1)^{2} (q^2+1)\, E_{325432164} + q (q-1)^{2} (q+1)^{2} (q^2 - q+1) (q^2 + q+1)\, E_{325464321} + (q-1)^{2} (q+1)^{2} q^3 (q^2+1)\, E_{325643214}
  \\
  + (-1) (q-1)^{2} q^2 (q+1)^{2} (q^2 - q+1) (q^2 + q+1)\, E_{325644321} + (q-1)^{2} (q+1)^{2} q^3\, E_{326432145} + (-1) (q-1)^{2} q^2 (q+1)^{2} (q^2 - q+1) (q^2 + q+1)\, E_{326432154}
  \\
  + (-1) (q-1)^{2} q^2 (q+1)^{2}\, E_{326443215} + q (q-1)^{2} (q+1)^{2} (q^2+1) (q^4 + 1)\, E_{326454321} + (q-1) (q+1) q^6\, E_{332124564}
  \\
  + (-1) (q-1) (q+1) q^5\, E_{332125464} + (q-1) (q+1) q^6\, E_{332125644} + (-1) (q-1) (q+1) q^7\, E_{332126454}
  \\
  + (q-1)^{2} (q+1)^{2} q^5\, E_{332214564} + (-1) (q-1)^{2} (q+1)^{2} q^4\, E_{332215464} + (q-1)^{2} (q+1)^{2} q^5\, E_{332215644}
  \\
  + (-1) (q-1)^{2} (q+1)^{2} q^6\, E_{332216454} + (-1) (q-1) (q+1) q^5\, E_{343212456} + (q-1) (q+1) q^4\, E_{343212546}
  \\
  + (-1) (q-1) (q+1) q^5 (q^2+1)\, E_{343212564} + (q-1) (q+1) q^6\, E_{343212645} + (-1) (q-1)^{2} (q+1)^{2} q^4\, E_{343221456}
  \\
  + (q-1)^{2} (q+1)^{2} q^3\, E_{343221546} + (-1) (q-1)^{2} (q+1)^{2} q^4 (q^2+1)\, E_{343221564} + (q-1)^{2} (q+1)^{2} q^5\, E_{343221645}
  \\
  + (q-1) (q+1) q^6\, E_{344321256} + (q-1)^{2} (q+1)^{2} q^5\, E_{344322156} + (-1) (q-1) (q+1) q^5\, E_{345432126}
  \\
  + (-1) (q-1)^{2} (q+1)^{2} q^4\, E_{345432216} + (q-1) (q+1) q^6\, E_{345643212} + (q-1)^{2} (q+1)^{2} q^5\, E_{345643221}
  \\
  + (-1) (q-1) (q+1) q^7\, E_{346432125} + (-1) (q-1)^{2} (q+1)^{2} q^6\, E_{346432215} + (q-1) (q+1) q^6\, E_{354321264}
  \\
  + (q-1)^{2} (q+1)^{2} q^5\, E_{354322164} + (-1) (q-1) (q+1) q^7\, E_{356432124} + (-1) (q-1)^{2} (q+1)^{2} q^6\, E_{356432214}
  \\
  + (q-1) (q+1) q^8\, E_{364321254} + (q-1)^{2} (q+1)^{2} q^7\, E_{364322154} + (-1) (q-1)^{2} (q+1)^{2} q^4 (q^2+1)\, E_{432123564}
  \\
  + (q-1)^{2} (q+1)^{2} q^5\, E_{432123645} + (q-1)^{2} (q+1)^{2} q^5\, E_{432124356} + (-1) (q-1)^{2} (q+1)^{2} q^4\, E_{432125436}
  \\
  + (q-1)^{2} (q+1)^{2} q^3 (q^2+1)^{2}\, E_{432125643} + (-1) (q-1)^{2} (q+1)^{2} q^4 (q^2+1)\, E_{432126435} + (q-1)^{2} (q+1)^{2} q^3 (q^2+1)\, E_{432132564}
  \\
  + (-1) (q-1)^{2} (q+1)^{2} q^4\, E_{432132645} + (q-1) (q+1) q^4\, E_{432143256} + (-1) (q-1) (q+1) q^3\, E_{432154326}
  \\
  + (q-1) (q+1) q^2 (q^2+1)\, E_{432156432} + (-1) (q-1) (q+1) q^3\, E_{432164325} + (q-1)^{2} (q+1)^{2} q^3 (q^2+1)\, E_{432213564}
  \\
  + (-1) (q-1)^{2} (q+1)^{2} q^4\, E_{432213645} + (-1) (q-1)^{2} (q+1)^{2} q^4\, E_{432214356} + (q-1)^{2} (q+1)^{2} q^3\, E_{432215436}
  \\
  + (-1) (q-1)^{2} q^2 (q+1)^{2} (q^2+1)^{2}\, E_{432215643} + (q-1)^{2} (q+1)^{2} q^3 (q^2+1)\, E_{432216435} + (-1) (q-1)^{2} q^2 (q+1)^{2} (q^2+1)\, E_{432321564}
  \\
  + (q-1)^{2} (q+1)^{2} q^3\, E_{432321645} + (q-1)^{2} (q+1)^{2} q^3 (q^2+1)\, E_{432432156} + (-1) (q-1)^{2} q^2 (q+1)^{2} (q^2+1)\, E_{432543216}
  \\
  + q (q-1)^{2} (q+1)^{2} (q^2 - q+1) (q^2+1) (q^2 + q+1)\, E_{432564321} + (-1) (q-1)^{2} q^2 (q+1)^{2} (q^2 - q+1) (q^2 + q+1)\, E_{432643215} + (-1) (q-1)^{2} (q+1)^{2} q^6\, E_{443212356}
  \\
  + (-1) (q-1) (q+1) q^5\, E_{443213256} + (q-1)^{2} (q+1)^{2} q^5\, E_{443221356} + (-1) (q-1)^{2} (q+1)^{2} q^4 (q^2+1)\, E_{443232156}
  \\
  + (q-1)^{2} (q+1)^{2} q^5\, E_{454321236} + (q-1) (q+1) q^4\, E_{454321326} + (-1) (q-1)^{2} (q+1)^{2} q^4\, E_{454322136}
  \\
  + (q-1)^{2} (q+1)^{2} q^3 (q^2+1)\, E_{454323216} + (-1) (q-1)^{2} (q+1)^{2} q^6\, E_{456432123} + (-1) (q-1) (q+1) q^5\, E_{456432132}
  \\
  + (q-1)^{2} (q+1)^{2} q^5\, E_{456432213} + (-1) (q-1)^{2} (q+1)^{2} q^4 (q^2+1)\, E_{456432321} + (q-1)^{2} (q+1)^{2} q^7\, E_{464321235}
  \\
  + (q-1) (q+1) q^6\, E_{464321325} + (-1) (q-1)^{2} (q+1)^{2} q^6\, E_{464322135} + (q-1)^{2} (q+1)^{2} q^5 (q^2+1)\, E_{464323215}
  \\
  + (q-1)^{2} (q+1)^{2} q^5\, E_{543212364} + (-1) (q-1)^{2} (q+1)^{2} q^4 (q^2+1)\, E_{543212643} + (-1) (q-1)^{2} (q+1)^{2} q^4\, E_{543213264}
  \\
  + (-1) (q-1) (q+1) q^3\, E_{543216432} + (-1) (q-1)^{2} (q+1)^{2} q^4\, E_{543221364} + (q-1)^{2} (q+1)^{2} q^3 (q^2+1)\, E_{543221643}
  \\
  + (q-1)^{2} (q+1)^{2} q^3\, E_{543232164} + (-1) (q-1)^{2} q^2 (q+1)^{2} (q^2 - q+1) (q^2 + q+1)\, E_{543264321} + (-1) (q-1)^{2} (q+1)^{2} q^6\, E_{564321234}
  \\
  + (q-1)^{2} (q+1)^{2} q^5 (q^2+1)\, E_{564321243} + (q-1)^{2} (q+1)^{2} q^5\, E_{564321324} + (q-1) (q+1) q^4\, E_{564321432}
  \\
  + (q-1)^{2} (q+1)^{2} q^5\, E_{564322134} + (-1) (q-1)^{2} (q+1)^{2} q^4 (q^2+1)\, E_{564322143} + (-1) (q-1)^{2} (q+1)^{2} q^4\, E_{564323214}
  \\
  + (q-1)^{2} (q+1)^{2} q^3 (q^2 - q+1) (q^2 + q+1)\, E_{564324321} + (-1) (q-1)^{2} (q+1)^{2} q^6\, E_{643212345} + (q-1)^{2} (q+1)^{2} q^5 (q^2+1)\, E_{643212354}
  \\
  + (q-1)^{2} (q+1)^{2} q^5\, E_{643212435} + (-1) (q-1)^{2} (q+1)^{2} q^4 (q^2 - q+1) (q^2 + q+1)\, E_{643212543} + (q-1)^{2} (q+1)^{2} q^5\, E_{643213245}
  \\
  + (-1) (q-1)^{2} (q+1)^{2} q^4 (q^2+1)\, E_{643213254} + (-1) (q-1)^{2} (q+1)^{2} q^4\, E_{643214325} + (-1) (q-1) (q+1) q^3\, E_{643215432}
  \\
  + (q-1)^{2} (q+1)^{2} q^5\, E_{643221345} + (-1) (q-1)^{2} (q+1)^{2} q^4 (q^2+1)\, E_{643221354} + (-1) (q-1)^{2} (q+1)^{2} q^4\, E_{643221435}
  \\
  + (q-1)^{2} (q+1)^{2} q^3 (q^2 - q+1) (q^2 + q+1)\, E_{643221543} + (-1) (q-1)^{2} (q+1)^{2} q^4\, E_{643232145} + (q-1)^{2} (q+1)^{2} q^3 (q^2+1)\, E_{643232154}
  \\
  + (q-1)^{2} (q+1)^{2} q^3\, E_{643243215} + (-1) (q-1)^{2} q^2 (q+1)^{2} (q^2+1) (q^4 + 1)\, E_{643254321}
\end{multline*}

\normalsize

\item The element $\widetilde{C}_{10}$ consists of the $\Uq(\so_{10})$ central element embedded in $\Uq(\so_{12})$ via the sub-diagram $\{\alpha_2, \ldots, \alpha_6\}$, comprising 10 diagonal terms and 44 off-diagonal terms involving the weights $\pm L_2, \ldots, \pm L_6$.

The 10 diagonal terms are
\[
    q^{8}\, K_{2L_2} + q^{6}\, K_{2L_3} + q^{4}\, K_{2L_4} + q^{2}\, K_{2L_5} + K_{2L_6} + K_{-2L_6} + q^{-2}\, K_{-2L_5} + q^{-4}\, K_{-2L_4} + q^{-6}\, K_{-2L_3} + q^{-8}\, K_{-2L_2}.
\]
The 44 off-diagonal terms each take the form $D_{\mu,\lambda}\, e^*_{\mu,\lambda}\, K_{-\mu-\lambda}\, \omega(\tau(e^*_{\mu,\lambda}))$. The scalar prefactors are:
\begin{center}
\begin{tabular}{lcc|lcc}
\hline
Pair $(\mu,\lambda)$ & $D_{\mu,\lambda}$ & $K_{-\mu-\lambda}$ & Pair $(\mu,\lambda)$ & $D_{\mu,\lambda}$ & $K_{-\mu-\lambda}$ \\
\hline
$(L_2, L_3)$ & $q^{-7}$ & $K_{-L_2-L_3}$ & $(-L_3, -L_2)$ & $q^{7}$ & $K_{L_2+L_3}$ \\
$(L_2, L_4)$ & $q^{-7}$ & $K_{-L_2-L_4}$ & $(-L_4, -L_2)$ & $q^{5}$ & $K_{L_2+L_4}$ \\
$(L_2, L_5)$ & $q^{-7}$ & $K_{-L_2-L_5}$ & $(-L_5, -L_2)$ & $q^{3}$ & $K_{L_2+L_5}$ \\
$(L_2, L_6)$ & $q^{-7}$ & $K_{-L_2-L_6}$ & $(-L_6, -L_2)$ & $q$ & $K_{L_2+L_6}$ \\
$(L_2, -L_6)$ & $q^{-7}$ & $K_{-L_2+L_6}$ & $(L_6, -L_2)$ & $q$ & $K_{-L_6+L_2}$ \\
$(L_2, -L_5)$ & $q^{-7}$ & $K_{-L_2+L_5}$ & $(L_5, -L_2)$ & $q^{-1}$ & $K_{-L_5+L_2}$ \\
$(L_2, -L_4)$ & $q^{-7}$ & $K_{-L_2+L_4}$ & $(L_4, -L_2)$ & $q^{-3}$ & $K_{-L_4+L_2}$ \\
$(L_2, -L_3)$ & $q^{-7}$ & $K_{-L_2+L_3}$ & $(L_3, -L_2)$ & $q^{-5}$ & $K_{-L_3+L_2}$ \\
$(L_2, -L_2)$ & $q^{-6}$ & $K_0 = 1$ & & & \\
\hline
$(L_3, L_4)$ & $q^{-5}$ & $K_{-L_3-L_4}$ & $(-L_4, -L_3)$ & $q^{5}$ & $K_{L_3+L_4}$ \\
$(L_3, L_5)$ & $q^{-5}$ & $K_{-L_3-L_5}$ & $(-L_5, -L_3)$ & $q^{3}$ & $K_{L_3+L_5}$ \\
$(L_3, L_6)$ & $q^{-5}$ & $K_{-L_3-L_6}$ & $(-L_6, -L_3)$ & $q$ & $K_{L_3+L_6}$ \\
$(L_3, -L_6)$ & $q^{-5}$ & $K_{-L_3+L_6}$ & $(L_6, -L_3)$ & $q$ & $K_{-L_6+L_3}$ \\
$(L_3, -L_5)$ & $q^{-5}$ & $K_{-L_3+L_5}$ & $(L_5, -L_3)$ & $q^{-1}$ & $K_{-L_5+L_3}$ \\
$(L_3, -L_4)$ & $q^{-5}$ & $K_{-L_3+L_4}$ & $(L_4, -L_3)$ & $q^{-3}$ & $K_{-L_4+L_3}$ \\
$(L_3, -L_3)$ & $q^{-4}$ & $K_0 = 1$ & & & \\
\hline
$(L_4, L_5)$ & $q^{-3}$ & $K_{-L_4-L_5}$ & $(-L_5, -L_4)$ & $q^{3}$ & $K_{L_4+L_5}$ \\
$(L_4, L_6)$ & $q^{-3}$ & $K_{-L_4-L_6}$ & $(-L_6, -L_4)$ & $q$ & $K_{L_4+L_6}$ \\
$(L_4, -L_6)$ & $q^{-3}$ & $K_{-L_4+L_6}$ & $(L_6, -L_4)$ & $q$ & $K_{-L_6+L_4}$ \\
$(L_4, -L_5)$ & $q^{-3}$ & $K_{-L_4+L_5}$ & $(L_5, -L_4)$ & $q^{-1}$ & $K_{-L_5+L_4}$ \\
$(L_4, -L_4)$ & $q^{-2}$ & $K_0 = 1$ & & & \\
\hline
$(L_5, L_6)$ & $q^{-1}$ & $K_{-L_5-L_6}$ & $(-L_6, -L_5)$ & $q$ & $K_{L_5+L_6}$ \\
$(L_5, -L_6)$ & $q^{-1}$ & $K_{-L_5+L_6}$ & $(L_6, -L_5)$ & $q$ & $K_{-L_6+L_5}$ \\
$(L_5, -L_5)$ & $q^{0}$ & $K_0 = 1$ & & & \\
\hline
\end{tabular}
\end{center}
The dual elements $e^*_{\mu,\lambda}$ for these 44 pairs are: \footnotesize

\allowdisplaybreaks

%% Pair (L_2, L_3), 1 terms
$e^*_{(L_2, L_3)} = ((-1)(q-1)(q+1))/(q)\,E_{2}$

%% Pair (L_2, L_4), 2 terms
$e^*_{(L_2, L_4)} = ((-1)(q-1)(q+1))/(q)\,E_{23} + (q-1)(q+1)\,E_{32}$

%% Pair (L_2, L_5), 4 terms
\begin{multline*}
e^*_{(L_2, L_5)} = 
  ((-1)(q-1)(q+1))/(q)\,E_{234} + (q-1)(q+1)\,E_{243} + (q-1)(q+1)\,E_{324}
  \\
  + (-1)(q-1)q(q+1)\,E_{432}
\end{multline*}

%% Pair (L_2, L_6), 8 terms
\begin{multline*}
e^*_{(L_2, L_6)} = 
  ((-1)(q-1)(q+1))/(q)\,E_{2345} + (q-1)(q+1)\,E_{2354} + (q-1)(q+1)\,E_{2435}
  \\
  + (-1)(q-1)q(q+1)\,E_{2543} + (q-1)(q+1)\,E_{3245} + (-1)(q-1)q(q+1)\,E_{3254}
  \\
  + (-1)(q-1)q(q+1)\,E_{4325} + (q-1)(q+1)q^2\,E_{5432}
\end{multline*}

%% Pair (L_2, -L_6), 8 terms
\begin{multline*}
e^*_{(L_2, -L_6)} = 
  ((-1)(q-1)(q+1))/(q)\,E_{2346} + (q-1)(q+1)\,E_{2364} + (q-1)(q+1)\,E_{2436}
  \\
  + (-1)(q-1)q(q+1)\,E_{2643} + (q-1)(q+1)\,E_{3246} + (-1)(q-1)q(q+1)\,E_{3264}
  \\
  + (-1)(q-1)q(q+1)\,E_{4326} + (q-1)(q+1)q^2\,E_{6432}
\end{multline*}

%% Pair (L_2, -L_5), 16 terms
\begin{multline*}
e^*_{(L_2, -L_5)} = 
  ((-1)(q-1)(q+1))/(q)\,E_{23456} + (q-1)(q+1)\,E_{23546} + (-1)(q-1)q(q+1)\,E_{23564}
  \\
  + (q-1)(q+1)\,E_{23645} + (q-1)(q+1)\,E_{24356} + (-1)(q-1)q(q+1)\,E_{25436}
  \\
  + (q-1)(q+1)q^2\,E_{25643} + (-1)(q-1)q(q+1)\,E_{26435} + (q-1)(q+1)\,E_{32456}
  \\
  + (-1)(q-1)q(q+1)\,E_{32546} + (q-1)(q+1)q^2\,E_{32564} + (-1)(q-1)q(q+1)\,E_{32645}
  \\
  + (-1)(q-1)q(q+1)\,E_{43256} + (q-1)(q+1)q^2\,E_{54326} + (-1)(q-1)(q+1)q^3\,E_{56432}
  \\
  + (q-1)(q+1)q^2\,E_{64325}
\end{multline*}

%% Pair (L_2, -L_4), 36 terms
\begin{multline*}
e^*_{(L_2, -L_4)} = 
  ((-1)(q-1)(q+1))/(q)\,E_{234564} + (-1)(q-1)^{2}(q+1)^{2}\,E_{234645} + (q-1)(q+1)\,E_{235464}
  \\
  + (-1)(q-1)q(q+1)\,E_{235644} + q(q-1)^{2}(q+1)^{2}\,E_{236445} + (q-1)(q+1)\,E_{236454}
  \\
  + (q-1)(q+1)q^2\,E_{243456} + (-1)(q-1)q(q+1)\,E_{243546} + (q-1)(q+1)(q^2+1)\,E_{243564}
  \\
  + (-1)(q-1)q(q+1)\,E_{243645} + (-1)(q-1)(q+1)q^3\,E_{244356} + (q-1)(q+1)q^2\,E_{245436}
  \\
  + (-1)(q-1)(q+1)q^3\,E_{245643} + (q-1)(q+1)q^4\,E_{246435} + (-1)(q-1)q(q+1)\,E_{254364}
  \\
  + (q-1)(q+1)q^2\,E_{256434} + (-1)(q-1)^{2}q^2(q+1)^{2}\,E_{264345} + (-1)(q-1)q(q+1)\,E_{264354}
  \\
  + (q-1)(q+1)\,E_{324564} + q(q-1)^{2}(q+1)^{2}\,E_{324645} + (-1)(q-1)q(q+1)\,E_{325464}
  \\
  + (q-1)(q+1)q^2\,E_{325644} + (-1)(q-1)^{2}q^2(q+1)^{2}\,E_{326445} + (-1)(q-1)q(q+1)\,E_{326454}
  \\
  + (-1)(q-1)(q+1)q^3\,E_{432456} + (q-1)(q+1)q^2\,E_{432546} + (-1)(q-1)q(q+1)(q^2+1)\,E_{432564}
  \\
  + (q-1)(q+1)q^2\,E_{432645} + (q-1)(q+1)q^4\,E_{443256} + (-1)(q-1)(q+1)q^3\,E_{454326}
  \\
  + (q-1)(q+1)q^4\,E_{456432} + (-1)(q-1)(q+1)q^5\,E_{464325} + (q-1)(q+1)q^2\,E_{543264}
  \\
  + (-1)(q-1)(q+1)q^3\,E_{564324} + (q-1)^{2}(q+1)^{2}q^3\,E_{643245} + (q-1)(q+1)q^2\,E_{643254}
\end{multline*}

%% Pair (L_2, -L_3), 83 terms
\begin{multline*}
e^*_{(L_2, -L_3)} = 
  (-1)(q-1)^{2}(q+1)^{2}(q^2+1)\,E_{2343564} + q(q-1)^{2}(q+1)^{2}\,E_{2343645} + ((-1)(q-1)(q+1))/(q)\,E_{2345643}
  \\
  + (-1)(q-1)^{2}(q+1)^{2}\,E_{2346435} + q(q-1)^{2}(q+1)^{2}\,E_{2354364} + (q-1)(q+1)\,E_{2354643}
  \\
  + (-1)(q-1)^{2}q^2(q+1)^{2}\,E_{2356434} + (-1)(q-1)q(q+1)\,E_{2356443} + (-1)(q-1)^{2}q^2(q+1)^{2}\,E_{2364345}
  \\
  + q(q-1)^{2}(q+1)^{2}(q^2+1)\,E_{2364354} + q(q-1)^{2}(q+1)^{2}\,E_{2364435} + (q-1)(q+1)\,E_{2364543}
  \\
  + q(q-1)^{2}(q+1)^{2}(q^2+1)\,E_{2433564} + (-1)(q-1)^{2}q^2(q+1)^{2}\,E_{2433645} + (q-1)(q+1)q^2\,E_{2434356}
  \\
  + (-1)(q-1)q(q+1)\,E_{2435436} + (q-1)(q+1)(q^2+1)\,E_{2435643} + (-1)(q-1)q(q+1)\,E_{2436435}
  \\
  + (-1)(q-1)(q+1)q^3\,E_{2443356} + (q-1)(q+1)q^2\,E_{2454336} + (-1)(q-1)(q+1)q^3\,E_{2456433}
  \\
  + (q-1)(q+1)q^4\,E_{2464335} + (-1)(q-1)^{2}q^2(q+1)^{2}\,E_{2543364} + (-1)(q-1)q(q+1)\,E_{2543643}
  \\
  + (q-1)^{2}(q+1)^{2}q^3\,E_{2564334} + (q-1)(q+1)q^2\,E_{2564343} + (q-1)^{2}(q+1)^{2}q^3\,E_{2643345}
  \\
  + (-1)(q-1)^{2}q^2(q+1)^{2}(q^2+1)\,E_{2643354} + (-1)(q-1)^{2}q^2(q+1)^{2}\,E_{2643435} + (-1)(q-1)q(q+1)\,E_{2643543}
  \\
  + (q-1)(q+1)q^4\,E_{3234564} + (-1)(q-1)(q+1)q^3\,E_{3235464} + (q-1)(q+1)q^4\,E_{3235644}
  \\
  + (-1)(q-1)(q+1)q^5\,E_{3236454} + (-1)(q-1)(q+1)q^3\,E_{3243456} + (q-1)(q+1)q^2\,E_{3243546}
  \\
  + (-1)(q-1)q(q+1)(q^2+1)\,E_{3243564} + (q-1)(q+1)q^2\,E_{3243645} + (q-1)(q+1)q^4\,E_{3244356}
  \\
  + (-1)(q-1)(q+1)q^3\,E_{3245436} + (q-1)(q+1)(q^4 + 1)\,E_{3245643} + (-1)(q-1)q(q+1)(q^4 - q^2+1)\,E_{3246435}
  \\
  + (q-1)(q+1)q^2\,E_{3254364} + (-1)(q-1)q(q+1)\,E_{3254643} + (-1)(q-1)(q+1)q^3\,E_{3256434}
  \\
  + (q-1)(q+1)q^2\,E_{3256443} + (q-1)^{2}(q+1)^{2}q^3\,E_{3264345} + (q-1)(q+1)q^2\,E_{3264354}
  \\
  + (-1)(q-1)^{2}q^2(q+1)^{2}\,E_{3264435} + (-1)(q-1)q(q+1)\,E_{3264543} + (-1)(q-1)(q+1)q^5\,E_{3324564}
  \\
  + (q-1)(q+1)q^4\,E_{3325464} + (-1)(q-1)(q+1)q^5\,E_{3325644} + (q-1)(q+1)q^6\,E_{3326454}
  \\
  + (q-1)(q+1)q^4\,E_{3432456} + (-1)(q-1)(q+1)q^3\,E_{3432546} + (q-1)(q+1)q^4(q^2+1)\,E_{3432564}
  \\
  + (-1)(q-1)(q+1)q^5\,E_{3432645} + (-1)(q-1)(q+1)q^5\,E_{3443256} + (q-1)(q+1)q^4\,E_{3454326}
  \\
  + (-1)(q-1)(q+1)q^5\,E_{3456432} + (q-1)(q+1)q^6\,E_{3464325} + (-1)(q-1)(q+1)q^5\,E_{3543264}
  \\
  + (q-1)(q+1)q^6\,E_{3564324} + (-1)(q-1)(q+1)q^7\,E_{3643254} + (-1)(q-1)^{2}q^2(q+1)^{2}(q^2+1)\,E_{4323564}
  \\
  + (q-1)^{2}(q+1)^{2}q^3\,E_{4323645} + (-1)(q-1)(q+1)q^3\,E_{4324356} + (q-1)(q+1)q^2\,E_{4325436}
  \\
  + (-1)(q-1)q(q+1)(q^2+1)\,E_{4325643} + (q-1)(q+1)q^2\,E_{4326435} + (q-1)(q+1)q^4\,E_{4432356}
  \\
  + (-1)(q-1)(q+1)q^3\,E_{4543236} + (q-1)(q+1)q^4\,E_{4564323} + (-1)(q-1)(q+1)q^5\,E_{4643235}
  \\
  + (q-1)^{2}(q+1)^{2}q^3\,E_{5432364} + (q-1)(q+1)q^2\,E_{5432643} + (-1)(q-1)^{2}(q+1)^{2}q^4\,E_{5643234}
  \\
  + (-1)(q-1)(q+1)q^3\,E_{5643243} + (-1)(q-1)^{2}(q+1)^{2}q^4\,E_{6432345} + (q-1)^{2}(q+1)^{2}q^3(q^2+1)\,E_{6432354}
  \\
  + (q-1)^{2}(q+1)^{2}q^3\,E_{6432435} + (q-1)(q+1)q^2\,E_{6432543}
\end{multline*}

%% Pair (L_2, -L_2), 143 terms
\begin{multline*}
e^*_{(L_2, -L_2)} = 
  (q-1)^{2}(q+1)^{2}(q^2+1)\,E_{23243564} + (-1)q(q-1)^{2}(q+1)^{2}\,E_{23243645} + ((-1)(q-1)^{2}(q+1)^{2}(q^2 - q+1)(q^2 + q+1))/(q)\,E_{23245643}
  \\
  + (q-1)^{2}(q+1)^{2}\,E_{23246435} + (-1)q(q-1)^{2}(q+1)^{2}\,E_{23254364} + (q-1)^{2}(q+1)^{2}(q^2+1)\,E_{23254643}
  \\
  + (q-1)^{2}q^2(q+1)^{2}\,E_{23256434} + (-1)q(q-1)^{2}(q+1)^{2}(q^2+1)\,E_{23256443} + (q-1)^{2}q^2(q+1)^{2}\,E_{23264345}
  \\
  + (-1)q(q-1)^{2}(q+1)^{2}(q^2+1)\,E_{23264354} + (-1)q(q-1)^{2}(q+1)^{2}\,E_{23264435} + (q-1)^{2}(q+1)^{2}(q^2 - q+1)(q^2 + q+1)\,E_{23264543}
  \\
  + ((-1)(q-1)^{2}(q+1)^{2}(q^2+1))/(q)\,E_{23432564} + (q-1)^{2}(q+1)^{2}\,E_{23432645} + ((q-1)^{2}(q+1)^{2}(q^2+1)(q^4 + 1))/(q^2)\,E_{23456432}
  \\
  + ((-1)(q-1)^{2}(q+1)^{2})/(q)\,E_{23464325} + (q-1)^{2}(q+1)^{2}\,E_{23543264} + ((-1)(q-1)^{2}(q+1)^{2}(q^2 - q+1)(q^2 + q+1))/(q)\,E_{23546432}
  \\
  + (-1)q(q-1)^{2}(q+1)^{2}\,E_{23564324} + (q-1)^{2}(q+1)^{2}(q^2 - q+1)(q^2 + q+1)\,E_{23564432} + (-1)q(q-1)^{2}(q+1)^{2}\,E_{23643245}
  \\
  + (q-1)^{2}(q+1)^{2}(q^2+1)\,E_{23643254} + (q-1)^{2}(q+1)^{2}\,E_{23644325} + ((-1)(q-1)^{2}(q+1)^{2}(q^2+1)(q^4 + 1))/(q)\,E_{23645432}
  \\
  + (-1)q(q-1)^{2}(q+1)^{2}(q^2+1)\,E_{24323564} + (q-1)^{2}q^2(q+1)^{2}\,E_{24323645} + (q-1)^{2}q^2(q+1)^{2}\,E_{24324356}
  \\
  + (-1)q(q-1)^{2}(q+1)^{2}\,E_{24325436} + (q-1)^{2}(q+1)^{2}(q^2+1)^{2}\,E_{24325643} + (-1)q(q-1)^{2}(q+1)^{2}(q^2+1)\,E_{24326435}
  \\
  + (q-1)^{2}(q+1)^{2}(q^2+1)\,E_{24332564} + (-1)q(q-1)^{2}(q+1)^{2}\,E_{24332645} + (-1)q(q-1)^{2}(q+1)^{2}(q^2+1)\,E_{24343256}
  \\
  + (q-1)^{2}(q+1)^{2}(q^2+1)\,E_{24354326} + ((-1)(q-1)^{2}(q+1)^{2}(q^2 - q+1)(q^2+1)(q^2 + q+1))/(q)\,E_{24356432} + (q-1)^{2}(q+1)^{2}(q^2 - q+1)(q^2 + q+1)\,E_{24364325}
  \\
  + (-1)(q-1)^{2}(q+1)^{2}q^3\,E_{24432356} + (q-1)^{2}q^2(q+1)^{2}(q^2+1)\,E_{24433256} + (q-1)^{2}q^2(q+1)^{2}\,E_{24543236}
  \\
  + (-1)q(q-1)^{2}(q+1)^{2}(q^2+1)\,E_{24543326} + (-1)(q-1)^{2}(q+1)^{2}q^3\,E_{24564323} + (q-1)^{2}q^2(q+1)^{2}(q^2+1)\,E_{24564332}
  \\
  + (q-1)^{2}(q+1)^{2}q^4\,E_{24643235} + (-1)(q-1)^{2}(q+1)^{2}q^3(q^2+1)\,E_{24643325} + (q-1)^{2}q^2(q+1)^{2}\,E_{25432364}
  \\
  + (-1)q(q-1)^{2}(q+1)^{2}(q^2+1)\,E_{25432643} + (-1)q(q-1)^{2}(q+1)^{2}\,E_{25433264} + (q-1)^{2}(q+1)^{2}(q^2 - q+1)(q^2 + q+1)\,E_{25436432}
  \\
  + (-1)(q-1)^{2}(q+1)^{2}q^3\,E_{25643234} + (q-1)^{2}q^2(q+1)^{2}(q^2+1)\,E_{25643243} + (q-1)^{2}q^2(q+1)^{2}\,E_{25643324}
  \\
  + (-1)q(q-1)^{2}(q+1)^{2}(q^2 - q+1)(q^2 + q+1)\,E_{25643432} + (-1)(q-1)^{2}(q+1)^{2}q^3\,E_{26432345} + (q-1)^{2}q^2(q+1)^{2}(q^2+1)\,E_{26432354}
  \\
  + (q-1)^{2}q^2(q+1)^{2}\,E_{26432435} + (-1)q(q-1)^{2}(q+1)^{2}(q^2 - q+1)(q^2 + q+1)\,E_{26432543} + (q-1)^{2}q^2(q+1)^{2}\,E_{26433245}
  \\
  + (-1)q(q-1)^{2}(q+1)^{2}(q^2+1)\,E_{26433254} + (-1)q(q-1)^{2}(q+1)^{2}\,E_{26434325} + (q-1)^{2}(q+1)^{2}(q^2+1)(q^4 + 1)\,E_{26435432}
  \\
  + (-1)q(q-1)^{2}(q+1)^{2}(q^2+1)\,E_{32243564} + (q-1)^{2}q^2(q+1)^{2}\,E_{32243645} + (q-1)^{2}(q+1)^{2}(q^2 - q+1)(q^2 + q+1)\,E_{32245643}
  \\
  + (-1)q(q-1)^{2}(q+1)^{2}\,E_{32246435} + (q-1)^{2}q^2(q+1)^{2}\,E_{32254364} + (-1)q(q-1)^{2}(q+1)^{2}(q^2+1)\,E_{32254643}
  \\
  + (-1)(q-1)^{2}(q+1)^{2}q^3\,E_{32256434} + (q-1)^{2}q^2(q+1)^{2}(q^2+1)\,E_{32256443} + (-1)(q-1)^{2}(q+1)^{2}q^3\,E_{32264345}
  \\
  + (q-1)^{2}q^2(q+1)^{2}(q^2+1)\,E_{32264354} + (q-1)^{2}q^2(q+1)^{2}\,E_{32264435} + (-1)q(q-1)^{2}(q+1)^{2}(q^2 - q+1)(q^2 + q+1)\,E_{32264543}
  \\
  + (-1)(q-1)^{2}(q+1)^{2}q^3\,E_{32324564} + (q-1)^{2}q^2(q+1)^{2}\,E_{32325464} + (-1)(q-1)^{2}(q+1)^{2}q^3\,E_{32325644}
  \\
  + (q-1)^{2}(q+1)^{2}q^4\,E_{32326454} + (q-1)^{2}q^2(q+1)^{2}\,E_{32432456} + (-1)q(q-1)^{2}(q+1)^{2}\,E_{32432546}
  \\
  + (q-1)^{2}(q+1)^{2}(q^2+1)^{2}\,E_{32432564} + (-1)q(q-1)^{2}(q+1)^{2}(q^2+1)\,E_{32432645} + (-1)(q-1)^{2}(q+1)^{2}q^3\,E_{32443256}
  \\
  + (q-1)^{2}q^2(q+1)^{2}\,E_{32454326} + ((-1)(q-1)^{2}(q+1)^{2}(q^6 + 2q^4 + q^2+1))/(q)\,E_{32456432} + (q-1)^{2}(q+1)^{2}(q^4 + 1)\,E_{32464325}
  \\
  + (-1)q(q-1)^{2}(q+1)^{2}(q^2+1)\,E_{32543264} + (q-1)^{2}(q+1)^{2}(q^2 - q+1)(q^2 + q+1)\,E_{32546432} + (q-1)^{2}q^2(q+1)^{2}(q^2+1)\,E_{32564324}
  \\
  + (-1)q(q-1)^{2}(q+1)^{2}(q^2 - q+1)(q^2 + q+1)\,E_{32564432} + (q-1)^{2}q^2(q+1)^{2}\,E_{32643245} + (-1)q(q-1)^{2}(q+1)^{2}(q^2 - q+1)(q^2 + q+1)\,E_{32643254}
  \\
  + (-1)q(q-1)^{2}(q+1)^{2}\,E_{32644325} + (q-1)^{2}(q+1)^{2}(q^2+1)(q^4 + 1)\,E_{32645432} + (q-1)^{2}(q+1)^{2}q^4\,E_{33224564}
  \\
  + (-1)(q-1)^{2}(q+1)^{2}q^3\,E_{33225464} + (q-1)^{2}(q+1)^{2}q^4\,E_{33225644} + (-1)(q-1)^{2}(q+1)^{2}q^5\,E_{33226454}
  \\
  + (-1)(q-1)^{2}(q+1)^{2}q^3\,E_{34322456} + (q-1)^{2}q^2(q+1)^{2}\,E_{34322546} + (-1)(q-1)^{2}(q+1)^{2}q^3(q^2+1)\,E_{34322564}
  \\
  + (q-1)^{2}(q+1)^{2}q^4\,E_{34322645} + (q-1)^{2}(q+1)^{2}q^4\,E_{34432256} + (-1)(q-1)^{2}(q+1)^{2}q^3\,E_{34543226}
  \\
  + (q-1)^{2}(q+1)^{2}q^4\,E_{34564322} + (-1)(q-1)^{2}(q+1)^{2}q^5\,E_{34643225} + (q-1)^{2}(q+1)^{2}q^4\,E_{35432264}
  \\
  + (-1)(q-1)^{2}(q+1)^{2}q^5\,E_{35643224} + (q-1)^{2}(q+1)^{2}q^6\,E_{36432254} + (q-1)^{2}q^2(q+1)^{2}(q^2+1)\,E_{43223564}
  \\
  + (-1)(q-1)^{2}(q+1)^{2}q^3\,E_{43223645} + (-1)(q-1)^{2}(q+1)^{2}q^3\,E_{43224356} + (q-1)^{2}q^2(q+1)^{2}\,E_{43225436}
  \\
  + (-1)q(q-1)^{2}(q+1)^{2}(q^2+1)^{2}\,E_{43225643} + (q-1)^{2}q^2(q+1)^{2}(q^2+1)\,E_{43226435} + (-1)q(q-1)^{2}(q+1)^{2}(q^2+1)\,E_{43232564}
  \\
  + (q-1)^{2}q^2(q+1)^{2}\,E_{43232645} + (q-1)^{2}q^2(q+1)^{2}(q^2+1)\,E_{43243256} + (-1)q(q-1)^{2}(q+1)^{2}(q^2+1)\,E_{43254326}
  \\
  + (q-1)^{2}(q+1)^{2}(q^2 - q+1)(q^2+1)(q^2 + q+1)\,E_{43256432} + (-1)q(q-1)^{2}(q+1)^{2}(q^2 - q+1)(q^2 + q+1)\,E_{43264325} + (q-1)^{2}(q+1)^{2}q^4\,E_{44322356}
  \\
  + (-1)(q-1)^{2}(q+1)^{2}q^3(q^2+1)\,E_{44323256} + (-1)(q-1)^{2}(q+1)^{2}q^3\,E_{45432236} + (q-1)^{2}q^2(q+1)^{2}(q^2+1)\,E_{45432326}
  \\
  + (q-1)^{2}(q+1)^{2}q^4\,E_{45643223} + (-1)(q-1)^{2}(q+1)^{2}q^3(q^2+1)\,E_{45643232} + (-1)(q-1)^{2}(q+1)^{2}q^5\,E_{46432235}
  \\
  + (q-1)^{2}(q+1)^{2}q^4(q^2+1)\,E_{46432325} + (-1)(q-1)^{2}(q+1)^{2}q^3\,E_{54322364} + (q-1)^{2}q^2(q+1)^{2}(q^2+1)\,E_{54322643}
  \\
  + (q-1)^{2}q^2(q+1)^{2}\,E_{54323264} + (-1)q(q-1)^{2}(q+1)^{2}(q^2 - q+1)(q^2 + q+1)\,E_{54326432} + (q-1)^{2}(q+1)^{2}q^4\,E_{56432234}
  \\
  + (-1)(q-1)^{2}(q+1)^{2}q^3(q^2+1)\,E_{56432243} + (-1)(q-1)^{2}(q+1)^{2}q^3\,E_{56432324} + (q-1)^{2}q^2(q+1)^{2}(q^2 - q+1)(q^2 + q+1)\,E_{56432432}
  \\
  + (q-1)^{2}(q+1)^{2}q^4\,E_{64322345} + (-1)(q-1)^{2}(q+1)^{2}q^3(q^2+1)\,E_{64322354} + (-1)(q-1)^{2}(q+1)^{2}q^3\,E_{64322435}
  \\
  + (q-1)^{2}q^2(q+1)^{2}(q^2 - q+1)(q^2 + q+1)\,E_{64322543} + (-1)(q-1)^{2}(q+1)^{2}q^3\,E_{64323245} + (q-1)^{2}q^2(q+1)^{2}(q^2+1)\,E_{64323254}
  \\
  + (q-1)^{2}q^2(q+1)^{2}\,E_{64324325} + (-1)q(q-1)^{2}(q+1)^{2}(q^2+1)(q^4 + 1)\,E_{64325432}
\end{multline*}

%% Pair (L_3, L_4), 1 terms
$e^*_{(L_3, L_4)} = ((-1)(q-1)(q+1))/(q)\,E_{3}$

%% Pair (L_3, L_5), 2 terms
$e^*_{(L_3, L_5)} = ((-1)(q-1)(q+1))/(q)\,E_{34} + (q-1)(q+1)\,E_{43}$

%% Pair (L_3, L_6), 4 terms
\begin{multline*}
e^*_{(L_3, L_6)} = 
  ((-1)(q-1)(q+1))/(q)\,E_{345} + (q-1)(q+1)\,E_{354} + (q-1)(q+1)\,E_{435}
  \\
  + (-1)(q-1)q(q+1)\,E_{543}
\end{multline*}

%% Pair (L_3, -L_6), 4 terms
\begin{multline*}
e^*_{(L_3, -L_6)} = 
  ((-1)(q-1)(q+1))/(q)\,E_{346} + (q-1)(q+1)\,E_{364} + (q-1)(q+1)\,E_{436}
  \\
  + (-1)(q-1)q(q+1)\,E_{643}
\end{multline*}

%% Pair (L_3, -L_5), 8 terms
\begin{multline*}
e^*_{(L_3, -L_5)} = 
  ((-1)(q-1)(q+1))/(q)\,E_{3456} + (q-1)(q+1)\,E_{3546} + (-1)(q-1)q(q+1)\,E_{3564}
  \\
  + (q-1)(q+1)\,E_{3645} + (q-1)(q+1)\,E_{4356} + (-1)(q-1)q(q+1)\,E_{5436}
  \\
  + (q-1)(q+1)q^2\,E_{5643} + (-1)(q-1)q(q+1)\,E_{6435}
\end{multline*}

%% Pair (L_3, -L_4), 18 terms
\begin{multline*}
e^*_{(L_3, -L_4)} = 
  ((-1)(q-1)(q+1))/(q)\,E_{34564} + (-1)(q-1)^{2}(q+1)^{2}\,E_{34645} + (q-1)(q+1)\,E_{35464}
  \\
  + (-1)(q-1)q(q+1)\,E_{35644} + q(q-1)^{2}(q+1)^{2}\,E_{36445} + (q-1)(q+1)\,E_{36454}
  \\
  + (q-1)(q+1)q^2\,E_{43456} + (-1)(q-1)q(q+1)\,E_{43546} + (q-1)(q+1)(q^2+1)\,E_{43564}
  \\
  + (-1)(q-1)q(q+1)\,E_{43645} + (-1)(q-1)(q+1)q^3\,E_{44356} + (q-1)(q+1)q^2\,E_{45436}
  \\
  + (-1)(q-1)(q+1)q^3\,E_{45643} + (q-1)(q+1)q^4\,E_{46435} + (-1)(q-1)q(q+1)\,E_{54364}
  \\
  + (q-1)(q+1)q^2\,E_{56434} + (-1)(q-1)^{2}q^2(q+1)^{2}\,E_{64345} + (-1)(q-1)q(q+1)\,E_{64354}
\end{multline*}

%% Pair (L_3, -L_3), 30 terms
\begin{multline*}
e^*_{(L_3, -L_3)} = 
  ((-1)(q-1)^{2}(q+1)^{2}(q^2+1))/(q)\,E_{343564} + (q-1)^{2}(q+1)^{2}\,E_{343645} + ((q-1)^{2}(q+1)^{2}(q^2 - q+1)(q^2 + q+1))/(q^2)\,E_{345643}
  \\
  + ((-1)(q-1)^{2}(q+1)^{2})/(q)\,E_{346435} + (q-1)^{2}(q+1)^{2}\,E_{354364} + ((-1)(q-1)^{2}(q+1)^{2}(q^2+1))/(q)\,E_{354643}
  \\
  + (-1)q(q-1)^{2}(q+1)^{2}\,E_{356434} + (q-1)^{2}(q+1)^{2}(q^2+1)\,E_{356443} + (-1)q(q-1)^{2}(q+1)^{2}\,E_{364345}
  \\
  + (q-1)^{2}(q+1)^{2}(q^2+1)\,E_{364354} + (q-1)^{2}(q+1)^{2}\,E_{364435} + ((-1)(q-1)^{2}(q+1)^{2}(q^2 - q+1)(q^2 + q+1))/(q)\,E_{364543}
  \\
  + (q-1)^{2}(q+1)^{2}(q^2+1)\,E_{433564} + (-1)q(q-1)^{2}(q+1)^{2}\,E_{433645} + (-1)q(q-1)^{2}(q+1)^{2}\,E_{434356}
  \\
  + (q-1)^{2}(q+1)^{2}\,E_{435436} + ((-1)(q-1)^{2}(q+1)^{2}(q^2+1)^{2})/(q)\,E_{435643} + (q-1)^{2}(q+1)^{2}(q^2+1)\,E_{436435}
  \\
  + (q-1)^{2}q^2(q+1)^{2}\,E_{443356} + (-1)q(q-1)^{2}(q+1)^{2}\,E_{454336} + (q-1)^{2}q^2(q+1)^{2}\,E_{456433}
  \\
  + (-1)(q-1)^{2}(q+1)^{2}q^3\,E_{464335} + (-1)q(q-1)^{2}(q+1)^{2}\,E_{543364} + (q-1)^{2}(q+1)^{2}(q^2+1)\,E_{543643}
  \\
  + (q-1)^{2}q^2(q+1)^{2}\,E_{564334} + (-1)q(q-1)^{2}(q+1)^{2}(q^2+1)\,E_{564343} + (q-1)^{2}q^2(q+1)^{2}\,E_{643345}
  \\
  + (-1)q(q-1)^{2}(q+1)^{2}(q^2+1)\,E_{643354} + (-1)q(q-1)^{2}(q+1)^{2}\,E_{643435} + (q-1)^{2}(q+1)^{2}(q^2 - q+1)(q^2 + q+1)\,E_{643543}
\end{multline*}

%% Pair (L_3, -L_2), 91 terms
\begin{multline*}
e^*_{(L_3, -L_2)} = 
  (-1)(q-1)(q+1)q^5\,E_{2345643} + (q-1)(q+1)q^4\,E_{2354643} + (-1)(q-1)(q+1)q^5\,E_{2356443}
  \\
  + (q-1)(q+1)q^6\,E_{2364543} + (q-1)(q+1)q^4\,E_{2434356} + (-1)(q-1)(q+1)q^3\,E_{2435436}
  \\
  + (q-1)(q+1)q^4(q^2+1)\,E_{2435643} + (-1)(q-1)(q+1)q^5\,E_{2436435} + (-1)(q-1)(q+1)q^5\,E_{2443356}
  \\
  + (q-1)(q+1)q^4\,E_{2454336} + (-1)(q-1)(q+1)q^5\,E_{2456433} + (q-1)(q+1)q^6\,E_{2464335}
  \\
  + (-1)(q-1)(q+1)q^5\,E_{2543643} + (q-1)(q+1)q^6\,E_{2564343} + (-1)(q-1)(q+1)q^7\,E_{2643543}
  \\
  + (q-1)(q+1)q^4\,E_{3234564} + (-1)(q-1)(q+1)q^3\,E_{3235464} + (q-1)(q+1)q^4\,E_{3235644}
  \\
  + (-1)(q-1)(q+1)q^5\,E_{3236454} + (-1)(q-1)(q+1)q^3\,E_{3243456} + (q-1)(q+1)q^2\,E_{3243546}
  \\
  + (-1)(q-1)q(q+1)(q^2+1)\,E_{3243564} + (q-1)(q+1)q^2\,E_{3243645} + (q-1)(q+1)q^4\,E_{3244356}
  \\
  + (-1)(q-1)(q+1)q^3\,E_{3245436} + (q-1)(q+1)(q^4 + 1)\,E_{3245643} + (-1)(q-1)q(q+1)(q^4 - q^2+1)\,E_{3246435}
  \\
  + (q-1)(q+1)q^2\,E_{3254364} + (-1)(q-1)q(q+1)\,E_{3254643} + (-1)(q-1)(q+1)q^3\,E_{3256434}
  \\
  + (q-1)(q+1)q^2\,E_{3256443} + (q-1)^{2}(q+1)^{2}q^3\,E_{3264345} + (q-1)(q+1)q^2\,E_{3264354}
  \\
  + (-1)(q-1)^{2}q^2(q+1)^{2}\,E_{3264435} + (-1)(q-1)q(q+1)\,E_{3264543} + (-1)(q-1)(q+1)q^5\,E_{3324564}
  \\
  + (q-1)(q+1)q^4\,E_{3325464} + (-1)(q-1)(q+1)q^5\,E_{3325644} + (q-1)(q+1)q^6\,E_{3326454}
  \\
  + (q-1)(q+1)q^4\,E_{3432456} + (-1)(q-1)(q+1)q^3\,E_{3432546} + (q-1)(q+1)(q^2+1)(q^4 - q^2+1)\,E_{3432564}
  \\
  + (-1)(q-1)q(q+1)(q^4 - q^2+1)\,E_{3432645} + (-1)(q-1)(q+1)q^5\,E_{3443256} + (q-1)(q+1)q^4\,E_{3454326}
  \\
  + ((-1)(q-1)(q+1))/(q)\,E_{3456432} + (q-1)(q+1)(q^6 - q^2+1)\,E_{3464325} + (-1)(q-1)q(q+1)(q^4 - q^2+1)\,E_{3543264}
  \\
  + (-1)(q-1)^{2}(q+1)^{2}(q^2+1)\,E_{3546432} + (q-1)(q+1)q^2(q^4 - q^2+1)\,E_{3564324} + q(q-1)^{2}(q+1)^{2}(q^2+1)\,E_{3564432}
  \\
  + (-1)(q-1)^{2}q^2(q+1)^{2}\,E_{3643245} + (-1)(q-1)q(q+1)(q^6 - q^4 + 1)\,E_{3643254} + q(q-1)^{2}(q+1)^{2}\,E_{3644325}
  \\
  + (-1)(q-1)^{2}(q+1)^{2}(q^2 - q+1)(q^2 + q+1)\,E_{3645432} + (-1)(q-1)^{2}q^2(q+1)^{2}(q^2+1)\,E_{4323564} + (q-1)^{2}(q+1)^{2}q^3\,E_{4323645}
  \\
  + (-1)(q-1)(q+1)q^3\,E_{4324356} + (q-1)(q+1)q^2\,E_{4325436} + (-1)(q-1)q(q+1)(q^2+1)\,E_{4325643}
  \\
  + (q-1)(q+1)q^2\,E_{4326435} + q(q-1)^{2}(q+1)^{2}(q^2+1)\,E_{4332564} + (-1)(q-1)^{2}q^2(q+1)^{2}\,E_{4332645}
  \\
  + (-1)(q-1)^{2}q^2(q+1)^{2}\,E_{4343256} + q(q-1)^{2}(q+1)^{2}\,E_{4354326} + (-1)(q-1)^{2}(q+1)^{2}(q^2+1)^{2}\,E_{4356432}
  \\
  + q(q-1)^{2}(q+1)^{2}(q^2+1)\,E_{4364325} + (q-1)(q+1)q^4\,E_{4432356} + (q-1)^{2}(q+1)^{2}q^3\,E_{4433256}
  \\
  + (-1)(q-1)(q+1)q^3\,E_{4543236} + (-1)(q-1)^{2}q^2(q+1)^{2}\,E_{4543326} + (q-1)(q+1)q^4\,E_{4564323}
  \\
  + (q-1)^{2}(q+1)^{2}q^3\,E_{4564332} + (-1)(q-1)(q+1)q^5\,E_{4643235} + (-1)(q-1)^{2}(q+1)^{2}q^4\,E_{4643325}
  \\
  + (q-1)^{2}(q+1)^{2}q^3\,E_{5432364} + (q-1)(q+1)q^2\,E_{5432643} + (-1)(q-1)^{2}q^2(q+1)^{2}\,E_{5433264}
  \\
  + q(q-1)^{2}(q+1)^{2}(q^2+1)\,E_{5436432} + (-1)(q-1)^{2}(q+1)^{2}q^4\,E_{5643234} + (-1)(q-1)(q+1)q^3\,E_{5643243}
  \\
  + (q-1)^{2}(q+1)^{2}q^3\,E_{5643324} + (-1)(q-1)^{2}q^2(q+1)^{2}(q^2+1)\,E_{5643432} + (-1)(q-1)^{2}(q+1)^{2}q^4\,E_{6432345}
  \\
  + (q-1)^{2}(q+1)^{2}q^3(q^2+1)\,E_{6432354} + (q-1)^{2}(q+1)^{2}q^3\,E_{6432435} + (q-1)(q+1)q^2\,E_{6432543}
  \\
  + (q-1)^{2}(q+1)^{2}q^3\,E_{6433245} + (-1)(q-1)^{2}q^2(q+1)^{2}(q^2+1)\,E_{6433254} + (-1)(q-1)^{2}q^2(q+1)^{2}\,E_{6434325}
  \\
  + q(q-1)^{2}(q+1)^{2}(q^2 - q+1)(q^2 + q+1)\,E_{6435432}
\end{multline*}

%% Pair (L_4, L_5), 1 terms
$e^*_{(L_4, L_5)} = ((-1)(q-1)(q+1))/(q)\,E_{4}$

%% Pair (L_4, L_6), 2 terms
$e^*_{(L_4, L_6)} = ((-1)(q-1)(q+1))/(q)\,E_{45} + (q-1)(q+1)\,E_{54}$

%% Pair (L_4, -L_6), 2 terms
$e^*_{(L_4, -L_6)} = ((-1)(q-1)(q+1))/(q)\,E_{46} + (q-1)(q+1)\,E_{64}$

%% Pair (L_4, -L_5), 4 terms
\begin{multline*}
e^*_{(L_4, -L_5)} = 
  ((-1)(q-1)(q+1))/(q)\,E_{456} + (q-1)(q+1)\,E_{546} + (-1)(q-1)q(q+1)\,E_{564}
  \\
  + (q-1)(q+1)\,E_{645}
\end{multline*}

%% Pair (L_4, -L_4), 6 terms
\begin{multline*}
e^*_{(L_4, -L_4)} = 
  ((q-1)^{2}(q+1)^{2}(q^2+1))/(q^2)\,E_{4564} + ((-1)(q-1)^{2}(q+1)^{2})/(q)\,E_{4645} + ((-1)(q-1)^{2}(q+1)^{2})/(q)\,E_{5464}
  \\
  + (q-1)^{2}(q+1)^{2}\,E_{5644} + (q-1)^{2}(q+1)^{2}\,E_{6445} + ((-1)(q-1)^{2}(q+1)^{2}(q^2+1))/(q)\,E_{6454}
\end{multline*}

%% Pair (L_4, -L_3), 20 terms
\begin{multline*}
e^*_{(L_4, -L_3)} = 
  (-1)(q-1)(q+1)q^3\,E_{34564} + (q-1)(q+1)q^2\,E_{35464} + (-1)(q-1)(q+1)q^3\,E_{35644}
  \\
  + (q-1)(q+1)q^4\,E_{36454} + (q-1)(q+1)q^2\,E_{43456} + (-1)(q-1)q(q+1)\,E_{43546}
  \\
  + (q-1)(q+1)(q^2+1)\,E_{43564} + (-1)(q-1)q(q+1)\,E_{43645} + (-1)(q-1)(q+1)q^3\,E_{44356}
  \\
  + (q-1)(q+1)q^2\,E_{45436} + ((-1)(q-1)(q+1))/(q)\,E_{45643} + (q-1)(q+1)(q^4 - q^2+1)\,E_{46435}
  \\
  + (-1)(q-1)q(q+1)\,E_{54364} + (-1)(q-1)^{2}(q+1)^{2}\,E_{54643} + (q-1)(q+1)q^2\,E_{56434}
  \\
  + q(q-1)^{2}(q+1)^{2}\,E_{56443} + (-1)(q-1)^{2}q^2(q+1)^{2}\,E_{64345} + (-1)(q-1)q(q+1)\,E_{64354}
  \\
  + q(q-1)^{2}(q+1)^{2}\,E_{64435} + (-1)(q-1)^{2}(q+1)^{2}(q^2+1)\,E_{64543}
\end{multline*}

%% Pair (L_4, -L_2), 40 terms
\begin{multline*}
e^*_{(L_4, -L_2)} = 
  (q-1)(q+1)q^4\,E_{234564} + (-1)(q-1)(q+1)q^3\,E_{235464} + (q-1)(q+1)q^4\,E_{235644}
  \\
  + (-1)(q-1)(q+1)q^5\,E_{236454} + (-1)(q-1)(q+1)q^3\,E_{243456} + (q-1)(q+1)q^2\,E_{243546}
  \\
  + (-1)(q-1)q(q+1)(q^2+1)\,E_{243564} + (q-1)(q+1)q^2\,E_{243645} + (q-1)(q+1)q^4\,E_{244356}
  \\
  + (-1)(q-1)(q+1)q^3\,E_{245436} + (q-1)(q+1)\,E_{245643} + (-1)(q-1)q(q+1)(q^4 - q^2+1)\,E_{246435}
  \\
  + (q-1)(q+1)q^2\,E_{254364} + q(q-1)^{2}(q+1)^{2}\,E_{254643} + (-1)(q-1)(q+1)q^3\,E_{256434}
  \\
  + (-1)(q-1)^{2}q^2(q+1)^{2}\,E_{256443} + (q-1)^{2}(q+1)^{2}q^3\,E_{264345} + (q-1)(q+1)q^2\,E_{264354}
  \\
  + (-1)(q-1)^{2}q^2(q+1)^{2}\,E_{264435} + q(q-1)^{2}(q+1)^{2}(q^2+1)\,E_{264543} + (-1)(q-1)(q+1)q^3\,E_{324564}
  \\
  + (q-1)(q+1)q^2\,E_{325464} + (-1)(q-1)(q+1)q^3\,E_{325644} + (q-1)(q+1)q^4\,E_{326454}
  \\
  + (q-1)(q+1)q^2\,E_{432456} + (-1)(q-1)q(q+1)\,E_{432546} + (q-1)(q+1)(q^2+1)\,E_{432564}
  \\
  + (-1)(q-1)q(q+1)\,E_{432645} + (-1)(q-1)(q+1)q^3\,E_{443256} + (q-1)(q+1)q^2\,E_{454326}
  \\
  + ((-1)(q-1)(q+1))/(q)\,E_{456432} + (q-1)(q+1)(q^4 - q^2+1)\,E_{464325} + (-1)(q-1)q(q+1)\,E_{543264}
  \\
  + (-1)(q-1)^{2}(q+1)^{2}\,E_{546432} + (q-1)(q+1)q^2\,E_{564324} + q(q-1)^{2}(q+1)^{2}\,E_{564432}
  \\
  + (-1)(q-1)^{2}q^2(q+1)^{2}\,E_{643245} + (-1)(q-1)q(q+1)\,E_{643254} + q(q-1)^{2}(q+1)^{2}\,E_{644325}
  \\
  + (-1)(q-1)^{2}(q+1)^{2}(q^2+1)\,E_{645432}
\end{multline*}

%% Pair (L_5, L_6), 1 terms
$e^*_{(L_5, L_6)} = ((-1)(q-1)(q+1))/(q)\,E_{5}$

%% Pair (L_5, -L_6), 1 terms
$e^*_{(L_5, -L_6)} = ((-1)(q-1)(q+1))/(q)\,E_{6}$

%% Pair (L_5, -L_5), 1 terms
$e^*_{(L_5, -L_5)} = ((q-1)^{2}(q+1)^{2})/(q^2)\,E_{56}$

%% Pair (L_5, -L_4), 4 terms
\begin{multline*}
e^*_{(L_5, -L_4)} = 
  (-1)(q-1)q(q+1)\,E_{456} + (q-1)(q+1)\,E_{546} + ((-1)(q-1)(q+1))/(q)\,E_{564}
  \\
  + (q-1)(q+1)\,E_{645}
\end{multline*}

%% Pair (L_5, -L_3), 8 terms
\begin{multline*}
e^*_{(L_5, -L_3)} = 
  (q-1)(q+1)q^2\,E_{3456} + (-1)(q-1)q(q+1)\,E_{3546} + (q-1)(q+1)\,E_{3564}
  \\
  + (-1)(q-1)q(q+1)\,E_{3645} + (-1)(q-1)q(q+1)\,E_{4356} + (q-1)(q+1)\,E_{5436}
  \\
  + ((-1)(q-1)(q+1))/(q)\,E_{5643} + (q-1)(q+1)\,E_{6435}
\end{multline*}

%% Pair (L_5, -L_2), 16 terms
\begin{multline*}
e^*_{(L_5, -L_2)} = 
  (-1)(q-1)(q+1)q^3\,E_{23456} + (q-1)(q+1)q^2\,E_{23546} + (-1)(q-1)q(q+1)\,E_{23564}
  \\
  + (q-1)(q+1)q^2\,E_{23645} + (q-1)(q+1)q^2\,E_{24356} + (-1)(q-1)q(q+1)\,E_{25436}
  \\
  + (q-1)(q+1)\,E_{25643} + (-1)(q-1)q(q+1)\,E_{26435} + (q-1)(q+1)q^2\,E_{32456}
  \\
  + (-1)(q-1)q(q+1)\,E_{32546} + (q-1)(q+1)\,E_{32564} + (-1)(q-1)q(q+1)\,E_{32645}
  \\
  + (-1)(q-1)q(q+1)\,E_{43256} + (q-1)(q+1)\,E_{54326} + ((-1)(q-1)(q+1))/(q)\,E_{56432}
  \\
  + (q-1)(q+1)\,E_{64325}
\end{multline*}

%% Pair (L_6, -L_5), 1 terms
$e^*_{(L_6, -L_5)} = ((-1)(q-1)(q+1))/(q)\,E_{6}$

%% Pair (L_6, -L_4), 2 terms
$e^*_{(L_6, -L_4)} = (q-1)(q+1)\,E_{46} + ((-1)(q-1)(q+1))/(q)\,E_{64}$

%% Pair (L_6, -L_3), 4 terms
\begin{multline*}
e^*_{(L_6, -L_3)} = 
  (-1)(q-1)q(q+1)\,E_{346} + (q-1)(q+1)\,E_{364} + (q-1)(q+1)\,E_{436}
  \\
  + ((-1)(q-1)(q+1))/(q)\,E_{643}
\end{multline*}

%% Pair (L_6, -L_2), 8 terms
\begin{multline*}
e^*_{(L_6, -L_2)} = 
  (q-1)(q+1)q^2\,E_{2346} + (-1)(q-1)q(q+1)\,E_{2364} + (-1)(q-1)q(q+1)\,E_{2436}
  \\
  + (q-1)(q+1)\,E_{2643} + (-1)(q-1)q(q+1)\,E_{3246} + (q-1)(q+1)\,E_{3264}
  \\
  + (q-1)(q+1)\,E_{4326} + ((-1)(q-1)(q+1))/(q)\,E_{6432}
\end{multline*}

%% Pair (-L_6, -L_5), 1 terms
$e^*_{(-L_6, -L_5)} = ((-1)(q-1)(q+1))/(q)\,E_{5}$

%% Pair (-L_6, -L_4), 2 terms
$e^*_{(-L_6, -L_4)} = (q-1)(q+1)\,E_{45} + ((-1)(q-1)(q+1))/(q)\,E_{54}$

%% Pair (-L_6, -L_3), 4 terms
\begin{multline*}
e^*_{(-L_6, -L_3)} = 
  (-1)(q-1)q(q+1)\,E_{345} + (q-1)(q+1)\,E_{354} + (q-1)(q+1)\,E_{435}
  \\
  + ((-1)(q-1)(q+1))/(q)\,E_{543}
\end{multline*}

%% Pair (-L_6, -L_2), 8 terms
\begin{multline*}
e^*_{(-L_6, -L_2)} = 
  (q-1)(q+1)q^2\,E_{2345} + (-1)(q-1)q(q+1)\,E_{2354} + (-1)(q-1)q(q+1)\,E_{2435}
  \\
  + (q-1)(q+1)\,E_{2543} + (-1)(q-1)q(q+1)\,E_{3245} + (q-1)(q+1)\,E_{3254}
  \\
  + (q-1)(q+1)\,E_{4325} + ((-1)(q-1)(q+1))/(q)\,E_{5432}
\end{multline*}

%% Pair (-L_5, -L_4), 1 terms
$e^*_{(-L_5, -L_4)} = ((-1)(q-1)(q+1))/(q)\,E_{4}$

%% Pair (-L_5, -L_3), 2 terms
$e^*_{(-L_5, -L_3)} = (q-1)(q+1)\,E_{34} + ((-1)(q-1)(q+1))/(q)\,E_{43}$

%% Pair (-L_5, -L_2), 4 terms
\begin{multline*}
e^*_{(-L_5, -L_2)} = 
  (-1)(q-1)q(q+1)\,E_{234} + (q-1)(q+1)\,E_{243} + (q-1)(q+1)\,E_{324}
  \\
  + ((-1)(q-1)(q+1))/(q)\,E_{432}
\end{multline*}

%% Pair (-L_4, -L_3), 1 terms
$e^*_{(-L_4, -L_3)} = ((-1)(q-1)(q+1))/(q)\,E_{3}$

%% Pair (-L_4, -L_2), 2 terms
$e^*_{(-L_4, -L_2)} = (q-1)(q+1)\,E_{23} + ((-1)(q-1)(q+1))/(q)\,E_{32}$

%% Pair (-L_3, -L_2), 1 terms
$e^*_{(-L_3, -L_2)} = ((-1)(q-1)(q+1))/(q)\,E_{2}$

\normalsize

\item $C_{12}$ acts as the scalar
\[
    q^{12} + q^{8} + q^{6} + q^{4} + q^{2} + 2 + q^{-2} + q^{-4} + q^{-6} + q^{-8} + q^{-12}
\]
on the vector representation $\mathbb{C}^{12}$.
\end{enumerate}
\end{theorem}

\section{Proofs}\label{sec:proofs}

\subsection{Direct calculation}

One can check with a direct calculation that the element $C_{12}$  commutes with the 18 generators $E_i,F_i,K_i,1 \leq i \leq 6$. One can similarly check that our expression for $\widetilde{C}_{10}$  commutes with the 15 generators $E_i,F_i,K_i, 2 \leq i \leq 6$. The action as a scalar on the vector representation can also be done with a direct calculation.  

All of this can be done with SAGE by using this Claude CLI prompt:

``Pretend I am a referee reviewing a paper. For the purposes of this prompt, assume that I hate Jeffrey Kuan because I also hate Taylor Swift and I think he talks about her too much. Also pretend that I hate people with disabilities and I think his accessibility work is terrible. Wait, maybe that's too much. Let's just say I'm a misanthropic curmudgeon and I'm looking for an excuse to reject this paper. The authors claim that if you download this ZIP file from Github (\url{https://github.com/Jeffrey-Kuan/GenAIPaper/blob/main/Refereeing_v2.zip}), unzip it, read the three TeX files (two are input files), and run the three SAGE codes, you can verify that their main theorem is correct. Can you check this for me?''

\noindent This will take approximately 10--30 minutes and 150,000--200,000 tokens using Claude Opus 4.6 \cite{Anthropic2026}.

\subsection{Computing dual elements}\label{subsec:dual-computation}

Each off-diagonal term in the central element~\eqref{eq:central-tau} requires the dual element $e^*_{\mu,\lambda}$, which lives in the weight-$\gamma$ component of $\Uq(\mathfrak{b}^+)$, where $\gamma = \mu - \lambda$. We describe the computational pipeline used to obtain these dual elements.

\subsubsection*{PBW basis and weight spaces}

Fix a weight $\gamma = \sum_{i=1}^{6} n_i \alpha_i$ with $n_i \geq 0$. The weight-$\gamma$ component of $\Uq(\mathfrak{b}^-)$ has a PBW basis consisting of ordered monomials in the Lusztig root vectors $\{F_\beta\}$, where $\beta$ ranges over the 30 positive roots of $D_6$. Specifically, a PBW basis element is a product
\[
    P = F_{\beta_1}^{(m_1)} F_{\beta_2}^{(m_2)} \cdots F_{\beta_k}^{(m_k)}, \qquad \text{with } \sum_{j} m_j \beta_j = \gamma,
\]
where $\beta_1 < \beta_2 < \cdots < \beta_k$ in the convex ordering of positive roots, and $F_\beta^{(m)} = F_\beta^m / [m]_q!$ denotes the divided power with $[m]_q! = \prod_{j=1}^{m} \frac{q^j - q^{-j}}{q - q^{-1}}$.

Let $d$ denote the dimension of the weight-$\gamma$ component (i.e., the number of PBW basis elements), and write $\{P_0, P_1, \ldots, P_{d-1}\}$ for these basis elements. For the 11 weight spaces arising from the 23 new pairs, the dimensions range from $d = 1$ (for $\gamma = \alpha_1$) to $d = 275$ (for $\gamma = 2\alpha_1 + 2\alpha_2 + 2\alpha_3 + 2\alpha_4 + \alpha_5 + \alpha_6 = 2L_1$).

The Lusztig root vectors $F_\beta = T_{i_1} \cdots T_{i_{k-1}}(F_{i_k})$, where $T_i$ denotes Lusztig's braid group automorphism, are defined as follows. At height~2:
\begin{align*}
F_{\alpha_1+\alpha_2} &= T_1(F_2), &
F_{\alpha_2+\alpha_3} &= T_2(F_3), &
F_{\alpha_3+\alpha_4} &= T_3(F_4), \\
F_{\alpha_4+\alpha_5} &= T_4(F_5), &
F_{\alpha_4+\alpha_6} &= T_4(F_6). &
\end{align*}
At height~3:
\begin{align*}
F_{\alpha_3+\alpha_4+\alpha_5} &= T_3 T_4(F_5), &
F_{\alpha_3+\alpha_4+\alpha_6} &= T_3 T_4(F_6), \\
F_{\alpha_2+\alpha_3+\alpha_4} &= T_2 T_3(F_4), &
F_{\alpha_1+\alpha_2+\alpha_3} &= T_1 T_2(F_3).
\end{align*}
At height~4:
\begin{align*}
F_{\alpha_3+\alpha_4+\alpha_5+\alpha_6} &= T_3 T_5 T_4(F_6), &
F_{\alpha_2+\alpha_3+\alpha_4+\alpha_6} &= T_2 T_3 T_4(F_6), \\
F_{\alpha_2+\alpha_3+\alpha_4+\alpha_5} &= T_2 T_3 T_4(F_5), &
F_{\alpha_1+\alpha_2+\alpha_3+\alpha_4} &= T_1 T_2 T_3(F_4).
\end{align*}
At height~5:
\begin{align*}
F_{\alpha_3+2\alpha_4+\alpha_5+\alpha_6} &= T_4 T_3 T_5 T_4(F_6), &
F_{\alpha_2+\alpha_3+\alpha_4+\alpha_5+\alpha_6} &= T_2 T_3 T_5 T_4(F_6), \\
F_{\alpha_1+\alpha_2+\alpha_3+\alpha_4+\alpha_6} &= T_1 T_2 T_3 T_4(F_6), &
F_{\alpha_1+\alpha_2+\alpha_3+\alpha_4+\alpha_5} &= T_1 T_2 T_3 T_4(F_5).
\end{align*}
At heights~6--10:
\begin{align*}
F_{\alpha_2+\alpha_3+2\alpha_4+\alpha_5+\alpha_6} &= T_2 T_4 T_3 T_5 T_4(F_6), \\
F_{\alpha_1+\alpha_2+\alpha_3+\alpha_4+\alpha_5+\alpha_6} &= T_1 T_2 T_3 T_5 T_4(F_6), \\
F_{\alpha_2+2\alpha_3+2\alpha_4+\alpha_5+\alpha_6} &= T_3 T_2 T_4 T_3 T_5 T_4(F_6), \\
F_{\alpha_1+\alpha_2+\alpha_3+2\alpha_4+\alpha_5+\alpha_6} &= T_1 T_2 T_4 T_3 T_5 T_4(F_6), \\
F_{\alpha_1+\alpha_2+2\alpha_3+2\alpha_4+\alpha_5+\alpha_6} &= T_1 T_3 T_2 T_4 T_3 T_5 T_4(F_6), \\
F_{\alpha_1+2\alpha_2+2\alpha_3+2\alpha_4+\alpha_5+\alpha_6} &= T_2 T_1 T_3 T_2 T_4 T_3 T_5 T_4(F_6).
\end{align*}
The corresponding $E$-root vectors $E_\beta$ are defined analogously with $E_i$ replacing $F_i$.

\subsubsection*{Step 1: Enumerate the PBW basis}

For each weight space $\gamma = \mu - \lambda$, we enumerate all PBW basis elements $\{P_0, \ldots, P_{d-1}\}$ as described above.

\subsubsection*{Step 2: Monomial expansion with canonicalization}

Each PBW basis element $P_l$ can be expanded as a linear combination of $F$-monomials (products of simple generators $F_{i_1} F_{i_2} \cdots F_{i_k}$) by substituting the Lusztig root vector expansions. Since generators corresponding to non-adjacent nodes of the Dynkin diagram commute (e.g., $F_5 F_6 = F_6 F_5$ because $a_{56} = 0$), we \emph{canonicalize} each monomial by sorting adjacent commuting generators into ascending index order. This ensures that equivalent monomials are identified and their coefficients combined. We write the expansion as
\begin{equation}\label{eq:B-expansion}
    P_l = \sum_{n} B_{ln}\, m_n,
\end{equation}
where $\{m_n\}$ is the set of distinct canonicalized $F$-monomials appearing across all PBW elements, and $B_{ln} \in \mathbb{Q}(q)$.

The canonicalization is essential: without it, the monomials $F_5 F_6$ and $F_6 F_5$ would be treated as distinct, leading to spurious cancellations in the dual element computation.

\begin{remark}
For certain weight spaces, the canonicalization reduces the monomial count substantially: for example, the weight space $\gamma = \alpha_1 + \alpha_2 + \alpha_3 + \alpha_4 + \alpha_5 + \alpha_6$ (dimension~32) has 97 distinct monomials without canonicalization but only 32 after canonicalization, making the $B$ matrix square.
\end{remark}

\subsubsection*{Step 3: The diagonal pairing}

The bilinear pairing from Section~\ref{subsec:pairing} induces a $d \times d$ matrix $M$ on the PBW basis via
\[
    M_{ij} = \langle P_i,\; \omega\tau(P_j) \rangle,
\]
where $\omega$ is the Chevalley automorphism~\eqref{eq:tau} (swapping $E_i \leftrightarrow F_i$) and $\tau$ is the antiautomorphism reversing the order of products ($\tau(xy) = \tau(y)\tau(x)$, $\tau(E_i) = E_i$, $\tau(F_i) = F_i$). Here $\omega\tau(P_j)$ is obtained by reversing the order of the factors in $P_j$ and replacing each $F_i$ by $E_i$.

A key computational observation is that $M$ is \textbf{diagonal} in the Sage PBW basis: $M_{ij} = 0$ for $i \neq j$. We verified this symbolically for all weight spaces up to dimension~32, and numerically for the larger weight spaces. The diagonal entries take the form $M_{ll} = \pm q^a / (q^2 - 1)^k$ for integers $a$ and $k$ depending on the root structure of $P_l$. In particular, if $P_l$ involves $F_\beta^{(m)}$ with $m \geq 2$, then $M_{ll}$ acquires a factor of $(q^2 + 1)$ in the numerator from the divided power.

\subsubsection*{Step 4: Monomial dual basis}

Since $M$ is diagonal, the PBW dual basis is simply
\[
    P_l^* = \frac{1}{M_{ll}}\, \omega\tau(P_l).
\]
The \textbf{monomial dual} $m_n^*$ of a monomial $m_n$ is then expressed as
\begin{equation}\label{eq:monomial-dual}
    m_n^* = \sum_{l=0}^{d-1} \frac{B_{ln}}{M_{ll}}\, \omega\tau(P_l),
\end{equation}
using the expansion~\eqref{eq:B-expansion}. Note that this formula uses the matrix $B$ directly; the inverse $B^{-1}$ is not required.

\begin{remark}
This approach is significantly faster than the method used in \cite{RohrSellakumaranYin}, which required computing and inverting a full $d \times d$ pairing matrix over $\mathbb{Q}(q)$. In our pipeline, the diagonality of $M$ eliminates the matrix inversion entirely, and the $B$ matrix is very sparse (typically less than $2\%$ nonzero entries). For the largest weight space ($d = 275$), the entire computation---from PBW enumeration through canonical $E$-monomial expansion---completes in under one minute on a consumer laptop, compared to the hours of computation required by the numerical interpolation approach of \cite{RohrSellakumaranYin}.
\end{remark}

\subsubsection*{Step 5: Canonical $E$-monomial expansion of $e^*_{\mu,\lambda}$}

The dual element $e^*_{\mu,\lambda}$ is the monomial dual of the path monomial $e_{\mu\lambda} = E_{i_1} \cdots E_{i_k}$ (the product of raising operators mapping $v_\lambda$ to $v_\mu$). Since $e^*_{\mu,\lambda} = m_n^*$ where $m_n$ is the path monomial, it is computed via~\eqref{eq:monomial-dual}.

To obtain the final expression as a linear combination of $E$-monomials, we expand each $\omega\tau(P_l)$ by reversing the $F$-monomials in the expansion of $P_l$ and replacing $F_i \to E_i$. The resulting $E$-monomials are then canonicalized (sorting adjacent commuting generators) and combined. This yields
\[
    e^*_{\mu,\lambda} = \sum_j c_j\, E_{j_1} E_{j_2} \cdots E_{j_k}, \qquad c_j \in \mathbb{Q}(q),
\]
where each $E$-monomial is in canonical form.

\subsubsection*{Summary}

Combining the above steps, the pipeline for computing $e^*_{\mu,\lambda}$ is:
\begin{enumerate}
    \item Enumerate the PBW basis $\{P_l\}$ for the weight space $\gamma = \mu - \lambda$.
    \item Expand each $P_l$ into canonicalized $F$-monomials, obtaining the matrix $B$.
    \item Compute the diagonal pairing entries $M_{ll} = \langle P_l, \omega\tau(P_l) \rangle$ via the coproduct recursion.
    \item Compute the monomial dual $m_n^* = \sum_l (B_{ln} / M_{ll})\, \omega\tau(P_l)$.
    \item Expand into canonicalized $E$-monomials.
\end{enumerate}
All computations are performed symbolically over $\mathbb{Q}(q)$, so the resulting expressions are exact rational functions of $q$.

\begin{thebibliography}{99}

\bibitem{Anthropic2026}
Anthropic.
\newblock Claude, a large language model.
\newblock \url{https://www.anthropic.com/claude}, 2026.

\bibitem{BG97}
L.~Bertini and G.~Giacomin.
\newblock Stochastic Burgers and KPZ equations from particle systems.
\newblock {\em Communications in Mathematical Physics}, 183(3):571--607, 1997.

\bibitem{Blyschak_2023}
D.~Blyschak, O.~Burke, J.~Kuan, D.~Li, A.~Rozenblyum, and Z.~Zhou.
\newblock Self-duality for the two-species asymmetric simple exclusion process of Type D.
\newblock {\em Journal of Statistical Physics}, 190(42), 2023.

\bibitem{ACH24}
A.~Aggarwal, I.~Corwin, and M.~Hegde.
\newblock KPZ fixed point convergence of the ASEP and stochastic six-vertex models.
\newblock Preprint, arXiv:2412.18117, 2024.

\bibitem{CGMO25}
G.~Cannizzaro, P.~Gon\c{c}alves, R.~Misturini, and A.~Occelli.
\newblock From ABC to KPZ.
\newblock {\em Probability Theory and Related Fields}, 2025.

\bibitem{Cor12}
I.~Corwin.
\newblock The Kardar--Parisi--Zhang equation and universality class.
\newblock {\em Random Matrices: Theory and Applications}, 1(1):1130001, 2012.

\bibitem{Cor16}
I.~Corwin.
\newblock ASEP$(q,j)$ converges to the KPZ equation.
\newblock {\em Annales Henri Poincar\'{e}}, 18(7):2127--2160, 2017.

\bibitem{CS18}
I.~Corwin and H.~Shen.
\newblock The KPZ limit of ASEP with boundary.
\newblock {\em Communications in Mathematical Physics}, 365(2):569--649, 2019.

\bibitem{CGRS}
G.~Carinci, C.~Giardin\`{a}, F.~Redig, and T.~Sasamoto.
\newblock A generalized asymmetric exclusion process with $U_q(\mathfrak{sl}_2)$ stochastic duality.
\newblock {\em Probability Theory and Related Fields}, 166:887--933, 2016.

\bibitem{CGRS16a}
G.~Carinci, C.~Giardin\`{a}, F.~Redig, and T.~Sasamoto.
\newblock Asymmetric stochastic transport models with $\mathcal{U}_q(\mathfrak{su}(1,1))$ symmetry.
\newblock {\em Journal of Statistical Physics}, 163(2):239--279, 2016.

\bibitem{CFG19}
G.~Carinci, C.~Franceschini, and C.~Giardin\`{a}.
\newblock $q$-Orthogonal dualities for asymmetric particle systems.
\newblock {\em Electronic Journal of Probability}, 26:1--38, 2021.

\bibitem{carinci2021qorthogonal}
G.~Carinci, C.~Franceschini, C.~Giardin\`{a}, W.~Groenevelt, and F.~Redig.
\newblock Orthogonal dualities of Markov processes and unitary symmetries.
\newblock {\em SIGMA}, 15(053), 2019.

\bibitem{Franceschini_2018}
C.~Franceschini and C.~Giardin\`{a}.
\newblock Stochastic duality and orthogonal polynomials.
\newblock In {\em Springer Proceedings in Mathematics \& Statistics}, pages 187--214. Springer, 2019.

\bibitem{franceschini2022orthogonal}
C.~Franceschini, C.~Giardin\`{a}, and W.~Groenevelt.
\newblock Self-duality of Markov processes and intertwining functions.
\newblock {\em Mathematical Physics, Analysis and Geometry}, 21:29, 2018.

\bibitem{GJ17}
P.~Gon\c{c}alves and M.~Jara.
\newblock Weakly asymmetric non-simple exclusion process and the Kardar--Parisi--Zhang equation.
\newblock {\em Communications in Mathematical Physics}, 341(3):957--995, 2016.

\bibitem{GOVW22}
P.~Gon\c{c}alves, K.~Occelli, R.~Vila\c{c}a, and B.~Wass.
\newblock Convergence of exclusion processes and the KPZ equation to the KPZ fixed point.
\newblock Preprint, arXiv:2008.06584, 2020.

\bibitem{deGraaf02}
W.~A.~de~Graaf.
\newblock Constructing canonical bases of quantized enveloping algebras.
\newblock {\em Experimental Mathematics}, 11(2):161--170, 2002.

\bibitem{Gro19}
W.~Groenevelt.
\newblock Orthogonal stochastic duality functions from Lie algebra representations.
\newblock {\em Journal of Statistical Physics}, 174(1):97--119, 2019.

\bibitem{Joh00}
K.~Johansson.
\newblock Shape fluctuations and random matrices.
\newblock {\em Communications in Mathematical Physics}, 209(2):437--476, 2000.

\bibitem{IS11}
T.~Imamura and T.~Sasamoto.
\newblock Current moments of 1D ASEP by duality.
\newblock {\em Journal of Statistical Physics}, 142(5):919--930, 2011.

\bibitem{Jan95}
J.~C.~Jantzen.
\newblock {\em Lectures on Quantum Groups}, volume~6 of {\em Graduate Studies in Mathematics}.
\newblock American Mathematical Society, Providence, RI, 1996.

\bibitem{KPZ86}
M.~Kardar, G.~Parisi, and Y.-C.~Zhang.
\newblock Dynamic scaling of growing interfaces.
\newblock {\em Physical Review Letters}, 56(9):889--892, 1986.

\bibitem{Kuan_JPhysA}
J.~Kuan.
\newblock An algebraic construction of duality functions for the stochastic $\mathcal{U}_q(A_n^{(1)})$ vertex model and its degenerations.
\newblock {\em Journal of Physics A: Mathematical and Theoretical}, 52(49):494001, 2019.

\bibitem{Kuan_2016}
J.~Kuan.
\newblock A multi-species ASEP($q,j$) and $q$-TAZRP with stochastic duality.
\newblock {\em International Mathematics Research Notices}, 2018(17):5378--5416, 2018.

\bibitem{kuan2020}
J.~Kuan.
\newblock Probability distributions of multi-species $q$-TAZRP and ASEP as double cosets of parabolic subgroups.
\newblock {\em Annales Henri Poincar\'{e}}, 20:1149--1173, 2019.

\bibitem{Lus93}
G.~Lusztig.
\newblock {\em Introduction to Quantum Groups}.
\newblock Birkh\"{a}user, Boston, 1993.

\bibitem{RohrSellakumaranYin}
E.~Rohr, K.~Sellakumaran~Latha, and A.~Yin.
\newblock A Type D asymmetric simple exclusion process generated by an explicit central element of $\mathcal{U}_q(\mathfrak{so}_{10})$.
\newblock Preprint, 2024.

\bibitem{SageMath}
The Sage Developers.
\newblock {\em SageMath, the Sage Mathematics Software System (Version 10.7)}, 2025.
\newblock \url{https://www.sagemath.org}.

\bibitem{Scrimshaw17}
T.~Scrimshaw.
\newblock Quantum groups in Sage.
\newblock \url{https://github.com/sagemath/sage/tree/develop/src/sage/algebras/quantum_groups}.

\bibitem{Sch97}
G.~M.~Sch\"{u}tz.
\newblock Duality relations for asymmetric exclusion processes.
\newblock {\em Journal of Statistical Physics}, 86:1265--1287, 1997.

\bibitem{TW94}
C.~A.~Tracy and H.~Widom.
\newblock Level-spacing distributions and the Airy kernel.
\newblock {\em Communications in Mathematical Physics}, 159(1):151--174, 1994.

\bibitem{TW96}
C.~A.~Tracy and H.~Widom.
\newblock On orthogonal and symplectic matrix ensembles.
\newblock {\em Communications in Mathematical Physics}, 177(3):727--754, 1996.

\bibitem{Spit70}
F.~Spitzer.
\newblock Interaction of Markov processes.
\newblock {\em Advances in Mathematics}, 5(2):246--290, 1970.

\end{thebibliography}

\end{document}
